\subsection{A completed comparison model}
\label{uc:model}

\paragraph{The task of the comparison model.}
The quantity to be estimated is a finite sum of absolute values of a
polynomial kernel. Its individual rows have two large cotangent peaks.
The height of each peak changes with the phase, and the overlap between
the peaks changes too. A useful approximation must keep both pieces of
information: the sizes of the peaks and the cancellation between them.
This subsection constructs such an approximation, denoted $M_\infty$.
Its coefficients come from completed infinite sums and integrals; its
degree dependence enters through $L=\log n$.

The next subsection compares this approximation with the actual finite
sum. We first estimate the model on the whole closed phase triangle,
including its edges. We then estimate the discrepancy using the actual
rows of the finite board. The final corner comparison will use a loss
of size $uL^2$, a smaller possible gain of size $uL$, and finite errors.
Keeping this hierarchy in view explains the many constants below.

A phase is an angular displacement within a mesh cell. After folding
signs, its coordinates are $x,y\geq0$. They retain different meanings
because the mesh sizes are $n$ and $m=n+1$. Increasing $x$ moves the
samples along a row, whereas increasing $y$ moves the row itself. The
region $x+y\leq\pi$ is a triangle in the $(x,y)$ plane. Its origin
corresponds to the exact corner value used for comparison.

Throughout the comparison argument $m=n+1$, $L=\log n$, $H=1+L$,
$H_n=\sum_{j=1}^n j^{-1}$, and
\[
 \begin{gathered}
 x=|\alpha|,\quad y=|\beta|,\quad \rho=x+y\leq\pi,\quad
 c_e=\cos(\rho/2),\quad c_d=\cos((x-y)/2),\\
 a=c_ec_d,\quad u=1-a,\quad z=(1-c_e)(1-c_d).
 \end{gathered}
\]
Also put $p=\sin x$, $q=\sin y$, $P=p+q$, and $d_+=(q-p)_+$.
Thus $a=(\cos x+\cos y)/2\geq0$ and $2-c_e-c_d=u+z$.

The identity for $a$ follows from
$2\cos A\cos B=\cos(A+B)+\cos(A-B)$ with
$A=(x+y)/2$, $B=(x-y)/2$. Thus $0\leq a\leq1$, $u\geq0$,
and both $c_e,c_d$ lie in $[0,1]$, so $z\geq0$. The relation involving
$z$ is the exact multiplication
\[
 (1-c_e)(1-c_d)=1-c_e-c_d+c_ec_d.
\]
It retains the connection between the two cosine losses. Taylor's
formula near zero gives the useful size guide
\[
 u=\frac{x^2+y^2}{4}+O(\rho^4),\qquad
 z=\frac{(x+y)^2(x-y)^2}{64}+O(\rho^6).
\]
The proof below uses exact inequalities.
Also $P=p+q\leq x+y=\rho$ by $\sin t\leq t$. The notation
$d_+=(q-p)_+$ means $\max(q-p,0)$; it identifies the surviving
central logarithmic coefficient.
The two mesh directions retain their original roles throughout.
For real $C,D$ and $0<e<\pi$, define the complete cancellation
\begin{align}\label{uc:cancellation}
 I(C,D,e)&=\int_0^\pi
 \{|C\cot t+D\cot(t+e)|-|C\cot t|-|D\cot(t+e)|\}\,dt,\notag\\
 J(v,s,e)&=I((v-s)/2,(v+s)/2,e).
\end{align}
Write $g(C,D;t)$ for this integrand, with the separation understood.
We integrate the combined integrand, whose pole limits are finite.
It is invariant under simultaneous coefficient reversal, exchange of
channels, and reflection of the separation.

\paragraph{What is being cancelled.}
For real $A,B$ the elementary identity
\[
 |A+B|-|A|-|B|=
 \begin{cases}0,&AB\geq0,\\-2\min(|A|,|B|),&AB<0\end{cases}
\]
explains the name cancellation. When the channels have opposite signs,
the smaller magnitude is lost from both sides of the triangle
inequality. For values $5$ and $-2$, for example, the cancellation is
$3-5-2=-4$. In particular $g\leq0$ everywhere it is defined.

Each separate integral of $|\cot t|$ through a pole diverges. The
combined expression behaves differently. Near $t=0$, if $C\ne0$, the
channel $C\cot t$ dominates the bounded channel $D\cot(t+e)$. On
each side the displayed identity reduces $g$ to either zero or
$-2|D\cot(t+e)|$, with finite one-sided limits. The argument at the
other pole is the same; if a coefficient vanishes, the complete
cancellation is zero. Thus $I$ integrates the combined locally bounded
expression. Splitting at its finitely many transitions gives ordinary
improper integrals with finite limits.

The coordinates $v,s$ express $C+D=v$ and $D-C=s$. For nonnegative
$v,s$ with $s\leq v$, the coefficients in $J$ have the same sign; for $s>v$ they have
opposite signs. A root below means an equality point of their channel
magnitudes. These points delimit the intervals on which one particular
cotangent supplies the smaller magnitude.

\begin{lemma}[Cancellation and root formulas]\label{uc:cancellation-lemma}
For $v,s\geq0$, $0<e\leq\pi/2$, write
\[
 J(v,s,e)=2(s-v)_+\log\sin e+R(v,s,e).
\]
Then $R\leq0$,
\begin{align}\label{uc:R-bounds}
 |R|&\leq\max(2v,s+v)\log2,
 &|R_s|&\leq\log\frac{v+s}{|v-s|},\\
 R_e&=\frac{2\{v\cos e-\sqrt{(v^2-s^2\sin^2e)_+}\}}{\sin e}
 &&(s>v),\qquad R_e=0\quad(s\leq v).\notag
\end{align}
The derivative in $s$ is used as a locally integrable bound at $s=v$.
If
\begin{equation}\label{uc:root-Z}
 E(v,s,e)=J(v,s,e)-2s\log\sin e,\qquad
 Z(v,s,e)=2\int_e^{\pi/2}
       \frac{\sqrt{(v^2-s^2\sin^2t)_+}}{\sin t}\,dt,
\end{equation}
then $E\leq Z$, with equality for $s\geq v$.
All complete formulas extend continuously through zero coefficients
and moving equality roots.
\end{lemma}
\begin{proof}
On a sign interval the integrand is $-2$ times the smaller absolute
channel, and its primitive is a linear combination of logarithms of
sines. For $s\leq v$, with $c=(v-s)/2$, $b=(v+s)/2$, this gives
$J=2\{c\log(c/v)+b\log(b/v)\}$. Logarithms appear because $(\log|\sin t|)'=\cot t$. On an interval
where the two channel signs and the identity of the smaller magnitude
are fixed, the cancellation is a fixed multiple of one cotangent. Its
primitive is therefore explicit. Adding these interval integrals gives
the displayed formulas.

For $v>0$ and $0\leq s\leq v$, set $r=c/v$. The expression just obtained becomes
\[
 J=2v\{r\log r+(1-r)\log(1-r)\},\qquad 0\leq r\leq\tfrac12.
\]
The bracket ranges from $-\log2$ to zero, as follows by differentiating
it twice or by elementary convexity. This proves the sign and the
$2v\log2$ bound in this regime. The limit $r\log r\to0$ at zero
supplies the coefficient-zero convention. At $v=s=0$ the complete
cancellation is zero directly, in agreement with continuity.

For $s>v$, put
$c=(s-v)/2$, $b=(s+v)/2$. If $s\sin e\geq v$, integration gives
$J=2s\log\sin e$. Otherwise the equality roots satisfy
$\sin(2t+e)=s\sin e/v$. With
$t_1=(\arcsin(s\sin e/v)-e)/2$, the result is
\[
 J=2s\log\sin e+2\left\{
 b\log\frac{\cos t_1}{\sin(t_1+e)}
 -c\log\frac{\cos(t_1+e)}{\sin t_1}\right\}.
\]
For opposite signs an equality on the relevant sign interval reads
$c\cot t=b\cot(t+e)$. Multiplication by the sine factors and the
addition formulas give
\[
 (b-c)\sin(2t+e)=(b+c)\sin e.
\]
Here $b-c=v$ and $b+c=s$, yielding the stated root equation. Its
right side can be reached only when $s\sin e\leq v$. At strict
inequality there are two equality roots; at equality they merge; beyond
that threshold the extra root interval disappears. The square root in
$R_e$ arises from the cosine of the angle in this equation.

The explicit term $2(s-v)_+\log\sin e$ carries the singular behavior
as the separation shrinks. The remainder $R$ stays bounded in terms of
the coefficient sizes. Individual logarithms in the root formula may
look singular as a root or coefficient tends to zero. Their multiplying
coefficients make the complete expressions converge.

\paragraph{The actual integration intervals.}
Here is a way to reproduce the formulas without guessing a primitive
for the absolute value. Assume first $0<s<v$, put
$\theta=\arcsin((s/v)\sin e)$, and set
$y_0=(e-\theta)/2$, $y_1=(e+\theta)/2$. The cotangents have
opposite signs only on $(\pi/2-e,\pi/2)$ and $(\pi-e,\pi)$.
Use $t=\pi/2-u$ in the first and $t=\pi-r$ in the second.
The positive quantity whose integral is multiplied by $-2$ is
\[
\begin{array}{c|c|c}
 \text{variable}&\text{range}&\text{smaller channel magnitude}\\ \hline
 u&0<u<y_1&c\tan u\\
 u&y_1<u<e&b\tan(e-u)\\
 r&0<r<y_0&b\cot(e-r)\\
 r&y_0<r<e&c\cot r.
\end{array}
\]
The first crossing solves $c\tan u=b\tan(e-u)$ and the second
$c\cot r=b\cot(e-r)$. The addition formulas give precisely $y_1$
and $y_0$. Since $\int\tan u\,du=-\log\cos u$, adding the four
integrals gives
\[
 J=-2v\log\sin e+
 2c\log(\sin y_0\cos y_1)+2b\log(\sin y_1\cos y_0).
\]
Now
$\sin y_0\cos y_1=(\sin e-\sin\theta)/2=c\sin e/v$ and
$\sin y_1\cos y_0=b\sin e/v$. Substitution cancels the logarithms
of $\sin e$, leaving the entropy formula. The cases $s=0$ and $s=v$
follow by taking the stated continuous limits.

For $s>v$, the coefficients are $(-c,b)$. Cancellation occurs on
$(0,\pi/2-e)$ and $(\pi/2,\pi-e)$, where the cotangents have the
same sign. On the second interval the smaller magnitude is always
$c|\cot t|$, and its contribution to $J$ is $2c\log\sin e$.
On the first interval the smaller magnitude is $b\cot(t+e)$ unless
there are equality roots. If there are none, its contribution is
$2b\log\sin e$, proving $J=2s\log\sin e$. If there are roots,
write $t_2=(\pi-\theta-e)/2$ with
$\theta=\arcsin(s\sin e/v)$; then
\[
 -\tfrac12 J_{\rm first}
 =b\int_0^{t_1}\cot(t+e)dt
  +c\int_{t_1}^{t_2}\cot t\,dt
  +b\int_{t_2}^{\pi/2-e}\cot(t+e)dt.
\]
Evaluating each integral at its two endpoints and using
$t_2=\pi/2-t_1-e$ gives the root formula above. This also identifies
the intervals that disappear when the roots merge.

\paragraph{Differentiating the remainder explicitly.}
Put $D_*=(v^2-s^2\sin^2e)_+^{1/2}$. For $s<v$ the entropy
formula gives
\[
 R_s=\log\frac{v+s}{v-s},\qquad R_e=0.
\]
For $s>v$, differentiation of the three definite integrals just written
cancels their moving-boundary terms. In the root regime the remaining
terms simplify to
\[
 R_s=\log\frac{s\cos e-D_*}{s\cos e+D_*},\qquad
 R_e=\frac{2(v\cos e-D_*)}{\sin e}.
\]
In the regime $s\sin e\geq v$, the explicit formula instead gives
$R=2v\log\sin e$, so $R_s=0$ and $R_e=2v\cot e$. For $s>v$,
$0\leq D_*\leq v\cos e$; consequently the root-regime quotient
lies between $(s-v)/(s+v)$ and one. This proves the stated bound for
$|R_s|$ and also $0\leq R_e\leq2v\cot e$.

For the size bound, the endpoint value for $s>v$ has $b=c+v$ and
can be written
\[
 R(v,s,0)=2b\left\{\frac cb\log\frac cb+
                         \frac vb\log\frac vb\right\}.
\]
The bracket is between $-\log2$ and zero. Monotonicity in $e$ and
the value zero at $\pi/2$ thus give $-(s+v)\log2\leq R\leq0$.
For $s\leq v$ the earlier entropy calculation gives
$-2v\log2\leq R\leq0$. These are exactly the two parts of the
single bound in the lemma.

These expressions have matching limits at every threshold, using
$0\log0=0$. Differentiation gives \eqref{uc:R-bounds}; the moving
endpoint terms cancel because the channel magnitudes agree there.
The endpoint cancellation under differentiation can be checked with
the Leibniz rule. If two integrands $h_1,h_2$ meet at a moving root
$r(\lambda)$, their endpoint contributions have sum
\[
 \{h_1(r(\lambda),\lambda)-h_2(r(\lambda),\lambda)\}r'(\lambda)=0.
\]
The equality defining the root is precisely what makes the bracket zero.
The logarithmic bound for $R_s$ is locally integrable at $s=v$, since
$\int_0^\eta|\log t|dt=\eta(1-\log\eta)$ for $\eta<1$.
Later parameter estimates integrate this locally integrable bound over
a range of $s$.

For $s>v$, $R$ increases from
$2(c\log c+v\log v-b\log b)$ at $e=0$ to zero at $e=\pi/2$.
These endpoints, and the entropy expression for $s\leq v$, give its
sign and size bounds. Subtracting $2s\log\sin e$ gives derivative
$-2\sqrt{(v^2-s^2\sin^2e)_+}/\sin e$ in the root regime and zero
beyond its threshold. Integration from that threshold proves $E=Z$
for $s\geq v$. For $s<v$, $J\leq0$ and
$\sqrt{v^2-s^2\sin^2t}\geq s\cos t$, whence
$Z\geq-2s\log\sin e\geq E$. The case $v=0$ is direct.

The advantage of $Z$ is its nonnegative integrand and explicit
support:
\[
 \sqrt{(v^2-s^2\sin^2t)_+}>0\quad\Longleftrightarrow\quad s\sin t<v.
\]
For $s\geq v$, integrating the derivative from the threshold where
both sides vanish proves $E=Z$. For $s<v$, the comparison in the
proof uses $v^2-s^2\sin^2t\geq s^2\cos^2t$, and hence
\[
 Z\geq2s\int_e^{\pi/2}\cot t\,dt=-2s\log\sin e\geq E.
\]
Nonnegative root integrals can be reordered by Tonelli's theorem: a
nonnegative double integral has the same value in either order, even
before its finiteness is established. That fact will let us estimate two
central root contributions together.
\end{proof}

\paragraph{From reciprocal sums to a scalar profile.}
Near a pole an absolute cotangent behaves like the reciprocal of its
distance. A mesh samples both sides of the pole. The parameter
$t\in(0,1)$ gives its fractional position between consecutive samples,
with distances $t,1-t,1+t,2-t,\ldots$. Pairing these distances and
subtracting their common harmonic growth leaves $\mathcal C$. For
fixed $t$,
\[
 \frac1{k+t}+\frac1{k+1-t}-\frac2k=O(k^{-2}),
\]
so the series converges. At $t=0$ or $1$ a reciprocal term diverges,
but the sine factor in $\Phi$ removes this pole:
$\sin(\pi t)/(\pi t)\to1$. The weighted profile fits the polynomial
rows because the corresponding cotangent coefficient contains the same
sine factor. Periodic extension joins its finite endpoint values.
Centering at $t=1/2$ gives the even function $f$; its joined endpoint
$w=\pi$ modulo $2\pi$ is the cusp, where its left and right slopes
differ.

The scalar reciprocal pair and its continuous periodic profile are
\begin{align}\label{uc:scalar-definitions}
 \mathcal C(t)&=2\gamma+\frac1t+\frac1{1-t}
 +\sum_{k\geq1}\left\{\frac1{k+t}+\frac1{k+1-t}-\frac2k\right\},
 \notag\\
 C_0&=\mathcal C(1/2)=2\gamma+4\log2,\notag\\
 \Phi(t)&=\frac{\sin\pi t}{\pi}\{\mathcal C(t)-2\log\pi\},
 \quad\Phi(0)=\Phi(1)=1,\notag\\
 f(w)=\Phi(1/2+w/(2\pi)).
\end{align}
Extend $\Phi$ with period one. Then $f$ is even with period $2\pi$.
We use $\psi(t)=-\gamma+\sum_{k\geq0}\{(k+1)^{-1}-(k+t)^{-1}\}$,
so $\mathcal C(t)=-\psi(t)-\psi(1-t)$.

\paragraph{A derivative measure and a finite part.}
A Lipschitz function satisfies $|f(w)-f(v)|\leq K|w-v|$. Its first
derivative can still jump. Here its ordinary second derivative is $g$
away from the cusp, and $f'$ has a downward jump of size $J_0$ at the
joined endpoint. The notation $df'=g(w)\,dw-J_0\delta_\pi$ includes
both changes. Integrating this measure over an interval gives the
change of $f'$ across that interval; the point mass $\delta_\pi$
contributes one if the cusp lies inside. Total variation adds the
absolute sizes of all changes. On a circle the joined endpoint must
be counted once.

The kernel $1/(2\sin(t/2))$ has reciprocal singularities at both
ends of $[0,2\pi]$. Its integral against $f(w+t)$ diverges like
$-2f(w)\log\epsilon$. The counterterm in $H$ subtracts that known
divergence, keeping a finite constant. This prescription is called a
finite part. A Cauchy principal value, by comparison, takes matched
cutoffs on opposite sides to cancel signed singularities. Here the
explicit logarithmic counterterm and its normalization by $\pi$
define the value completely. They must be carried together when the
constant is compared at different phases.

In the series estimates, $\zeta(s)=\sum_{k\geq1}k^{-s}$ for $s>1$
denotes the convergent reciprocal-power sum.

\begin{lemma}[The scalar maximum, including its cusp]\label{uc:scalar-maximum}
The function $f$ is Lipschitz, $f'$ has bounded variation on the circle,
and, for its ordinary second derivative $g$ on $(-\pi,\pi)$,
\begin{equation}\label{uc:scalar-bounds}
 \begin{gathered}
 \frac12<f\leq1,\quad \operatorname{Lip}(f)<2,\quad
 |g|<1,\quad |f'''|<13/20,\quad
 \operatorname{Var}_{\rm per}f'<4,\\
 df'=g(w)\,dw-J_0\delta_\pi\leq\frac3{\pi^2}\,dw,
 \quad J_0=\frac{2(\log\pi-\gamma)}\pi>\frac1\pi.
 \end{gathered}
\end{equation}
The endpoint atom is counted once. The finite part
\begin{equation}\label{uc:H}
 H(w)=\frac1\pi\lim_{\epsilon\downarrow0}
 \left\{\int_\epsilon^{2\pi-\epsilon}
 \frac{f(w+t)}{2\sin(t/2)}\,dt+2f(w)\log\frac\epsilon\pi\right\}
\end{equation}
satisfies $H(w)\leq H_0:=H(0)$ for all real $w$.
\end{lemma}
\begin{proof}
On $0<t\leq1/2$ the defining series becomes
\[
 \mathcal C(t)=t^{-1}+2\gamma+Q(t),\qquad
 Q(t)=2\sum_{j\geq1}\zeta(2j+1)t^{2j}.
\]
The expansion of $Q$ follows by geometric series. For example,
\[
 \frac1{k+t}+\frac1{k-t}-\frac2k
 =\frac{2t^2}{k(k^2-t^2)}
 =2\sum_{j\geq1}\frac{t^{2j}}{k^{2j+1}}.
\]
For $t\leq1/2$, uniform geometric majorants justify summing and
differentiating the needed series. Summing over $k$ gives the zeta
values. Integral comparison bounds a decreasing positive sum by
integrals over adjacent unit intervals. This supplies simple rational
bounds for the zeta values and their tails. The constants used here
are deliberately broad enough to make subsequent error estimates easy
to check.

Integral comparison gives $\zeta(3)<5/4$, hence
$Q\leq4t^2$, $Q'\leq5t/(1-t^2)^2\leq(32/3)t$, and
$Q''\leq Q''(1/2)=28\zeta(3)-16<19$
(the weaker absolute bound $608/27$ is also convenient).
With $k=2\gamma-2\log\pi$,
\[
 \Phi(t)=\frac{\sin\pi t}{\pi t}
             +k\frac{\sin\pi t}\pi+Q(t)\frac{\sin\pi t}\pi.
\]
Since $k+Q<0$, this is at most one. For the lower bound, put
$k_0=C_0-2\log\pi$ and $c_s=\cos(\pi s/2)$. Reciprocal pairing gives
\begin{equation}\label{uc:f-series}
 f(w)=\frac{\cos(w/2)}\pi
 \left\{k_0+4w^2\sum_{\ell\geq1,\ \ell\ {\rm odd}}
       \frac1{\ell(\ell^2\pi^2-w^2)}\right\},\quad |w|<\pi.
\end{equation}
The inequality $c_s\geq(3/4)(1-s^2)$ follows by differentiating
$\sin(\pi r/2)-(3/4)r(2-r)$, $r=1-s$: its second derivative
changes sign once, and its first derivative ends at zero. Keeping
$\ell=1$ in \eqref{uc:f-series}, and using
$1-c_s\leq\pi^2s^2/8$, $k_0<2$, $\pi^2<10$, gives
$\pi f(\pi s)\geq c_sk_0+4c_ss^2/(1-s^2)\geq k_0$.
The elementary bounds $\gamma>23/40$, $\log2>69/100$,
$\log\pi<23/20$ give $k_0>161/100>\pi/2$.

The centered reciprocal series also explains the lower bound. Its
$\ell=1$ term is positive and becomes large near the endpoint, while
the prefactor $c_s=\cos(\pi s/2)$ removes its pole. In detail,
\[
 c_sk_0+\frac{4c_ss^2}{1-s^2}-k_0
 \geq-\frac{k_0\pi^2s^2}{8}+3s^2\geq0,
\]
using $c_s\geq(3/4)(1-s^2)$, $k_0<2$, and $\pi^2<10$.
Thus $f(w)\geq k_0/\pi>1/2$, with endpoint values supplied by
continuity. The upper bound comes from the endpoint representation of
$\Phi$. Positivity will matter when weighted reciprocal estimates
are multiplied later.

\paragraph{A table for the scalar derivative bounds.}
The reciprocal formula for $Q$ and its first two derivatives is
\[
 \begin{aligned}
 Q(t)&=2t^2\sum_{k\geq1}\frac1{k(k^2-t^2)},\\
 Q'(t)&=4t\sum_{k\geq1}\frac{k}{(k^2-t^2)^2},\\
 Q''(t)&=4\sum_{k\geq1}\frac{k}{(k^2-t^2)^2}
           +16t^2\sum_{k\geq1}\frac{k}{(k^2-t^2)^3}.
 \end{aligned}
\]
On $t\leq1/2$, $k^2-t^2\geq3k^2/4$. Using
$\sum k^{-3}<3/2$ and $\sum k^{-5}<5/4$ in the last line gives
$Q''\leq32/3+320/27=608/27$. The sharper endpoint value follows
by differentiating the reciprocal pair twice:
$\mathcal C''(1/2)=4\sum_{k\geq0}(k+1/2)^{-3}=28\zeta(3)$.
The second derivative of $1/t$ there is $16$, so subtracting it gives
$Q''(1/2)=28\zeta(3)-16$. Positivity of the power-series coefficients
makes $Q''$ increasing. For a fully rational estimate,
\[
 \zeta(3)\leq1+\frac18+\frac1{27}+\int_3^\infty r^{-3}dr<\frac54.
\]
Summing the geometric series then gives
$Q\leq(5/2)t^2/(1-t^2)\leq4t^2$ and
$Q'\leq5t/(1-t^2)^2\leq16/3$. Also $k<-1$, from
$\gamma<3/5$, $\log\pi>11/10$, while $Q\leq1$; hence $k+Q<0$
as used for the upper bound on $\Phi$.

Let $T(t)=\sin(\pi t)/(\pi t)=\int_0^1\cos(\pi tr)dr$.
Differentiating this integral bounds $|T'|$ by $\pi/2<2$ and
$|T''|$ by $\pi^2/3<10/3$. The other differentiated terms of
$\Phi=T+k\sin(\pi t)/\pi+Q\sin(\pi t)/\pi$ give
\[
\begin{array}{c|c}
 \text{term in }\Phi''&\text{absolute upper bound}\\ \hline
 T''&10/3\\
 -k\pi\sin(\pi t)&10\\
 Q''\sin(\pi t)/\pi&608/81\\
 2Q'\cos(\pi t)&32/3\\
 -\pi Q\sin(\pi t)&10/3.
\end{array}
\]
Their sum is $2822/81<35$. Similarly the first derivatives of the
three summands are bounded respectively by $2$, $5$, and $4$:
for the last use $Q'/\pi+Q<16/9+1<4$. The factor $1/(2\pi)$
in a $w$ derivative then gives the claimed Lipschitz constant.

For the upper curvature bound retain signs. The sinc integral and
$\cos r\geq1-r^2/2$ give
\[
 T''\leq-\pi^2\left(\frac13-\frac{\pi^2}{40}\right).
\]
Moreover $\cos\pi t\leq\pi(1/2-t)$, and
\[
 3+34t^2+3t^4-20t
 \geq34(t-5/17)^2+1/17>0.
\]
Thus $2Q'\cos\pi t\leq3\pi/4$. Dropping the nonpositive last
term in the table and using $-k<5/4$ yields
$\Phi''\leq-\pi^2/12+2\pi+19/\pi<12$.
The earlier sine comparison can likewise be checked explicitly:
for $h(r)=\sin(\pi r/2)-(3/4)r(2-r)$, $h(0)=0$, $h'(1)=0$,
and $h''$ decreases through zero once. Therefore $h'$ first increases
and then decreases to zero, staying positive, and $h\geq0$.

Differentiating the three terms of $\Phi$ bounds their first
 derivatives by $2,5,4$, respectively. Their five second-derivative
contributions have absolute bounds
$10/3,10,608/81,32/3,10/3$, whose sum is less than $35$.
Thus $\operatorname{Lip}(f)<6/\pi<2$ and $|g|<35/(4\pi^2)<1$.
For the one-sided bound use
\[
 \left(\frac{\sin\pi t}{\pi t}\right)''
 \leq-\pi^2(1/3-\pi^2/40)\leq-\pi^2/12,
 \quad \frac{t(1/2-t)}{(1-t^2)^2}\leq\frac3{40}.
\]
The latter follows from $3+34t^2+3t^4-20t>0$.
Since $-k<5/4$, differentiation gives
$\Phi''\leq-\pi^2/12+2\pi+19/\pi<12$.
The absolute second-derivative bound and the upper curvature bound
serve different purposes. The absolute bound controls ordinary changes;
the upper bound retains favorable negative terms. The scaling
$t=1/2+w/(2\pi)$ divides second derivatives by $4\pi^2$, so
$\Phi''<12$ gives $g<3/\pi^2$. A downward jump only strengthens
this upper bound. Concretely, for a nonnegative continuous periodic
function $\chi$ the measure inequality says
\[
 \int\chi\,df'\leq\frac3{\pi^2}\int\chi(w)\,dw.
\]
The total change of a periodic derivative is zero. Its positive and
negative masses therefore have equal magnitudes, so bounding its
positive mass bounds its total variation by twice that amount.

At an integer $\Phi'$ has jump $2k<0$; reflection matches its
ordinary derivative at the half integer. Scaling proves the measure
bound in \eqref{uc:scalar-bounds}, including the stated atom.
The positive mass of $df'$ is at most $6/\pi$, its total mass is
zero, and therefore its variation is at most $12/\pi<4$.

For the remaining derivative facts an endpoint-cancelled formula is
useful:
\begin{equation}\label{uc:f-integral}
 f(w)=f(0)\cos(w/2)+2\int_0^{1/2}
       \frac{\cos(wt)-\cos(w/2)}{\cos\pi t}\,dt,\qquad |w|\leq\pi.
\end{equation}
At $t=1/2$ the numerator in this integral vanishes along with
$\cos\pi t$. The mean value theorem bounds
$|\cos(wt)-\cos(w/2)|$ by $|w|(1/2-t)$, while
$\cos\pi t\geq2(1/2-t)$. The quotient is consequently bounded.
Each of its first three differentiated numerators in $w$ still
vanishes at $t=1/2$. The same argument, on the compact range
$|w|\leq\pi$, supplies integrable majorants for those derivatives.

The factors $r^{2j+1}$ below first make the trigonometric series
absolutely convergent for $r<1$. The sum may then be integrated term
by term. Taking $r\uparrow1$ is justified by the displayed common
majorant and dominated convergence: pointwise convergence under one
integrable bound permits passage to the limit in the integral. This
proves the formula at the joined endpoints as well as in the interior.

To derive it, integrate the numerator against $\cos((2j+1)\pi t)$;
the result is
$(-1)^j\cos(w/2)w^2/[(2j+1)\pi((2j+1)^2\pi^2-w^2)]$.
The Abel sum of the kernels is
\[
 \sum_{j\geq0}(-1)^jr^{2j+1}\cos((2j+1)\pi t)
 =\frac{r(1+r^2)\cos\pi t}{(1-r^2)^2+4r^2\cos^2\pi t}
 \longrightarrow\frac1{2\cos\pi t}.
\]
For $r\geq1/2$ it is bounded by $5/(8\cos\pi t)$, while the
numerator is $O(1/2-t)$ and $\cos\pi t\geq2(1/2-t)$.
Dominated convergence proves \eqref{uc:f-integral}, including its
endpoints, and permits three derivatives because each differentiated
numerator also vanishes at $t=1/2$. In particular
\[
 f'''(w)=\frac{f(0)}8\sin(w/2)
 +2\int_0^{1/2}\frac{t^3\sin(tw)-\sin(w/2)/8}{\cos\pi t}\,dt.
\]
The numerator's $t$ derivative has magnitude at most
$3/4+\pi/8<23/20$, giving $|f'''|<23/40+3/40=13/20$.
For $0<w\leq\pi$, monotonicity of $r\cot r$ gives
$\sin(tw)/\sin(w/2)\leq\sin\pi t$. Since $f(0)<3/5$, split
$[0,1/2]$ at $1/4,3/8$, use respectively $\sin\pi t\leq(22/7)t$
and $t^3\sin\pi t\leq27/512$, and discard the negative last part.
This yields
\[
 \frac{f'''(w)}{\sin(w/2)}
 <\frac3{40}+\frac{44}{35\cdot1024}-\frac1{16}-\frac{37}{2048}
 =-\frac{311}{71680}<0.
\]
Consequently $g$ is even and decreases radially.

The sine-ratio step follows from ordinary differentiation. For fixed
$0<t<1/2$ set $R_t(w)=\sin(tw)/\sin(w/2)$. On $0<w<\pi$,
\[
 (\log R_t)'=t\cot(tw)-\tfrac12\cot(w/2)>0,
\]
because $(r\cot r)'=(\sin r\cos r-r)/\sin^2r<0$ for
$0<r\leq\pi/2$. Hence $R_t(w)\leq R_t(\pi)=\sin\pi t$.
After this substitution the integral bounding
$f'''(w)/\sin(w/2)$ has negative numerator
$t^3\sin\pi t-1/8$. Dividing by $\cos\pi t\leq1$ makes that
negative number smaller. Thus the first two interval contributions are
\[
 \begin{aligned}
 2\int_0^{1/4}\frac{t^3\sin\pi t-1/8}{\cos\pi t}dt
 &\leq2\int_0^{1/4}\left(\frac{22}{7}t^4-\frac18\right)dt
 =\frac{44}{35\cdot1024}-\frac1{16},\\
 2\int_{1/4}^{3/8}\frac{t^3\sin\pi t-1/8}{\cos\pi t}dt
 &\leq2\left(\frac38-\frac14\right)
                    \left(\frac{27}{512}-\frac18\right)
 =-\frac{37}{2048}.
 \end{aligned}
\]
The last interval contributes a negative amount, which may be dropped
for an upper bound. Together with $f(0)/8<3/40$, these evaluated
integrals give exactly the rational negative coefficient above.

The sign of the third derivative says that $g=f''$ decreases as
$|w|$ increases from zero to $\pi$. This is the meaning of radially
decreasing on the circle. It is the order property needed for the
finite-part maximum. Pairing the endpoints removes the logarithmic
divergence directly: the Lipschitz bound gives
\[
 |f(w+t)+f(w-t)-2f(w)|\leq2\operatorname{Lip}(f)t,
\]
and $2\sin(t/2)$ is comparable to $t$ on $[0,\pi]$. Thus the next
quotient is integrable, including at a cusp. At an ordinary smooth
point its numerator is even $O(t^2)$, but the weaker uniform estimate
is sufficient.

Pairing endpoints in \eqref{uc:H} gives the absolutely convergent
formula
\[
 H(w)=\lambda_0f(w)+\frac1\pi\int_0^\pi
 \frac{f(w+t)+f(w-t)-2f(w)}{2\sin(t/2)}\,dt,
 \qquad \lambda_0=\frac2\pi\log\frac4\pi<\frac1{2\pi}.
\]
The bound follows from $\pi>25/8$ and $e^{1/4}>32/25$.
Set $W(t)=\int_{|t|}^\pi\log\cot(s/4)\,ds$ on the circle.
It is bounded, nonnegative and radially decreasing. Integrating the
triangular-kernel representation of the second difference against
$df'$ gives
\[
 H(w)=\lambda_0f(w)+\frac1\pi\int_{-\pi}^{\pi}W(t)g(w+t)\,dt
                  -\frac{J_0}\pi W(\pi-|w|),\quad |w|\leq\pi.
\]
For a smooth function the second-difference identity is
\[
 f(w+t)+f(w-t)-2f(w)
 =\int_{-t}^{t}(t-|s|)f''(w+s)\,ds.
\]
It follows by integrating the derivative twice. For a derivative of
bounded variation the same identity holds with its finite measure
$df'$ in place of $f''\,ds$; the triangular weight is continuous.
Integrating this identity against the reciprocal-sine kernel gives
$W$. Its boundedness and the finite total variation justify the
interchange of the two integrals. The ordinary curvature yields the
convolution term; the cusp yields the separate negative term.

The layer-cake argument can be expressed using indicator functions:
$h(t)=\int_0^\infty\mathbf1_{h(t)>r}dr$ for $h\geq0$.
For an even radially decreasing function each superlevel set is a
centered arc. A product integral becomes an integral of lengths of
intersections of two arcs. Each intersection is at most the length of
the smaller arc, and equal centers attain that value. Translation can
therefore only decrease the centered convolution. Subtracting the
minimum of $g$ makes it nonnegative without changing which translate
maximizes its convolution with $W$.

Fubini is valid against the finite variation measure; $W(\pi)=0$
counts the atom only once. Subtract the minimum of $g$ to make it
nonnegative. Layer cake reduces its convolution with $W$ to lengths
of intersections of two centered circular arcs; a translate cannot
increase such a length. The convolution is therefore largest at zero.
The measure bound implies $f(w)-f(0)\leq3w^2/(2\pi^2)$, and
$W(\pi-|w|)\geq w^2/4$, since the derivative of
$\log\cot((\pi-s)/4)$ is $1/(2\cos(s/2))\geq1/2$.
Thus $H(w)-H_0\leq-(\pi-3)w^2/(4\pi^3)\leq0$.
Continuity and periodicity include the closed endpoints.

The final quadratic comparison also follows by integration. Since
$f'(0)=0$, the upper measure bound gives
$f(w)-f(0)\leq3w^2/(2\pi^2)$. The estimate for $W$ is
\[
 W(\pi-w)=\int_0^w\log\cot((\pi-s)/4)\,ds\geq w^2/4,
 \qquad 0\leq w\leq\pi.
\]
With $\lambda_0<1/(2\pi)$ and $J_0>1/\pi$, these yield the
negative coefficient $(3-\pi)/(4\pi^3)$ in $H(w)-H_0$.
This explains how the finite part $H$ is largest at zero although
$f$ itself reaches its largest value at the joined endpoint: curvature
and the cusp contribution enter the finite part as well.

The weight $W$ can be obtained explicitly from the triangular
identity. For a curvature contribution at displacement $s\geq0$,
its integrated weight is
\[
 K(s)=\int_s^\pi\frac{t-s}{2\sin(t/2)}dt.
\]
Differentiating in $s$ gives
$K'(s)=-\int_s^\pi dt/(2\sin(t/2))=-\log\cot(s/4)$,
and $K(\pi)=0$. Therefore
$K(s)=\int_s^\pi\log\cot(t/4)dt=W(s)$. This evaluates the
kernel produced by the change of integration order and explains the
specific logarithm in the convolution formula.
\end{proof}

\paragraph{The row profile.}
The variable $v$ gives the angular position of a row. The coefficients
$C(v),D(v)$ are its two cotangent amplitudes. Near an endpoint they
contain a term proportional to $a/v$. Averaging a cotangent peak along
one mesh produces a logarithm, and summing such row averages over the
other mesh produces another logarithm. This accounts for the leading
$L^2$ term of the completed model.

The coefficient $A(v)$ records the $L$ part of a single-row average.
Its constant part $\mathcal B(v)$ combines the two weighted scalar
pole profiles with their complete cancellation. An approximation is
called completed here because it includes the full infinite endpoint
reference and a specified subtracted integral for the constant. This
fixes the constant term as well as the leading growth. A bounded but
unspecified constant error would be too large for the strict corner
comparison later in the argument.

For $0<v<2\pi$ define the row coefficients and profile
\begin{align}\label{uc:row-profile}
 C(v)&=(p-q-2a\cot(v/2))/4,\quad D(v)=(p+q+2a\cot(v/2))/4,\notag\\
 \eta_-(v)&=\tfrac12+(x-y+v)/(2\pi),\quad
 \eta_+(v)=\tfrac12+(x+y-v)/(2\pi),\notag\\
 A(v)&=2(|C(v)|+|D(v)|)/\pi,\notag\\
 \mathcal B(v)&=\frac{c_e\Phi(\eta_-(v))+c_d\Phi(\eta_+(v))}
                         {2\sin(v/2)}+\frac{I(C(v),D(v),v)}\pi.
\end{align}
Write $\mathcal H(L,v)=A(v)L+\mathcal B(v)$ and
$U_0=c_ef(x-y)+c_df(x+y)$, $B_0=4a/\pi^2$,
$D_0=U_0/\pi+B_0\log\pi$. The completed model is
\begin{equation}\label{uc:completed-model}
 \begin{gathered}
 M_\infty(x,y,L)=\frac{4a}{\pi^2}L^2+b_{\rm full}L+C_\infty,\quad
 b_{\rm full}=\frac{2U_0}\pi+
 \frac{8a\log2+4a\log a+2(xp+yq)}{\pi^2},\\
 E_\infty=\frac a3+
 \frac{a}{\pi^2}\mathcal C(1/2+(x-y)/(2\pi))
                         \mathcal C(1/2+(x+y)/(2\pi)),\\
 C_\infty=E_\infty+\lim_{\epsilon\downarrow0}
 \left\{B_0\log^2\frac\epsilon\pi+2D_0\log\frac\epsilon\pi
            +\frac1\pi\int_\epsilon^{2\pi-\epsilon}\mathcal B(v)\,dv\right\}.
 \end{gathered}
\end{equation}
Weighted endpoint products and $a\log a$ have their continuous values.

The endpoint convention is visible in one variable:
$\mathcal C(t)\sim1/t$ and $\sin\pi t\sim\pi t$ as $t\downarrow0$,
so the weighted product tends to $\pi$. In $E_\infty$ the factor
$a=c_ec_d$ supplies the two needed cosine weights. The complete
product stays finite on an edge where an unweighted reciprocal factor
diverges. Similarly $a\log a\to0$ as $a\downarrow0$.

Two counterterms occur because the endpoint profile has both a
reciprocal and a logarithmic-reciprocal part. On each folded side its
leading expression is
\[
 \frac{U_0}{t}+\frac{4a}{\pi}\frac{\log t}{t},
\]
whose primitives are $U_0\log t$ and $(2a/\pi)\log^2t$.
There are two sides, and the integral has an outer factor $1/\pi$.
This gives the coefficients $B_0$ and $2D_0$ after changing the
normalization from $\log\epsilon$ to $\log(\epsilon/\pi)$.
In particular $B_0\log\pi$ inside $D_0$ is the cross term from
this change. The pieces of the finite part must be kept together
because this normalization fixes its constant.

At the origin, $a=c_e=c_d=1$, $p=q=0$, and
$U_0=2f(0)=2k_0/\pi$. Substitution in the model's linear coefficient
gives
\[
 b_{\rm full}(0)=\frac{4k_0+8\log2}{\pi^2}
 =\frac{8\{\gamma+\log(8/\pi)\}}{\pi^2}=b_c.
\]
This check ties the completed row normalization to the corner
normalization used later.

The limit exists uniformly on the closed phase triangle. We give an
explicit envelope, also needed for the finite mesh. For either folded
side $v_L(t)=t$, $v_R(t)=2\pi-t$, $0<t\leq\pi/4$,
\begin{equation}\label{uc:endpoint-envelope}
 \left|\mathcal B(v_s(t))-
       \frac{U_0+(4a/\pi)\log t}{t}\right|
 \leq7+\frac P\pi\left(3+2\log\frac{4\pi}{t}\right).
\end{equation}
\paragraph{Uniform control near an endpoint.}
A limit is uniform in phase if one cutoff works for every point of
the closed triangle. This is needed because a maximizing phase may
move with $n$ or approach an edge. The envelope just displayed gives
this control after the explicit endpoint singularity is removed. Its
right side is integrable, since
\[
 \int_0^\eta\log(A/t)\,dt=\eta\{1+\log(A/\eta)\}\longrightarrow0.
\]
The omitted short intervals therefore have contributions tending to
zero uniformly in phase. This proves that the completed constants
are uniformly Cauchy and later bounds the mesh intervals that contain
no retained sample.

The reference coefficients $(-w,w)$ preserve the large endpoint
amplitudes. Their cancellation is explicitly integrable. The phase
adds coefficient displacements of size at most $P/4$. Integrating a
bound involving the absolute cotangents over the whole interval would
lose finiteness at their poles. The proof first uses channel dominance
on short pole arcs and then uses the ordinary Lipschitz inequality for
absolute values on their complement.

To prove this, compare $I(C,D,t)$ with $I(-w,w,t)$, where
$w=a\cot(t/2)/2$ and each coefficient displacement is at most $P/4$.
For general displacements at most $\varepsilon$, remove pole arcs of
radius $t/8$. If $w\geq4\varepsilon$, the own channel dominates there
and the two cancellations differ by at most $8\varepsilon/7$ in
integral. If $w<4\varepsilon$, both complete cancellations are bounded
by twice the other channel, costing at most $72\varepsilon/7$.
On the complement their difference is at most
$2\varepsilon(|\cot z|+|\cot(z+t)|)$, whose integral is at most
$8\varepsilon\log(4\pi/t)$. Thus
\[
 |I(C,D,t)-I(-w,w,t)|\leq\varepsilon[11+8\log(4\pi/t)],
 \qquad I(-w,w,t)=4w\log\sin t.
\]
On the complement, changing each coefficient by at most $\varepsilon$
changes the complete cancellation by at most
$2\varepsilon(|\cot z|+|\cot(z+t)|)$, by two applications of
$||A|-|B||\leq|A-B|$. Deleting arcs of radius $t/8$ leaves a
logarithmic cotangent integral. On an arc, either the reference
coefficient is large enough to dominate the bounded other channel,
or all coefficients have size comparable to $\varepsilon$. This
explains the two cases $w\geq4\varepsilon$ and $w<4\varepsilon$
and their finite costs proportional to $\varepsilon$.

For the reference, the first lemma with $v=0$, $s=2w$ gives
$I(-w,w,t)=4w\log\sin t$. Replacing it by its singular part uses
$\cot(t/2)-2/t=O(t)$ and $\log(\sin t/t)=O(t^2)$.
Multiplication by the remaining reciprocal factor leaves integrable
errors of size $O(t(1+|\log t|))$. The numerator estimate below uses
the global Lipschitz bound for $f$, so it also works when a shifted
scalar argument crosses its cusp.

For the pole numerator $N_s$, $|N_s(t)-U_0|\leq12t/\pi$,
$|N_s|\leq8$. The sine chord gives a pole remainder at most
$6+\pi t/6$. Finally
$|2a\cot(t/2)\log\sin t-(4a/t)\log t|
\leq at(1+|\log t|)$. These estimates give \eqref{uc:endpoint-envelope}.
Its integrable right side makes the finite parts uniformly Cauchy.
The cancellation lemma bounds all central root transitions by an
integrable multiple of $1+|\log\sin t|$; hence their limits are
continuous in phase, including the axes and vertices.

On any truncated interval the complete cancellation depends
continuously on phase, including when equality roots merge. The first
lemma supplies continuous formulas and a common integrable bound.
Dominated convergence gives continuity there. The endpoint envelope
makes the rest uniformly small, proving continuity of the completed
constant on the closed triangle.

The next calculation independently checks the logarithmic coefficient.
The absolute coefficient sum is a piecewise maximum. Its transition
points can be found before integrating, and each piece has an elementary
primitive. Using $X=\max(x,y)$ and $Y=\min(x,y)$ organizes this scalar
calculation; the row roles of $x$ and $y$ are still retained in the
finite argument. At $a=0$ individual logarithmic primitives can look
singular, but the resulting expression has the specified continuous
$a\log a$ value. The integrable subtracted maximum provides the
justification for passage to that edge.

For clarity the logarithmic coefficient can be checked independently.
Put $X=\max(x,y)$, $Y=\min(x,y)$ and $T_t=2a\cot(t/2)$. Then
$\pi(A(t)+A(2\pi-t))$ is respectively
$2T_t$, $T_t+p+q$, $2\sin X$ on the intervals cut by
$\pi-X-Y$ and $\pi-X+Y$. Integrating with
$\int T_t\,dt=4a\log\sin(t/2)$ gives
\[
 \int_0^\pi\{A(t)+A(2\pi-t)-8a/(\pi t)\}\,dt
 =\{8a\log(2/\pi)+4a\log a+2(xp+yq)\}/\pi.
\]
Here is the piecewise integration. The identity
$|p-h|+|p+h|=2\max(p,|h|)$ gives
\[
 \pi A(t)=\max(p,q+T_t),\qquad
 \pi A(2\pi-t)=\max(p,|q-T_t|).
\]
Also $\sin X\geq\sin Y$, because
$\sin X-\sin Y=2\cos((X+Y)/2)\sin((X-Y)/2)\geq0$.
The changes of branch occur when $T_t=p+q$ and $T_t=|p-q|$.
The addition formulas put these respectively at
$t_1=\pi-X-Y$ and $t_2=\pi-X+Y$. First take $x+y<\pi$ and
$0<\epsilon<t_1$. Then, before subtracting
the reciprocal integral, the sum of the three pieces on $[\epsilon,\pi]$
is
\[
 \begin{aligned}
 &\pi\int_\epsilon^\pi\{A(t)+A(2\pi-t)\}dt\\
 &=2\int_\epsilon^{t_1}T_tdt+\int_{t_1}^{t_2}T_tdt\\
 &\quad +(p+q)(t_2-t_1)+2\sin X(\pi-t_2)\\
 &=4a\log a-8a\log\sin(\epsilon/2)+2(xp+yq).
 \end{aligned}
\]
Here $\sin(t_1/2)=c_e$, $\sin(t_2/2)=c_d$, and
$(p+q)2Y+2\sin X(X-Y)=2(xp+yq)$. Subtract
$8a\int_\epsilon^\pi dt/t$, then let $\epsilon\downarrow0$ using
$\sin(\epsilon/2)/\epsilon\to1/2$. This gives the displayed
integral and its logarithmic coefficient.

This proves the coefficient in \eqref{uc:completed-model}.
The max formulas bound each subtracted integrand by one, so dominated
convergence includes $a=0$.

\paragraph{Keeping phase losses together.}
The leading coefficient $4a/\pi^2$ is largest at the origin. The
other coefficients may increase away from zero, so the useful question
is how much they can increase at the same phase. The first estimate
in the next lemma bounds the linear increase by $3u/\pi^2$. The
quadratic and linear terms together therefore lose at least
\[
 \frac{4uL^2-3uL}{\pi^2}=uT.
\]
The constant adds a further loss $ju$, with either a possible sine gain
or the explicitly isolated central term $E_c$.

These bounds are called correlated because $a,c_e,c_d,p,q$ all come
from the same two phases. Their separate largest possible values need
not occur together. Retaining their exact relations gives losses that
would disappear if each factor were bounded separately. The
nonnegative quantity $z$ stores one such additional loss.

\begin{lemma}[Correlated model bounds]\label{uc:master}
Put
\[
 \begin{gathered}
 b_*=C_0-7/2,\quad B=C_0^2/\pi^2,\quad d_0=C_0b_*/\pi^2,
 \quad v_0=7b_*/(2\pi^2),\\
 g_0=4(\log^2\pi-\log^22)/\pi^2,\quad
 j=H_0+d_0-g_0,\quad t_0=2\log2/\pi.
 \end{gathered}
\]
With $C_*=C_\infty(0,0)$ and $b_c=8(\gamma+\log(8/\pi))/\pi^2$,
\begin{align}\label{uc:model-loss}
 b_{\rm full}-b_c&\leq3u/\pi^2,& b_{\rm full}&<79/(5\pi^2),\notag\\
 C_\infty-C_*&\leq-ju+t_0p-(H_0+v_0)z\leq-ju+t_0p,
 &C_\infty&<9/4.
\end{align}
On $\rho\leq\pi/2$ there is the exact decomposition
\begin{equation}\label{uc:central-decomposition}
 C_\infty=E_\infty+c_eH(x-y)+c_dH(x+y)-(1/3+g_0)a+E_c,
\end{equation}
where
\begin{equation}\label{uc:Ec}
 \begin{aligned}
 E_c=\frac1{\pi^2}\int_0^{\pi/2}
 \bigl\{&J(p/2,a\tan(t/2)+q/2,t)\\
 &+J(p/2,|a\tan(t/2)-q/2|,t)-4a\tan(t/2)\log\sin t\bigr\}\,dt.
 \end{aligned}
\end{equation}
Consequently
\begin{equation}\label{uc:regular-loss}
 (C_\infty-E_c)-C_*\leq-ju-(H_0+v_0)z\leq0.
\end{equation}
\end{lemma}
\begin{proof}
First the reciprocal-pair series implies, for $|\theta|\leq\pi$,
\begin{equation}\label{uc:weighted-C}
 \cos(\theta/2)\mathcal C(1/2+\theta/(2\pi))
       \leq C_0-b_*\{1-\cos(\theta/2)\}.
\end{equation}
\paragraph{Obtaining the infinite cotangent formula from a finite identity.}
We give the series justification used in this proof and in the residual
bounds below. The finite sine product is proportional to $\sin z$:
\[
 \prod_{k=0}^{N-1}\sin((z+k\pi)/N)=c_N\sin z,
 \qquad c_N\ne0\ \text{independent of }z.
\]
To check proportionality, write each sine using
$\sin u=(e^{\mathrm i u}-e^{-\mathrm i u})/(2\mathrm i)$.
With $w=e^{2\mathrm iz/N}$ and $\zeta=e^{2\pi\mathrm i/N}$,
the finite polynomial identity
$\prod_{k=0}^{N-1}(1-w\zeta^k)=1-w^N$ leaves a constant times
$e^{-\mathrm iz}(e^{2\mathrm iz}-1)$, which is a constant times
$\sin z$. The polynomial identity follows because both sides have
the same $N$ roots and the same constant term. Here
$\mathrm i^2=-1$; this finite algebra is the only use of complex
exponentials in the argument. Differentiating the logarithm of the
absolute product on any real interval avoiding its zeros gives
\[
 \cot z=\frac1N\sum_{k=0}^{N-1}\cot((z+k\pi)/N).
\]
Take odd $N=2M+1$. Pair the residues $k$ and $N-k$, using cotangent's
period $\pi$, to write
\[
 \cot z=\frac1N\cot(z/N)
   +\frac1N\sum_{k=1}^M
       \{\cot((z+k\pi)/N)+\cot((z-k\pi)/N)\}.
\]
For fixed $k$, the paired term tends to
$1/(z+k\pi)+1/(z-k\pi)$. To control the whole tail, fix
$|z|\leq R<\pi$. For $1\leq k\leq M$, the absolute sine
arguments $(k\pi\pm z)/N$ lie between zero and $2\pi/3$.
The chord bound $\sin u\geq u/\pi$ on that interval gives
\[
 \begin{aligned}
 \left|\frac1N\{\cot((z+k\pi)/N)+\cot((z-k\pi)/N)\}\right|
 &=\frac{|\sin(2z/N)|}
       {N|\sin((z+k\pi)/N)\sin((z-k\pi)/N)|}\\
 &\leq\frac{2\pi^2R}{k^2\pi^2-R^2}\leq\frac{C_R}{k^2}.
 \end{aligned}
\]
The constant $C_R=2\pi^2R/(\pi^2-R^2)$ is independent of $N,k$.
Thus the tail after $K$ is at most $C_R\sum_{k>K}k^{-2}\leq C_R/K$.
Pass first $N\to\infty$ in the finite initial part, then
$K\to\infty$ in this uniform tail bound. This proves
\[
 \cot z=\frac1z+\sum_{k\geq1}
       \left(\frac1{z+k\pi}+\frac1{z-k\pi}\right)
 =\frac1z-2z\sum_{k\geq1}\frac1{\pi^2k^2-z^2},\qquad |z|<\pi.
\]
After subtracting $1/z$ the formula also has its continuous value at
zero. On compact subintervals the rational series and its needed
derivatives have summable bounds, which justify differentiation.
For example, on $|z|\leq\pi/2$, each denominator is at least
$3\pi^2k^2/4$. Comparing the coefficient of $z$ at zero with
$\cot z=1/z-z/3+O(z^3)$ also gives
$\sum_{k\geq1}k^{-2}=\pi^2/6$. Thus the series identities here
are consequences of finite products and elementary tail estimates.

Here is a direct verification with its endpoint normalization.
For $0<s<1$, $c=\cos(\pi s/2)$, the tangent partial fraction gives
\[
 \frac{d}{ds}\log\frac{1-s^2}{c}
 =\sum_{\ell\geq3,\ \ell\ {\rm odd}}\frac{2s}{\ell^2-s^2}
 \geq2s(\pi^2/8-1).
\]
For clarity, the partial fraction used here is
\[
 \frac\pi2\tan\frac{\pi s}{2}
 =\sum_{\ell\geq1,\ \ell\ \mathrm{odd}}\frac{2s}{\ell^2-s^2}.
\]
It follows by substituting the cotangent partial fraction in
$\tan z=\cot z-2\cot(2z)$; the remaining series is absolutely
convergent on compact subsets of $|s|<1$. Differentiating at zero
gives $\sum_{\ell\ \mathrm{odd}}\ell^{-2}=\pi^2/8$.
Put $Q_*(s)=(1-s^2)/c$ and $k_*'=\pi^2/8-1>0$. Integrating
$(\log Q_*)'\geq2k_*'s$ gives $Q_*\geq e^{k_*'s^2}\geq1+k_*'s^2$.
Equivalently
\[
 1-c\geq\frac{(1+k_*')s^2}{1+k_*'s^2},\qquad
 \frac{cs^2}{(1-c)(1-s^2)}\leq\frac1{1+k_*'}=\frac8{\pi^2}.
\]
Explicitly,
\[
 \frac{cs^2}{(1-c)(1-s^2)}
 =\frac{s^2}{Q_*-1+s^2}\leq\frac1{1+k_*'}.
\]
The denominator is positive for $s>0$. This supplies the
rearrangement behind the short integration step.

Integration yields $cs^2/[(1-c)(1-s^2)]\leq8/\pi^2$.
Also $\sum_{\ell\geq3,\ \ell\ {\rm odd}}1/[\ell(\ell^2-1)]
=3/4-\log2$, by partial fractions. Therefore

The summand identity is
\[
 \frac1{\ell(\ell^2-1)}
 =\frac12\left(\frac1{\ell-1}+\frac1{\ell+1}\right)-\frac1\ell.
\]
Summing for $\ell=3,5,\ldots,2N+1$ and using
$H_{2N}-H_N\to\log2$ gives $3/4-\log2$. In the reciprocal
increment the $\ell=1$ term is $1/(1-s^2)$, while all remaining
terms are at most this constant sum. The just-proved cosine inequality
bounds their weighted prefactor by $8/\pi^2$; hence
\[
 \frac{c\{\mathcal C(1/2+s/2)-C_0\}}{1-c}
 \leq\frac{32}{\pi^2}\{1+3/4-\log2\}.
\]
This displays separately the endpoint pole and its convergent tail.
This gives the following displayed estimate.
\[
 \frac{c\{\mathcal C(1/2+s/2)-C_0\}}{1-c}
 \leq\frac{32}{\pi^2}(7/4-\log2)<7/2.
\]
The endpoint weighted increment tends to $\pi<7/2$; symmetry and
continuity prove \eqref{uc:weighted-C} on the closed interval.
Here is the multiplication explicitly. Put $h_e=1-c_e$ and
$h_d=1-c_d$. The positive weighted reciprocal factors are bounded by
$C_0-b_*h_e$ and $C_0-b_*h_d$. Their product divided by $\pi^2$
is
\[
 B-d_0(h_e+h_d)+\frac{b_*^2}{\pi^2}h_eh_d.
\]
Use $h_e+h_d=u+z$, $h_eh_d=z$, and $C_0-b_*=7/2$. The coefficient
of $z$ becomes $-b_*(C_0-b_*)/\pi^2=-v_0$. Adding $a/3$
yields the next bound. Keeping the product before simplifying is what
preserves this negative $z$ term.

Multiplication of its positive factors gives
\begin{equation}\label{uc:endpoint-loss}
 E_\infty\leq a/3+B-d_0u-v_0z.
\end{equation}
If $D=\mathcal C(1/2+(x+y)/(2\pi))+
\mathcal C(1/2+(x-y)/(2\pi))-2C_0$, then $aD\leq(7/2)u$.
Substitution in $b_{\rm full}$ gives exactly
\[
 b_{\rm full}-b_c=\frac2{\pi^2}
 \{aD+2a\log a+xp+yq-4(\gamma+\log(8/\pi))u\}.
\]
The logarithm estimates can be obtained from an elementary positive
series. For $|r|<1$ define
\[
 \operatorname{atanh}r=\frac12\log\frac{1+r}{1-r}
 =\sum_{k\geq0}\frac{r^{2k+1}}{2k+1}.
\]
The equality follows by integrating the geometric series for
$1/(1-r^2)$ from zero to $r$. With $r=17/39$ the logarithm is
$\log(28/11)$, so retaining its first two positive terms gives the
lower bound invoked below. With $r=1/3$ the same formula gives
$\log2$, and three terms give $\log2>69/100$.

Use $xp+yq\leq4u$, $a\log a\leq0$, and
$3/2<\gamma+\log(8/\pi)<8/5$. For the lower bound the first two
terms of $2\operatorname{atanh}(17/39)$ give
$\log(28/11)>37/40$; the remaining constants follow from
$23/40<\gamma<3/5$, $3<\pi<22/7$. This proves both coefficient bounds.

The bound $xp+yq\leq4u$ also has a direct trigonometric proof.
For $0\leq t\leq\pi$,
\[
 t\sin t\leq2(1-\cos t),
\]
because after factoring $2\sin(t/2)$ it becomes
$t\cos(t/2)\leq2\sin(t/2)$, which follows from
$\tan(t/2)\geq t/2$. Summing for $x,y$ gives $xp+yq\leq4u$.
Consequently the linear coefficient is bounded by a loss quadratic
near zero. In its exact difference formula, the bounds
$aD\leq(7/2)u$, $xp+yq\leq4u$, $a\log a\leq0$, and
$\gamma+\log(8/\pi)>3/2$ leave at most
$(2/\pi^2)(3u/2)=3u/\pi^2$.

The primitive below uses the function
$\operatorname{Li}_2(s)=\sum_{k\geq1}s^k/k^2$ on $[0,1]$.
Differentiating its series for $0<s<1$ gives
$\operatorname{Li}_2'(s)=-\log(1-s)/s$, and its endpoint value
is $\operatorname{Li}_2(1)=\pi^2/6$. These identities will verify
the primitive and its endpoint constant.

Split the completed finite part into its scalar pole term
$c_eH(x-y)+c_dH(x+y)$ and cancellation term. The cancellation
parameters at $0<t<\pi$ are
$J(p/2,(q+T_t)/2,\epsilon)$ and
$J(p/2,|q-T_t|/2,\epsilon)$, $\epsilon=\min(t,\pi-t)$.
Since $R\leq0$, their sum is at most
\[
 4a\cot(t/2)\log\sin t+2p|\log\sin t|.
\]
Indeed their logarithmic coefficient is
$(q+T_t-p)_++(|q-T_t|-p)_+\geq2T_t-2p$.
The integral of $|\log\sin t|$ is $\pi\log2$. A primitive of the
first term is
$4a\log^2(2\sin(t/2))-4a\log^22-2a\operatorname{Li}_2(\sin^2(t/2))$.
Since $\operatorname{Li}_2(1)=\sum k^{-2}=\pi^2/6$, its endpoint
counterterm leaves $-(1/3+g_0)a$.
The sign of $\log\sin t$ is nonpositive. A lower bound on its
coefficient therefore provides an upper bound on the product. For
$T\geq0$, the positive parts used here satisfy
\[
 (q+T-p)_++(|q-T|-p)_+\geq2T-2p.
\]
Each positive part is at least its argument, and
$q+T+|q-T|\geq2T$. Together with $R\leq0$, this proves the stated
cancellation envelope. Integrating its possible gain uses
$\int_0^\pi|\log\sin t|dt=\pi\log2$.

\paragraph{A cutoff calculation for the completed constant.}
Write $I_L(t)+I_R(t)$ for the two folded cancellation integrals.
The pole-profile part of the completed integral is
$c_eH(x-y)+c_dH(x+y)$: insert the two shifted $f$ numerators in the
definition of $H$, reflect the second integration variable, and use
its periodicity. Their common endpoint value is $U_0$, so their
counterterm is $(2U_0/\pi)\log(\epsilon/\pi)$.
Subtracting this counterterm from the one in $C_\infty$ leaves
\[
 B_0\log^2(\epsilon/\pi)+2B_0\log\pi\log(\epsilon/\pi)
 =B_0(\log^2\epsilon-\log^2\pi).
\]
Thus the remaining cancellation finite part is exactly
\[
 \lim_{\epsilon\downarrow0}
 \left\{B_0(\log^2\epsilon-\log^2\pi)
     +\frac1{\pi^2}\int_\epsilon^\pi(I_L(t)+I_R(t))dt\right\}.
\]
For the explicit comparison integrand
$4a\cot(t/2)\log\sin t$, its primitive at $\pi$ is
$-a\pi^2/3$, and at $\epsilon$ is
$4a\log^2\epsilon-4a\log^22+o(1)$. Adding the counterterm
therefore leaves
\[
 -\frac a3+\frac{4a}{\pi^2}(\log^22-\log^2\pi)
 =-(1/3+g_0)a.
\]
On $\rho\leq\pi/2$ the endpoint half equals this comparison
integrand exactly. On the other half use $t\mapsto\pi-t$;
then $\cot((\pi-t)/2)=\tan(t/2)$, and the two actual coefficient
parameters become $q/2+a\tan(t/2)$ and
$|q/2-a\tan(t/2)|$. The difference from the comparison integrand
is exactly the integral defining $E_c$. This proves the central
decomposition with its prescribed finite-part normalization.

The finite-part counterterms cancel the endpoint logarithms in this
primitive, leaving $-(1/3+g_0)a$.

On $\rho\leq\pi/2$, $2a\geq p+q$, and the endpoint half is
root-free: $(2a\cot(t/2)-q)\sin t\geq2a-q\geq p$ for
$0<t\leq\pi/2$. Reflecting the central half gives exactly
\eqref{uc:central-decomposition}. The cancellation lemma bounds its
integrand by $10|\log\sin t|+4\log2$, including equality at the
closed boundary. Dominated convergence justifies all boundary limits.
Apply $H\leq H_0$ and \eqref{uc:endpoint-loss}, then use
$2-c_e-c_d=u+z$. At zero this gives
$C_*=B-g_0+2H_0$ and proves both correlated losses.

On the central phase region, the endpoint half of the cancellation
has no equality roots: its parameters satisfy the root-free inequality
displayed in the proof. Its integral is exact. Reflecting the other
half isolates all remaining root motion in $E_c$. The ordinary regular
part can then be estimated by $H\leq H_0$ and the weighted product
bound, giving
\[
 (C_\infty-E_c)-C_*
 \leq-(H_0+d_0-g_0)u-(H_0+v_0)z.
\]
This explains the definition of $j$. At zero, $p=q=0$ and
$J(0,s,t)=2s\log\sin t$, so the two terms in $E_c$ cancel its
explicit subtraction. Thus $E_c(0,0)=0$ and the origin constant is
$C_*=B-g_0+2H_0$. The identity obtained next equates it with the
independently normalized corner expansion. The finite error estimate
proved in the next subsection tends to zero and justifies this
identification; that estimate does not require the numerical bounds
for $H_0$ that are being deduced here.

The corner expansion identifies the same constant as
$C_*=4(\gamma+\log(8/\pi))^2/\pi^2+\kappa$:
at zero $S_n=\Lambda_c(n)$ and the finite error proved below tends
to zero. Hence
\begin{equation}\label{uc:origin-identity}
 H_0=\kappa/2+(2k_0/\pi^2)\log(2/\pi).
\end{equation}
The coarse certified input
\begin{equation}\label{uc:coarse-input}
 9/10<\kappa<23/25,\qquad C_*<1847/1000
\end{equation}
is needed here. We specify its finite reduction. Define
\begin{align*}
 \mathfrak a_0&=1,\qquad
 \mathfrak a_p=1-2\sum_{r=1}^p\frac{\zeta(2r)}{4^r}\quad(p\geq1),\\
 \mathfrak b_0&=7\zeta(3)/8,\\
 \mathfrak b_q&=(1-2^{-2q-3})\zeta(2q+3)
              -(1-2^{-2q-1})\zeta(2q+1)\quad(q\geq1),\\
 c_k&=\sum_{p+q=k}\mathfrak a_p\mathfrak b_q.
\end{align*}
Cancel the endpoint pole before multiplying the two power series:
\begin{equation}\label{uc:kappa-series}
 \begin{gathered}
 \frac{t^2\cot(\pi t/2)}{1-t^2}
       =\frac{2t}\pi\sum_{p\geq0}\mathfrak a_pt^{2p},\qquad
 (1-t^2)\sum_{\ell\geq1,\ \ell\text{ odd}}
       \frac1{\ell(\ell^2-t^2)}=\sum_{q\geq0}\mathfrak b_qt^{2q},\\
 \kappa=\frac8{\pi^2}\sum_{k\geq0}\frac{c_k}{k+1},\qquad
 \left|\kappa-\frac8{\pi^2}\sum_{k=0}^K\frac{c_k}{k+1}\right|
       \leq\frac8{3(K+2)4^{K+1}}.
 \end{gathered}
\end{equation}
Indeed the cotangent partial fraction
$\cot(\pi t/2)=2/(\pi t)-(4t/\pi)\sum_{j\geq1}(4j^2-t^2)^{-1}$
and $\sum_{j\geq1}(4j^2-1)^{-1}=1/2$ give the first series and
$0<\mathfrak a_p\leq4^{-p}$. Reciprocal expansion gives the second;
its $\ell=1$ pole cancels, $\mathfrak b_0<10/9$, and for $q\geq1$,
$|\mathfrak b_q|\leq9^{-(q-1)}(1/27+\int_3^\infty s^{-3}ds)<9^{-q}$.
Consequently $|c_k|\leq4^{-k}(10/9+\sum_{q\geq1}(4/9)^q)
<2\cdot4^{-k}$. These bounds give uniform absolute convergence on
$[0,1]$, so substitution into \eqref{eq:main-kappa} and termwise
integration prove \eqref{uc:kappa-series}; its tail follows by summing
the geometric majorant, using $8/\pi^2<1$.
\paragraph{Reducing the constant enclosure to a finite sum.}
The factor $1-t^2$ removes the $\ell=1$ pole before multiplying
the series. Their resulting geometric majorants give uniform absolute
convergence on $[0,1]$. Thus the product can be expanded and integrated
term by term. The denominator $k+1$ comes from integrating $t^{2k+1}$
together with the prefactors of the original integral.

For the tail, $|c_k|<2\cdot4^{-k}$. When $k\geq K+1$, replace
$1/(k+1)$ by $1/(K+2)$ and sum the geometric majorant. With
$8/\pi^2<1$, this yields the stated exact tail radius. A finite sum
computed with rigorous lower and upper arithmetic bounds, enlarged by
that radius, therefore encloses the infinite constant.

Outward arithmetic means rounding every lower endpoint downward and
every upper endpoint upward at each operation. It produces an interval
containing the exact value. The inherited certificate described next
supplies this finite arithmetic input. The series and tail explain its
mathematical reduction; this paragraph asserts no new computation.

The inherited certificate evaluates this explicit finite sum through
$K=64$ by outward Arb arithmetic at $320$ bits and adds the exact tail
radius $8/(3\cdot66\cdot4^{65})$. It also evaluates
$C_*=4(\gamma+\log(8/\pi))^2/\pi^2+\kappa$ outwardly. Its recorded
enclosures and independent verification imply \eqref{uc:coarse-input};
the reproducibility record describes the arithmetic receipt. This input
encloses the limiting constants used in the global comparison.
Using $161/100<k_0<9/5$, $\log(\pi/2)<23/50$ and
$\log(\pi/2)>\log(25/16)>18/41$ gives
\begin{equation}\label{uc:scalar-constants}
 \begin{gathered}
 7/25<H_0<1633/5125<8/25,\quad d_0>39/250,\\
 g_0<43/125,\quad j>23/250,\quad 0<t_0<221/500.
 \end{gathered}
\end{equation}
For example $H_0>9/20-828/4925>7/25$;
$C_0>39/10$, $b_*>2/5$, $\pi^2<10$ give the $d_0$ bound;
and $g_0<(80/197)((23/20)^2-(69/100)^2)<43/125$.
For the last constant use $\log2<6932/10000$, $\pi>157/50$.
The positive bound for $j$ supplies a strict regular loss. The
coefficient $t_0$ bounds the possible sine gain in the global constant
estimate, and the enclosure for $C_*$ makes the universal $9/4$ bound
quantitative. For that last maximization, Cauchy--Schwarz gives
\[
 -ju+t_0\sin x
 \leq-\frac j2+\frac j2\cos x+t_0\sin x
 \leq-\frac j2+\sqrt{j^2/4+t_0^2}.
\]
The vector $(\cos x,\sin x)$ has length one, which explains the
square root. Writing $\delta=9/4-C_*>0$, the desired inequality is
equivalent, after squaring positive quantities, to
$\delta^2+\delta j>t_0^2$. This is exactly the rational comparison
made at the end of the proof.

Finally $u\geq(1-\cos x)/2$ gives
$C_\infty\leq C_*-j/2+\sqrt{j^2/4+t_0^2}$.
With $\delta=9/4-C_*>403/1000$, one has
$\delta^2+\delta j>(403/1000)^2+(403/1000)(23/250)
>(221/500)^2$. Squaring after adding $j/2>0$ proves $C_\infty<9/4$.
\end{proof}

\paragraph{Why the central power is three halves.}
The central root threshold near zero has shape $p\approx at^2$,
because $a\tan(t/2)\sin t=a(1-\cos t)$. Root activity is therefore
confined to a region with width comparable to $\sqrt p$ and height
comparable to $p$. Their product suggests $p^{3/2}$. The next lemma
makes this precise while allowing roots to merge. It also retains the
favorable $y$ term $B_\beta$.

The number $b$ is the point where $a\tan(t/2)=q/2$. On $[0,b]$
the difference $q/2-a\tan(t/2)$ is nonnegative. Since
$-\log\sin t\geq0$, $B_\beta\geq0$. The identity at $x=0$
will show that this term is an exact boundary loss. It is particularly
useful along the $y$ axis, where the positive $p^{3/2}$ term vanishes.

\begin{lemma}[The paired central roots]\label{uc:central-pair}
On $x,y\geq0$, $\rho\leq\pi/2$, put
\[
 b=2\arctan(q/(2a)),\qquad
 B_\beta=\frac4{\pi^2}\int_0^b
       (q/2-a\tan(t/2))[-\log\sin t]\,dt.
\]
Then $0\leq b\leq\rho$, $B_\beta\geq0$, $E_c(0,y)=-B_\beta$, and
\begin{equation}\label{uc:central-pair-general}
 E_c\leq\frac{4I_4}{\pi^2}\frac{p^{3/2}}{\sqrt a}-B_\beta,
 \qquad I_4=\int_0^1\sqrt{1-t^4}\,dt<319/360.
\end{equation}
In particular
\begin{equation}\label{uc:central-pair-bounds}
 E_c\leq\begin{cases}3x^{3/2}/8-B_\beta,&\rho\leq1/2,\\
                       5x^{3/2}/12-B_\beta,&\rho\leq1.
       \end{cases}
 \qquad B_\beta\geq1679y^2/8000\quad(\rho\leq1/2).
\end{equation}
\end{lemma}
\begin{proof}
The phase identity $(\cos x+\cos y)\tan(\rho/2)=p+q$ gives
$b\leq\rho$. Using \eqref{uc:root-Z}, the exact integrand rearranges
as two $E$ terms minus $B_\beta$. Put $v=p/2$, $\delta=q/2$.
To see the subtraction, let $s_+=a\tan(e/2)+q/2$ and
$s_-=|a\tan(e/2)-q/2|$. Then
\[
 s_++s_-=2a\tan(e/2)+2(q/2-a\tan(e/2))_+.
\]
Replacing each $J(v,s_\pm,e)$ by
$E(v,s_\pm,e)+2s_\pm\log\sin e$ in $E_c$ leaves two $E$
terms and $4(q/2-a\tan(e/2))_+\log\sin e$. The integral of
the latter, divided by $\pi^2$, is $-B_\beta$. Thus $E\leq Z$
reduces the remaining task to a sum of nonnegative root integrals.

Tonelli and $w=a\tan(e/2)$ give for the sum of their $Z$ majorants
\[
 2\int_0^{\pi/2}\frac1{\sin t}\int_{\mathbb R}
 \frac{2a\,\mathbf1_{|w|\leq a\tan(t/2)}}{a^2+w^2}
        \sqrt{(v^2-(w-\delta)^2\sin^2t)_+}\,dw\,dt.
\]
The substitution $w=a\tan(e/2)$ has Jacobian
$de=2a\,dw/(a^2+w^2)$. The condition $e<t$ in the definition
of $Z$ becomes $w\leq a\tan(t/2)$. Combining the two choices
$s_\pm$ extends the $w$ integral to positive and negative values
with a translation by $\delta=q/2$. The factors are a truncated
Cauchy weight $2a/(a^2+w^2)$ and a square root with support in an
interval. Before translation both are even and decrease with $|w|$.
Their nonnegative superlevel sets are centered intervals. The same
intersection-length argument used for the circular convolution shows
that their product integral is largest when the centers coincide.
This simultaneously bounds both central roots while preserving the
subtracted $B_\beta$.

Both inner factors are even and radially decreasing. Layer cake for
interval intersections shows that translation by $\delta$ cannot
increase the integral. It is therefore at most
$2F(v,a)$, where $F(v,a)=\int_0^{\pi/2}Z(v,a\tan(e/2),e)\,de$.
For $\tau=v/a\leq1$, a second Tonelli interchange, preserving
$e<t$, gives
\[
 F=4v\tau\int_0^1\sqrt{1-r^2}\,K(\tau r)\,dr,
 \quad K(k)=\int_{\arccos(1-k)}^{\pi/2}\frac{dt}{\sin^2t+k^2}
 \leq\frac1{\sqrt{2k}}.
\]
The lower integration limit in $K$ can be checked directly. In the
unshifted root integral put $w=vr/\sin t$, $0\leq r\leq1$.
The inequality $w\leq a\tan(t/2)$ becomes
\[
 \tau r\leq\sin t\tan(t/2)=1-\cos t.
\]
Hence $t\geq\arccos(1-\tau r)$. The Jacobian and the outside
factor $1/\sin t$ combine into $\sin^2t+\tau^2r^2$ in the
denominator. This preserves the condition $e<t$ while changing the
order of integration, explaining the particular integral $K$.

Indeed $w=\cot t$ bounds $K(k)$ by
$(1-k)/\sqrt{k(2-k)}\leq1/\sqrt{2k}$, using
$2(1-k)^2\leq2-k$. Thus
$F\leq4\sqrt2 I_4v^{3/2}/\sqrt a$,
$I_4=\int_0^1\sqrt{1-r^4}\,dr<319/360$; the latter follows by
squaring $1-u/2-u^2/8\geq\sqrt{1-u}$ and integrating $u=r^4$.
The resulting bound for $E_c+B_\beta$ is
$4I_4p^{3/2}/(\pi^2\sqrt a)$.

The change $r=s^2$ gives
\[
 \int_0^1r^{-1/2}\sqrt{1-r^2}\,dr=2I_4.
\]
Also $1-u/2-u^2/8$ is nonnegative for $0\leq u\leq1$, and its
square minus $1-u$ is $u^3/8+u^4/64\geq0$. Substituting $u=s^4$
and integrating bounds $I_4$ by
$1-1/10-1/72=319/360$, strictly since the pointwise inequality is
strict on an interval. Finally $v=p/2$ converts the bound for $F$
into the coefficient stated for $E_c+B_\beta$.
For $\rho\leq1/2$, $a\geq15/16>(29/30)^2$; for $\rho\leq1$,
$a\geq49/64$. Together with $\pi^2>197/20$ these give $3/8$
and $5/12$, respectively, and $p\leq x$. Finiteness also follows
from $Z\leq2v\log(2/e)$, so zero parameters are included.
At $v=0$, $E=Z=0$ and the asserted boundary identity is exact.
Finally $2a\leq1+\cos y$ gives $b\geq y$.
Convexity of tangent implies $a\tan(t/2)\leq qt/(2b)$ on $[0,b]$.

A convex function lies below the chord joining two points of its
graph. Here those points are $(0,0)$ and $(b,q/2)$ for
$t\mapsto a\tan(t/2)$. Since $\sin t\leq t$, this gives
\[
 B_\beta\geq\frac{2q}{\pi^2}\int_0^b(1-t/b)(-\log t)dt
 =\frac{qb}{\pi^2}\{\log(1/b)+3/2\}.
\]
For the integral, set $t=br$ and use
$\int_0^1(1-r)dr=1/2$ and
$\int_0^1(1-r)(-\log r)dr=3/4$. The inequalities $b\geq y$,
$b\leq1/2$, and $q\geq(23/24)y$ on the near domain then give
the rational quadratic loss displayed in the statement.
For $b\leq1/2$,
$B_\beta\geq qb(\log(1/b)+3/2)/\pi^2
\geq(23/24)y^2(69/100+3/2)/10=1679y^2/8000$.
\end{proof}

\subsection{One finite-row telescope}
\label{uc:finite-comparison}

\paragraph{From a continuous model to the actual board.}
The model has now been defined on the closed triangle, with explicit
phase losses. We next compare it with $S_n$, the finite sum introduced
above. The comparison will use one exact decomposition and several
estimates for its terms. Every estimate has a particular purpose: row
errors measure sampling along a row, outer errors measure sampling
across rows, retained-row errors protect special endpoint and central
rows, and mesh terms record the mismatch between finite harmonic sums
and their completed values.

There are two kinds of final bound. An absolute bound on
$|S_n-M_\infty|$ is useful on the whole triangle. Near the origin,
a bound for the difference between the two finite errors is much
stronger. The same exact identity will produce both bounds. Keeping
the error terms separate preserves the stronger comparison near the origin.

The integer restriction on $n$ describes the actual board. Allowances
are also viewed as functions of a real variable $n\geq32$ so their
monotonicity can be checked by differentiation. This use of real $n$
only concerns the explicit scalar bounds, never a nonintegral board.

Hereafter $n\geq32$. All allowances are explicit scalar functions of
real $n$; the interpolation degree remains integral. Define once
\begin{align}\label{uc:constants}
 T&=(4L^2-3L)/\pi^2,\quad b_0=4/\pi^2,\quad r_0=\log(5/2)+1/18,\notag\\
 D&=\frac{(52/\pi^2)H+7}{n}+\frac{10}{n^2}+\frac{20H}{\pi n^2},\notag\\
 \mathcal M&=\frac{b_0(2L+2r_0)}n
 +\frac{4+b_0(L+1/n)}{4n^2-1}+\frac{b_0}{(2n-1)^2}\notag\\
 &\qquad+\frac{5+b_0(\log(n+1/2)+2)}{4(n-1/2)^2},\notag\\
 \bar e&=D+\mathcal M+1/n,
 \quad K_D=\frac{[6+16/(3\pi)]H+2/\pi}{n},\notag\\
 A_4&=4/n+40H/(\pi n^2),\quad C_4=A_4+20/n^2,
 \quad d_\beta=(17L+27)/n^2,\notag\\
 J_N&=((3/2)H+16/5)/n+78/n^2,
 \quad J_W=((3/2)H+6)/n+148/n^2,\notag\\
 k_{Nx}&=((19/2)H+(23/3)L+677/10)/n+82/n^2,
 \quad k_{Ny}=k_{Nx}+d_\beta,\notag\\
 k_{Wx}&=((19/2)H+(23/3)L+155/2)/n+152/n^2,
 \quad k_{Wy}=k_{Wx}+d_\beta,\notag\\
 U_N&=T-\bar e+A_4+\rho J_N+j,
 \quad U_W=T-\bar e+A_4+\rho J_W+j.
\end{align}
The constants in the first display are a list of allowances that will
be assembled at the end. Their meanings can be read in groups.
The quantity $T$ is the quadratic-logarithmic model loss. The quantity
$D$ bounds the row contribution at the origin, and $\mathcal M$
bounds its mesh correction. The added $1/n$ in $\bar e$ allows for
the remaining one-sided origin error. The term $K_DP$ accounts for
ordinary coefficient perturbations away from the symmetric reference.

The symbols $N,W$ refer to the two nested domains $\rho\leq1/2$
and $\rho\leq1$. On the smaller domain, denominators of the retained
rows stay farther from zero, giving the smaller constants $J_N$ and
$k_{Nx},k_{Ny}$. The larger domain covers phases that must later be
excluded with a different margin. The expressions $U_N,U_W$ collect
the usable quadratic phase loss after the finite corrections are paid.
No optimization problem has been replaced by these constants: they
are explicit upper or lower estimates derived below.

For the phase-dependent absolute allowance put
\begin{align}\label{uc:phase-allowance}
 D_a&=((52/\pi^2)H+7)/n+20H/(\pi n^2),\quad
 D_\phi=aD_a+10/n^2,\notag\\
 \epsilon&=3\pi/(2m),\quad \ell=-\log\sin\epsilon,\quad
 \overline s=1/2+2a/\epsilon,\quad
 \Gamma_\phi=p[1+\log(1+2\overline s/p)],\notag\\
 M_A^\phi&=\max(p,q)/\pi+4a/\pi^2,\quad V_A^\phi=2M_A^\phi,
 \quad C_B^\phi=6+(4+4a/\mathrm e)/\pi,\notag\\
 V_B^\phi&=23+\{3\Gamma_\phi+(2p+6P)\ell+26a+20+2d_+\log2\}/\pi,
 \notag\\
 G_\phi&=\frac{LV_A^\phi+V_B^\phi}{2m}+\frac{6LM_A^\phi}{m}
 +\frac4m\left\{7+\frac P\pi(5+2\log(2m))\right\}
 +\frac{4d_+\log\pi}{\pi m}+\frac{2C_B^\phi}{m},\notag\\
 C_{6,\phi}&=\frac{21+21a+2L+\{(1+48p)H+24q\}/\pi}{n}
                         +\frac{317+5/12}{n^2},\notag\\
 \overline D_\phi&=4(c_e+c_d)/\pi+ab_0\log\pi,\quad
 \overline Z_\phi=\overline D_\phi+ab_0/n,\notag\\
 \mathcal M_\phi&=\frac{ab_0(2L+2r_0)}n
 +\frac{\overline D_\phi+ab_0(L+1/n)}{4n^2-1}
 +\frac{ab_0}{(2n-1)^2}\notag\\
 &\qquad+\frac{\overline Z_\phi+ab_0(\log(n+1/2)+2)}{4(n-1/2)^2},\notag\\
 W_n&=(3\log m+\log2+2\log\pi+2)/(\pi m),\notag\\
 C_{E,\phi}&=D_\phi+G_\phi+C_{6,\phi}+\mathcal M_\phi+d_+W_n.
\end{align}
Set $\Gamma_\phi=0$ at $p=0$. For its uniform version define
\begin{equation}\label{uc:uniform-allowance}
 \begin{gathered}
 g_n=1+\log(2+8m/(3\pi)),\qquad M_A=1/\pi+4/\pi^2,\\
 G_g=\frac1m\left\{7M_A L+\frac{111}2+
 \frac{63+\log2+4\log\pi+7\log m+16\log(2m)+(3/2)g_n}\pi\right\},\\
 \overline C_E=D+G_g+(64+16L)/n+318/n^2+\mathcal M+W_n.
 \end{gathered}
\end{equation}

The phase-dependent list preserves factors such as $a,p,q$ because
their largest values are incompatible. The uniform list is simpler
and is used when only a global bound is needed. The expression
$p[1+\log(1+2\overline s/p)]$ has continuous value zero at $p=0$:
for bounded positive $A$, $p\log(1+A/p)\to0$. Its logarithm measures
an integrated parameter singularity near a moving equality root.

The proposition has three distinct outputs. The first is an absolute
approximation and a global norm bound. The second is a lower estimate
at the actual finite origin. The last compares $S_n$ with the exact
finite corner $\Lambda_c$. The lower origin estimate is used only
for $n\geq160$, which explains that restriction in the last result.
The two-radius form of the last bound keeps distinct constants on the
two domains, with tighter bounds on the smaller domain.

\begin{proposition}[Finite comparison bounds]\label{uc:finite-bounds}
On the closed phase triangle,
\begin{equation}\label{uc:global-error}
 |S_n-M_\infty|\leq C_{E,\phi}+K_DP,
 \qquad C_{E,\phi}<\overline C_E.
\end{equation}
The quantities $D_\phi,G_\phi,C_{6,\phi},\mathcal M_\phi,W_n$
and $C_{E,\phi}$ decrease at a fixed phase for real $n\geq32$.
On the closed signed phase diamond the same bounds
hold with separate absolute phases. Every complete signed board sum,
and hence every complete signed cardinal sum, is bounded by
\begin{equation}\label{uc:analytic-norm}
 M(n)=\frac{4L^2+(79/5)L}{\pi^2}+\frac94+\overline C_E+2K_D.
\end{equation}
For the actual finite origin, writing $E_n=S_n-M_\infty$,
\begin{equation}\label{uc:origin-lower}
 E_n(0)>-\bar e\qquad(n\geq160).
\end{equation}
Finally, for $n\geq160$,
\begin{equation}\label{uc:relative-comparison}
 S_n-\Lambda_c\leq
 \begin{cases}
 E_c-uU_N+C_4+k_{Nx}x+k_{Ny}y,&\rho\leq1/2,\\
 E_c-uU_W+C_4+k_{Wx}x+k_{Wy}y,&\rho\leq1.
 \end{cases}
\end{equation}
\end{proposition}

We prove these statements with one exact identity. The actual rows are
\[
 \begin{gathered}
 \delta_r=(1-r)\bmod2,\quad
 u_{rk}=(\delta_r\pi+x+2k\pi)/n,\quad v_r=(r\pi+y)/m,\\
 R_r=\sum_{k=0}^{n-1}|G_n(u_{rk},v_r)/2|,
 \qquad S_n=(nm)^{-1}\sum_{r=0}^{2m-1}R_r.
 \end{gathered}
\]
\paragraph{The two meshes and their linked collision values.}
For each row index $r$, the bit $\delta_r$ is zero or one according
to its parity. Thus the row samples are spaced by $2\pi/n$, with
the two parity classes offset by $\pi/n$. There are exactly $n$
samples per row and $2m$ rows. The normalization $(nm)^{-1}$ is the
normalization of the signed board from the preceding construction;
it must be used with these multiplicities.

A regular sample is one at which the cotangent formula has no zero
denominator. A collision occurs when an argument of a cotangent is an
integer multiple of $\pi$. The original $G_n$ is a polynomial kernel
in trigonometric coordinates, so its value there is finite. The apparent
singularity comes from rewriting that polynomial as cotangents. A
linked limit means varying the actual phase variables together and
taking the limit of the entire polynomial identity. This fixes the
value uniquely. Taking separate limits of a vanishing coefficient and
a divergent cotangent without preserving their relation could give an
incorrect sample value.

At regular samples the exact polynomial identity is
\[
 (-1)^{\delta_r}G_n(u,v_r)/2
 =C(v_r)\cot((u-v_r)/2)+D(v_r)\cot((u+v_r)/2).
\]
At collisions we use the complete polynomial continuation.
Retain
$A_6=\{0,1,2m-2,2m-1,m-1,m\}$, with
$J_6=(nm)^{-1}\sum_{r\in A_6}R_r=J_{\rm end}+J_{\rm cen}$.

The first four indices are the rows nearest the periodic endpoint,
while $m-1,m$ are the two rows nearest the center. They are retained
as finite rows throughout the argument. The endpoint rows contain the
largest peaks and their removable values. The central rows can carry
linked collisions that are too large relative to their local model to
hide in an ordinary quadrature error. Removing these six rows from
the outer sampling leaves a regular range on which centered-cell
estimates work uniformly.

The endpoint reference $C_2$ will match these four endpoint rows and
the completed harmonic tail. It is a finite expression with a prescribed
continuation at its apparent poles. The two central rows have no such
endpoint reference; they enter $E_{\rm ret}$ directly through $J_6$.
This assignment keeps every actual row in exactly one part of the
error identity.
For $t=y/\pi$, $h=\pi/m$, set
\begin{align}\label{uc:mesh-series}
 h_j(t)&=-\log(j+t-1/2)+\sum_{k=j}^\infty
 \left\{\frac1{k+t}-\log\frac{k+t+1/2}{k+t-1/2}\right\},\notag\\
 g_j(t)&=-\tfrac12\log^2(j+t-1/2)\notag\\
 &\quad+\sum_{k=j}^\infty
 \left\{\frac{\log(k+t)}{k+t}
       -\frac{\log^2(k+t+1/2)-\log^2(k+t-1/2)}2\right\}.
\end{align}
The mesh series are built from the error of replacing a sample by
the integral over its centered unit cell. For $w=k+t$,
\[
 \int_{w-1/2}^{w+1/2}\frac{ds}{s}
   =\log\frac{w+1/2}{w-1/2},\qquad
 \int_{w-1/2}^{w+1/2}\frac{\log s}{s}ds
   =\frac{\log^2(w+1/2)-\log^2(w-1/2)}2.
\]
The braces in $h_j,g_j$ are precisely the sample minus this cell
integral. The linear Taylor term integrates to zero on a centered cell;
the second derivative bounds the remaining error. For these two
functions the defects are respectively $O(w^{-3})$ and
$O((1+\log w)w^{-3})$, giving absolute convergence. The initial
logarithm terms record the lower edge of the cell union, so subtracting
two tail series recovers an exact finite sum.

We use $\operatorname{sinc}t=\sin t/t$ for $t\ne0$ and
$\operatorname{sinc}0=1$, the continuous value supplied by
$\sin t/t\to1$.

Both series converge absolutely. The four-row endpoint reference is
\[
 \begin{gathered}
 C_2=\operatorname{sinc}y+\ell_1(x,y)+\ell_1(x,-y)+\ell_2(x,-y),
 \\ \ell_r(x,b)=\frac{\sin b}{w}
   +8a\sum_{j=-r+1,-r+3,\ldots,r-1}\frac1{w^2-(j\pi+x)^2},
 \quad w=r\pi+b.
 \end{gathered}
\]
Its linked sinc continuation is specified below. Set
$\mathcal H_j=h_j(t)+h_j(1-t)$, $\mathcal Q_j=g_j(t)+g_j(1-t)$,
$E_2=C_2+D_0\mathcal H_2+B_0\mathcal Q_2$,
$\Delta_2=E_\infty-E_2$, $\epsilon_n=\log(m/n)$,
$Z_n=D_0-B_0\epsilon_n$, and
\begin{align*}
 F_s(L,z)&=\mathcal H(L,v_s(z))-
 \frac{U_0+(4a/\pi)(L+\log z)}z-\frac{d_+}\pi\log(\pi-z),\\
 Q_\tau[F]&=\frac1m\sum_{k=2}^{m-2}F(h(k+\tau))
                      -\frac1\pi\int_0^\pi F(z)\,dz,\\
  E_{\rm row}&=S_n-J_6-\frac1m\sum_{r\notin A_6}\mathcal H(L,v_r)-\Delta_2,\\
  E_{\rm out}&=Q_t[F_L]+Q_{1-t}[F_R],\quad E_{\rm ret}=J_6-C_2,\\
 R_{\rm mesh}&=-B_0\epsilon_n\mathcal H_2-Z_n\mathcal H_n
               -B_0\mathcal Q_n-B_0L^2-2D_0L,\\
 R_{\log}&=\frac1m\sum_{k=2}^{m-2}
 \{\log(\pi-h(k+t))+\log(\pi-h(k+1-t))\}-2(\log\pi-1).
\end{align*}
\paragraph{Reading the five-term decomposition.}
The function $F_s$ is the outer profile after subtracting its endpoint
reciprocal terms and its central logarithm. The operator $Q_\tau$
compares its actual finite sample average with its integral. It includes
the omitted intervals at both ends; it is not merely the error on the
union of sampled cells. The terms in the decomposition have the
following concrete roles.

The row term compares actual nonretained rows with their continuous
row profiles, with the completed endpoint-tail correction $\Delta_2$
removed. The outer term compares the regularized profiles with their
integrals. The retained term compares the six exact rows with the
four-row reference. The mesh term is the remaining explicitly summed
reciprocal and logarithmic-reciprocal contribution. Finally $R_{\log}$
keeps the central logarithm as its own finite sum. This last term has a
phase factor $d_+$ and will be charged once throughout the proof.

These definitions are identities before any bounds are applied. In
particular the finite reference $C_2$ appears once with each sign:
inside the retained-row discrepancy and inside the completed correction.
This is how the exact finite endpoint values are connected to the
constant of the model.

All subtracted integrals concern their complete expressions. Direct
subtraction of the series at indices $2$ and $n=m-1$ gives
\begin{equation}\label{uc:telescope}
 E_n=E_{\rm row}+E_{\rm out}+E_{\rm ret}+R_{\rm mesh}
                       +(d_+/\pi)R_{\log}.
\end{equation}
Indeed the singular sample is exactly
$Z_n/(k+t)+B_0\log(k+t)/(k+t)$; its cell primitives telescope.
The central logarithm is kept as its finite sum. The completed
finite part then cancels $\Delta_2$ and the two lower mesh constants.
This derivation preserves $m=n+1$ and all finite origin values.

\paragraph{Collecting the finite sums and the completed integral.}
Here is the full algebra that produces $R_{\rm mesh}$. Let
\[
 \mathcal I_F=\frac1\pi\int_0^\pi(F_L+F_R)dz,
 \qquad \mathcal L_c=2(\log\pi-1).
\]
The latter is the integral average of the two central logarithms.
By the definitions of $F_s$ and the completed finite part,
\[
 \mathcal I_F=(b_{\rm full}-2D_0)L+C_\infty-E_\infty
                         -\frac{d_+}{\pi}\mathcal L_c.
\]
For its $L$ coefficient, integrate
$A(z)+A(2\pi-z)-8a/(\pi z)$ and use the coefficient calculation
above: after division by $\pi$ its integral is
$b_{\rm full}-2U_0/\pi-8a\log\pi/\pi^2=b_{\rm full}-2D_0$.
For its constant coefficient, integrate the two subtracted singular
terms up to $\pi$. Their endpoint primitives cancel exactly the
counterterms defining $C_\infty$, leaving $C_\infty-E_\infty$.
Finally the subtracted central logarithm contributes
$-(d_+/\pi)\mathcal L_c$.

Let $B_{\rm ord}$ denote the ordinary-row profile sum in the
definition of $E_{\rm row}$. The finite harmonic recurrences give
\[
 \begin{aligned}
 B_{\rm ord}
 &=E_{\rm out}+\mathcal I_F
       +Z_n(\mathcal H_2-\mathcal H_n)
       +B_0(\mathcal Q_2-\mathcal Q_n)
       +\frac{d_+}{\pi}(R_{\log}+\mathcal L_c),\\
 S_n&=E_{\rm row}+J_6+B_{\rm ord}+E_\infty-E_2.
 \end{aligned}
\]
Insert $E_2=C_2+D_0\mathcal H_2+B_0\mathcal Q_2$ in the second
line. The $C_2$ term combines with $J_6$ into $E_{\rm ret}$;
the two lower logarithmic series cancel; the lower harmonic coefficient
is $(Z_n-D_0)\mathcal H_2=-B_0\epsilon_n\mathcal H_2$.
Moreover
\[
 E_\infty+\mathcal I_F+\frac{d_+}{\pi}\mathcal L_c
 =M_\infty-B_0L^2-2D_0L.
\]
Subtracting $M_\infty$ therefore leaves, besides the row, outer,
retained, and central-log terms, exactly
\[
 -B_0\epsilon_n\mathcal H_2-Z_n\mathcal H_n-B_0\mathcal Q_n
                                      -B_0L^2-2D_0L.
\]
This is the displayed $R_{\rm mesh}$, including both its finite upper
index and its constant normalization.
The same identity holds after applying
$\delta_aX=X(x,y)-aX(0,0)$.

\paragraph{Algebra of the telescope.}
The word telescope refers to cancellation of neighboring cell
primitives. With $z=h(k+t)=\pi(k+t)/m$, the singular sample, including
the factor $1/m$, is exactly
\[
 \frac1m\frac{U_0+(4a/\pi)(L+\log z)}z
 =\frac{D_0-B_0\log(m/n)}{k+t}
       +B_0\frac{\log(k+t)}{k+t}.
\]
This is the stated expression with $Z_n=D_0-B_0\epsilon_n$.
The identity uses $m=n+1$ in $\epsilon_n=\log(m/n)$.
For either offset the series definitions give
\[
 h_2(t)-h_n(t)=\sum_{k=2}^{n-1}\frac1{k+t},\qquad
 g_2(t)-g_n(t)=\sum_{k=2}^{n-1}\frac{\log(k+t)}{k+t}.
\]
Indeed all intermediate logarithmic cell boundaries cancel exactly.
Pairing the two offsets produces $\mathcal H_2-\mathcal H_n$ and
$\mathcal Q_2-\mathcal Q_n$. The lower terms combine with $E_2$;
the upper terms and the explicit integral counterterms produce
$R_{\rm mesh}$. The identity thus keeps the small distinction between
$\log m$ and $\log n$ as a known term.

The operator $\delta_a$ is linear for a fixed phase: it subtracts
$a$ times the origin value. Since the leading model coefficient is
also multiplied by $a$, this weighted subtraction cancels large common
terms before estimates are taken. The final conversion from this
weighted difference to an ordinary phase difference will use the
one-sided finite origin bound.

\paragraph{Pole norms and the matched endpoint reference.}
For $B_n(t)=\sum_{k=0}^{n-1}|\cot(\pi(k+t)/n)|$, $0<t<1$,
\begin{equation}\label{uc:pole-error}
 \left|\frac{B_n(t)}n-
 \frac{2L+\mathcal C(t)-2\log\pi}{\pi}\right|\leq\frac5{n^2}.
\end{equation}
\paragraph{Why a uniform pole estimate is possible.}
As the fractional offset $t$ approaches zero, the unweighted pole sum
and $\mathcal C(t)$ both diverge. Their difference nevertheless has a
uniform bound because the same reciprocal singularities have been
subtracted from each. This is a recurring principle of the comparison:
identify the singular expression explicitly, cancel it exactly, and
estimate its regular remainder.

The weighted version has a still better interpretation at an integer
offset. Multiplication by $\sin\pi t$ replaces the divergent reciprocal
sample by its removable limit. The limit one is the cardinal peak
value in this normalization. Thus the bound continues to the collision
values when the corresponding polynomial row is used.

Here is a proof uniform in the offset. Subtract $1/w+1/(\pi-w)$
from $|\cot w|$ on $(0,\pi)$. The residual $b$ has integral
$-2\log\pi$, is symmetric, and $b-|w-\pi/2|$ has Lipschitz
first derivative with second derivative bounded by two.
This follows directly by twice differentiating the cotangent partial
fraction: $|(\cot w-1/w)''|<1$ on $[0,\pi/2]$.
For a smooth function $b$ with $|b''|\leq M$, Taylor's formula
with integral remainder bounds the deviation from its linear part on
a cell by $M$ times a quadratic distance. Pairing the offsets cancels
the linear contributions because $b$ is symmetric about $\pi/2$.
The residual has one corner at $\pi/2$, represented by
$|w-\pi/2|$. For even $n$ its cell placement makes the paired error
zero; for odd $n$ its single crossing cell gives the extra explicit
allowance. Thus the error is $O(n^{-2})$ even though the original
cotangent has poles. The reciprocal part is evaluated separately
through the defining harmonic series, with an independently bounded
tail.

On each cell, pairing offsets $t,1-t$ cancels the linear term;
its error is at most $h^2/3$, $h=\pi/n$. The cusp contributes zero
for even $n$ and at most $h/(4n)$ for odd $n$. Thus the residual
average error is at most $(\pi^2/3+\pi/4)/n^2$.
The remaining harmonic-tail error is at most $1/(2\pi n^2)$:
its $k$th paired summand is bounded by $1/(3k^3)$ using the
midpoint and trapezoid errors for $(k+t)^{-1}$, and integral comparison
sums the tail. Their total is less than $5/n^2$.
\paragraph{Cell kernels and the numerical pole allowance.}
The elementary cotangent residual $c(w)=\cot w-1/w$ has
\[
 c''(w)=-4w\sum_{j\geq1}
       \frac{3\pi^2j^2+w^2}{(\pi^2j^2-w^2)^3},\qquad 0<w\leq\pi/2.
\]
This follows by differentiating the convergent partial fraction of
cotangent after its $1/w$ pole has been removed. Each positive term
in $-c''$ increases with $w$, so its maximum is attained at $\pi/2$.
Direct differentiation at that endpoint gives $-c''(\pi/2)=16/\pi^3<1$.
For $w<\pi/2$, the residual used in the pole proof is
$b(w)=c(w)-1/(\pi-w)$; hence $|b''|\leq1+16/\pi^3<2$ there.
At $\pi/2$ its slopes are $-1$ and $1$. Subtracting
$|w-\pi/2|$ removes that jump.

For a cell of width $h$ and midpoint $c_*$, Taylor's integral
remainder gives a sample error at offset distance at most $h/2$
bounded by $Mh^2/8$, and a cell-average error bounded by $Mh^2/24$
when $|F''|\leq M$. Pairing offsets $t,1-t$ removes the linear
part. Taking $M=2$ gives $h^2/3$. For odd $n$ the one cusp cell
has paired error $h(|t-1/2|-1/4)$; after averaging over $n$ cells
this costs at most $\pi/(4n^2)$. For even $n$ each cell sees a linear
cusp function and its paired error is zero.

For completeness, write $H_n^*(t)=\sum_{k=0}^{n-1}(k+t)^{-1}$.
Its paired tail identity is
\[
 H_n^*(t)+H_n^*(1-t)-2\log n-\mathcal C(t)
 =\sum_{k=n}^\infty\left\{2\log(1+1/k)-\frac1{k+t}-\frac1{k+1-t}\right\}.
\]
Convexity places the paired reciprocal values between the midpoint
pair and the endpoint pair on $[k,k+1]$. The trapezoid error is at
most $\sup|F''|/12$ and the midpoint error at most
$\sup|F''|/24$. For $F(w)=1/w$ one has
$\sup|F''|=2/k^3$, so the absolute value of the braces is at most
$1/(3k^3)$. Summing bounds it by
$1/(6n^2)+1/(3n^3)\leq1/(2n^2)$, and the pole normalization
adds the factor $1/\pi$. Thus the total coefficient is
\[
 \frac{\pi^2}{3}+\frac\pi4+\frac1{2\pi}<5,
\]
which is the allowance used in the pole estimate.

We will also use the first-order cell identity for a continuous
function of bounded variation. On a cell $[a,a+h]$ with center $c_*$,
\[
 F(c_*)-\frac1h\int_a^{a+h}F(z)dz
 =\int_a^{a+h}K(z)\,dF(z),\quad
 K(z)=\begin{cases}(z-a)/h,&z<c_*,\\(z-a-h)/h,&z>c_*.
 \end{cases}
\]
The fundamental theorem of calculus proves it first for a smooth
function; integration against its finite change measure gives the same
formula for bounded variation. Since $|K|\leq1/2$, disjoint cells
have total sample-minus-integral defect at most half the variation.
For a piecewise smooth function with jumps away from $c_*$, apply the
identity on each continuity interval and include each jump as a point
mass in $dF$. The same half-variation bound for an arbitrary assigned
sample value follows directly by splitting the cell into its two halves:
\[
 \begin{aligned}
 \left|F(c_*)-\frac1h\int_a^{a+h}F(z)dz\right|
 &\leq\tfrac12\operatorname{Var}_{[a,c_*]}F
            +\tfrac12\operatorname{Var}_{[c_*,a+h]}F\\
 &=\tfrac12\operatorname{Var}_{[a,a+h]}F.
 \end{aligned}
\]
On each half, $|F(c_*)-F(z)|$ is bounded by that half's variation;
integrating over its length $h/2$ proves the inequality. This form
also covers the finite values assigned at removable pole samples.
This concrete kernel explains all first-order BV quadrature factors
below.

Multiplication by $\sin\pi t$ gives a continuous periodic bound,
with value one and zero error at integer offsets. Since
\[
 C(v)=-\frac{c_e\sin\pi\eta_-(v)}{2\sin(v/2)},\qquad
 D(v)= \frac{c_d\sin\pi\eta_+(v)}{2\sin(v/2)},
\]
the two weighted row errors, summed and normalized, are at most
$10/n^2+20aH_n/(\pi n^2)$. Indeed
$|C|+|D|\leq1+a|\cot(v/2)|$ and the two endpoint distance
indices give $\sum|\cot(v_r/2)|\leq4mH_n/\pi$.
At a collision this calculation supplies the profile; the complete
polynomial row supplies the actual value.

\paragraph{Matching the endpoint tail.}
A half-line defect is the difference between samples and an integral
along one ray starting at an endpoint. Here it is matched cell by cell:
the sample and the integral use the same finite cutoff before a limit
is taken. This matching matters because reciprocal samples and their
integrals each have logarithmically divergent tails. Their difference
has a finite limit only after the common tail is cancelled.

For a function of bounded variation, the sum of absolute changes over
any disjoint intervals is at most its total variation. Comparing one
sample with the average over its own cell therefore bounds the
sample-minus-integral defect by that variation. Subtracting the
reciprocal channel gives the residual below, whose variation is
$4/\pi$. The reciprocal part can then be summed using the defining
series of $\psi$. This explains both the uniform constant and the
appearance of the digamma function in the matched reference.

For the cancellation reference $(-w,w)$ the half-line defect is
particularly simple. If $h_*=\pi/n$, $d=v/h_*$ and $0<v<\pi/2$,
\[
 \left|\sum_{v+h_*(k+t)<\pi/2}\cot(v+h_*(k+t))
       +\frac{\log\sin v}{h_*}
       -\frac{\log d-\psi(d+t)}{h_*}\right|\leq4/\pi.
\]
To prove it, extend $\cot(v+w)$ by zero beyond $\pi/2$ and subtract
$(v+w)^{-1}$. The residual decreases to $-2/\pi$ and increases to
zero, so its variation is at most $4/\pi$. Compare each sample with
its own cell average, first in a finite matched sum. The reciprocal
part tends to $h_*^{-1}(\log d-\psi(d+t))$ by the defining series.
Applying this matched-sum limit to the two
nonzero reference intervals gives a defect $T_0$ and
\begin{equation}\label{uc:halfline}
 \Theta=\frac{2w}{h_*}
 \{\psi(d+t_-)+\psi(d+1-t_+)-2\log d\},
 \qquad |T_0-\Theta|\leq16|w|/\pi,
\end{equation}
where $t_-,t_+$ are the actual fractional pole offsets.
In row $r$ take $w_r=a\cot(v_r/2)/2$ and
$X_{rk}=(u_{rk}-v_r)/2$. The notation $T_{0,r}$ means
$\sum_k g(-w_r,w_r;X_{rk})-(n/\pi)\int_0^\pi
g(-w_r,w_r;X)\,dX$, with separation $v_r$;
reflection also allows negative $w_r$.

We now identify the endpoint reference as $\Delta_2$. For $j\geq2$, $|s|+x/\pi\leq1$,
put $q_\pm=(1+s\pm x/\pi)/2$, $\nu=j+s$ and
\[
 B_j(s)=\psi(j+q_+)+\psi(j+q_-)-2\log\nu,\qquad
 \Xi_j(s)=\frac{2a}{\pi^2\nu}B_j(s).
\]
The temporary row index $j$ describes a farther endpoint row. Its
reference amplitude decays like $1/(j+s)$, while its remaining harmonic
bracket $B_j$ is also $O((j+s)^{-1})$. Their product $\Xi_j$ therefore
has a summable $j^{-2}$ tail. This summability is what allows a finite
endpoint reference to be matched with a completed infinite one. The
shifts $q_\pm$ still retain both phase directions.

The elementary harmonic estimate
$0\leq\log w-(2w)^{-1}-\psi(w)\leq(8w^2)^{-1}$ follows by
unit-cell trapezoid integration of $1/w$. It implies
$|B_j(s)|\leq2/\nu$, $|\Xi_j(s)|\leq4a/(\pi^2\nu^2)$.
The harmonic estimate itself is a sum of elementary trapezoid
errors. For $F(z)=1/z$, the error on $[w+k,w+k+1]$ is
\[
 \frac{F(w+k)+F(w+k+1)}2-\int_{w+k}^{w+k+1}F(z)dz
 =\int_0^1\frac{r(1-r)}2 F''(w+k+r)dr\geq0.
\]
Summing and taking the defining limit for $\psi$ gives
$\log w-(2w)^{-1}-\psi(w)$. Since $r(1-r)/2\leq1/8$,
the sum is at most $(1/8)\int_w^\infty F''=1/(8w^2)$.
For the present arguments $j+q_\pm\in[\nu,\nu+1]$;
monotonicity of $\psi$ and this estimate then give
$|\psi(j+q_\pm)-\log\nu|\leq1/\nu$, proving the asserted
$B_j$ and $\Xi_j$ bounds.

Partial fractions in the finite reference give
\[
 \ell_j(x,b)=\frac{D_0+(2a/\pi^2)
                \{\psi(j+q_-)+\psi(j+q_+)\}}{j+s},\qquad s=b/\pi.
\]
The constant follows from
$\psi(1-t)-\psi(t)=\pi\cot\pi t$, itself the cotangent partial
fraction.

To expand that partial-fraction step, for the $j$ finite reference
indices put $z_k=(-j+1+2k)\pi+x$, $0\leq k<j$. Then
\[
 \begin{gathered}
 w-z_k=2\pi(j-1-k+q_-),\qquad w+z_k=2\pi(k+q_+),\\
 \frac1{w^2-z_k^2}=\frac1{2w}\left(\frac1{w-z_k}+\frac1{w+z_k}\right).
 \end{gathered}
\]
Summation gives
\[
 \ell_j=\frac1{\pi\nu}\left[\sin b+\frac{2a}{\pi}
   \{\psi(j+q_-)+\psi(j+q_+)-\psi(q_-)-\psi(q_+)\}\right].
\]
The numerator constant is $\pi D_0$. Indeed
$D_0=a\{\mathcal C(t_d)+\mathcal C(t_e)\}/\pi^2$ follows by
inserting $f$ into $U_0$. For $b=y$, $q_-=1-t_d$, $q_+=t_e$,
so reflection reduces the difference of the proposed constants to
\[
 \frac{\sin y}{\pi}+\frac a\pi
              \{\cot(\pi t_e)-\cot(\pi t_d)\}=0.
\]
The cotangent difference is $-\sin y/a$ by the sine addition formula.
For $b=-y$ the two complementary arguments exchange roles and the
same cancellation holds. Initially these calculations use nonzero
weights; their completed products extend continuously to the boundary.

Increasing a temporary cutoff from $j$ to $j+1$ adds
$\ell_j(x,y)+\ell_{j+1}(x,-y)$ to its finite reference. At the same
time $h_j(t)-h_{j+1}(t)=1/(j+t)$ and
$g_j(t)-g_{j+1}(t)=\log(j+t)/(j+t)$. Inserting the just-derived
formula for $\ell_j$ cancels $D_0$ and the logarithmic parts, leaving
exactly $\Xi_j(t)+\Xi_{j+1}(-t)$. This proves the cutoff recurrence
before its limit is taken. The $h_j,g_j$ recurrences therefore show that increasing
the temporary cutoff $j$ adds $\Xi_j(t)+\Xi_{j+1}(-t)$.
To identify its limit, sum
\[
 \frac{H_j(q_-)+H_j(q_+)}{j+s}
 =\sum_{k+l=j-1}\frac1{(k+q_-)(l+q_+)},\qquad
 H_j(q)=\psi(j+q)-\psi(q).
\]
The finite convolution identity can be verified term by term. When
$k+l=j-1$, the denominators satisfy
$(k+q_-)+(l+q_+)=j+s$, because $q_-+q_+=1+s$.
Consequently
\[
 \frac{1/(k+q_-)+1/(l+q_+)}{j+s}
 =\frac1{(k+q_-)(l+q_+)}.
\]
Summing gives the displayed identity. Summing again over $j$ fills a
triangle of lattice indices. The product of two harmonic sums fills a
square. Their missing triangular region yields the integral described
next. This geometric distinction supplies the $\pi^2/6$ constant;
it cannot be dropped in identifying $E_\infty$.

The triangular convolution differs from the square harmonic product by
$\pi^2/6+o(1)$: after scaling, the missing triangle has integral
$\int_0^1-\log(1-w)\,dw/w=\pi^2/6$. Restricting first to both
indices at least $\eta j$ justifies ordinary Riemann sums; the omitted
sides are $O(\eta)$, since a side index $k$ has only $k+2$ partners.
The thin-side estimate in that Riemann limit can be made explicit.
Let $p_*,q_*>0$ be fixed small shifts and compare
\[
 T_R=\sum_{k+l\leq R-1}\frac1{(k+p_*)(l+q_*)}
 \quad\text{with}\quad
 H_R(p_*)H_R(q_*).
\]
In the missing upper triangle with $k\leq\eta R$, one has
$l\geq R-k-O(1)$, at most $k+2$ possible partners, and
$l+q_*\geq R(1-\eta)-O(1)$. Its contribution is bounded by
\[
 \frac{C}{R}\sum_{k\leq\eta R}\frac{k+2}{k+p_*}
 \leq C\eta+C_{p_*}\frac{1+\log R}{R}.
\]
The other thin side is identical. Away from both sides ordinary
Riemann sums apply. Letting first $R\to\infty$ and then
$\eta\downarrow0$ gives
\[
 H_R(p_*)H_R(q_*)-T_R\longrightarrow
 \int_0^1\frac{-\log(1-r)}rdr=\sum_{k\geq1}\frac1{k^2}=\frac{\pi^2}{6}.
\]
The series follows by integrating the nonnegative power series for
$-\log(1-r)$ term by term. Also
$H_R(q)=\psi(R+q)-\psi(q)=\log R-\psi(q)+O(1/R)$,
by the displayed harmonic estimate. The product error is consequently
$O(\log R/R)$ and tends to zero. Multiplying these two expressions and subtracting the
triangle constant supplies the product terms in the limit below.
The signed sine contribution is explicitly
\[
 \frac{\sin y}{y}+\frac{\sin y}{\pi}
       \sum_{r\geq1}\left(\frac1{r+y/\pi}-\frac1{r-y/\pi}\right)
 =\cos y,
\]
again by the cotangent partial fraction. This displays each constant
left after the logarithmic cutoff terms cancel.

Also $h_j=-\log j+o(1)$, $g_j=-\tfrac12\log^2j+o(1)$ by their
cell tails, and the signed sine sum including $\ell_0$ tends to
$\cos y$. If $t_d=1/2+(x-y)/(2\pi)$,
$t_e=1/2+(x+y)/(2\pi)$, the resulting limit is
\[
 \cos y+\frac{2a}{\pi^2}
 \{\psi(1-t_d)\psi(t_e)+\psi(1-t_e)\psi(t_d)-\pi^2/3\}.
\]
Reflection replaces the two products in braces by
$\{\mathcal C(t_d)\mathcal C(t_e)-\pi^2\cot(\pi t_d)
\cot(\pi t_e)\}/2$ before the last subtraction.
Since $a\cot(\pi t_d)\cot(\pi t_e)=(\cos y-\cos x)/2$,
this is exactly $E_\infty$.
The uniform $\Xi_j$ tail includes the closed boundary, since its
weighted endpoint products tend to $\pi$. Thus
\begin{equation}\label{uc:reference-match}
 \begin{gathered}
 \Delta_2=\sum_{j=2}^\infty\{\Xi_j(t)+\Xi_{j+1}(-t)\},\\
 \left|\frac1{nm}\sum_{r\notin A_6}T_{0,r}-\Delta_2\right|
 \leq\frac{[52a/\pi^2+16P/(3\pi)]H_n+7a+2P/\pi}{n}.
 \end{gathered}
\end{equation}
For completeness, the finite charges in the second assertion are as
follows. On either endpoint side $v=\pi\nu/m<\pi/2$, the actual
arguments in \eqref{uc:halfline} are
$j+q_\pm-\nu/(2m)-k_\pm$, with $k_\pm\in\{0,1\}$.
Before a wrap, $\psi'\leq2/j$ bounds the displacement and
$|\log(n/m)|\leq1/n$, giving $5/n$ in the braces.
The normalized prefactor is
$2a\tau(v/2)/(\pi^2\nu)$, $\tau(w)=w\cot w$;
$0\leq1-\tau(w)\leq w^2$ costs $a/m^2$ per row.
A wrap requires $\rho\geq\pi-v$, hence
$a=(P/2)\cot(\rho/2)\leq Pv/2$; its recurrence cost is at most
$8P/(3\pi m(j-1))$ per row. The two endpoint sums consequently
cost $[20a/\pi^2+16P/(3\pi)]H_n/n+2a/n$.
Their missing infinite tails cost $2a/n$, omitted collisions cost
$2P/(\pi n)$, and the residual in \eqref{uc:halfline} costs
$32aH_n/(\pi^2n)$. Centrally the reference variation is
$8|w\cot v|$, so centered-cell quadrature gives
$|T_0|\leq4|w|/|\sin v|=a/\sin^2(v/2)\leq2a$;
its normalized sum costs $3a/n$. These are all the terms in
\eqref{uc:reference-match}.

Their normalized sum can be checked in a short ledger:
\[
\begin{array}{c|c}
 \text{source}&\text{upper charge}\\ \hline
 \text{digamma displacement}&20aH_n/(\pi^2n)\\
 \tau\text{ prefactor}&2a/n\\
 \text{offset wraps}&16PH_n/(3\pi n)\\
 \text{missing infinite tails}&2a/n\\
 \text{omitted collision references}&2P/(\pi n)\\
 \text{half-line residual}&32aH_n/(\pi^2n)\\
 \text{central reference rows}&3a/n.
\end{array}
\]
For the displacement, $\delta=\nu/(2m)$ moves each argument by
at most $\delta$ and $\psi'\leq2/j$, while the common logarithm
moves by $2\log(m/n)\leq2/n$. This gives $5/n$ in the braces,
then $10a/(\pi^2n\nu)$ per row. The combined endpoint reciprocal sum is
at most $2H_n$. The prefactor difference uses
$1-\tau(z/2)\leq z^2/4$ and $|B_j|\leq2/\nu$, giving
$a/m^2$ per row.
For a wrap, use its actual digamma argument
$z_{\rm arg}=j+q_\pm-\nu/(2m)\geq j-1/4$. Then
$z_{\rm arg}-1\geq j-5/4\geq3(j-1)/4$, and its prefactor is
at most $P/(\pi m)$; at most two wraps give the listed harmonic
charge. The reference tail starts at $\nu\geq m/2$ and
$\Xi_j=O(a\nu^{-2})$, producing $2a/n$. Each omitted collision
reference uses the same bound together with $a\leq Pv/2$.
The half-line residual is at most $16|w|/\pi$ per row and
$w\leq am/(\pi\nu)$, producing its $32/\pi^2$ harmonic factor.
The central angular interval has length $\pi$ and row spacing
$\pi/m$, so it contains at most $m+1$ rows. Each reference error
is at most $2a$, giving $2a(m+1)/(nm)<3a/n$. Adding the ledger gives $20+32=52$ and
$2+2+3=7$, as required.

This accounting keeps the finite indexing effects visible. A wrap
means that an offset has crossed the end of its unit interval; applying
$\psi(z+1)=\psi(z)+1/z$ restores the usual representative and produces
a reciprocal correction. Such a crossing forces a phase near the
triangle edge, where the factor $a$ is small relative to $P$ and the
row distance. The displayed inequality $a\leq Pv/2$ is what makes
the correction summable. The missing infinite tails, omitted collision
values, half-line residuals, and central rows have different sources,
so each has a separate allowance before they are added. Their sum
matches the coefficient $52a/\pi^2$ and the $P$ terms in the final
row bound.

\paragraph{Controlling a perturbed cancellation by variation.}
The total variation of a function is the supremum of sums of absolute
successive changes over all finite partitions. For a piecewise monotone
function it is the sum of the rises and falls, including jumps.
A shifted sampling grid can be compared with centered cells by
integrating a bounded sawtooth weight against these changes. Its
magnitude is at most $1/2$, so the defect is at most half the variation
in the normalization used here. This is the bounded-variation
quadrature principle used repeatedly below.

The perturbation estimate must use the difference of two complete
cancellations. A separate estimate for each large reference channel
would lose the small coefficient $P$. Preserving the common reference
first lets the variation depend only on the coefficient displacement
$\varepsilon=P/4$. The finite number of equality roots and ordinary
critical points bounds the number of monotone branches.

The ordinary perturbation has the sharp coefficient needed here.

In the general variation bound let $M=\max(|C|,|D|)$. For equal
coefficient signs, cancellation occurs where the cotangents have
opposite signs. On the two supports the smaller absolute cotangent
has maximal heights $\tan(e/2)$ and $\cot(e/2)$, attained where
the two magnitudes agree. Multiplication by the coefficients can only
reduce the bound $M$ times those heights. Each nonnegative minimum
increases and then decreases, and multiplying by $-2$ gives variation
at most four times its height. Their sum is
\[
 4M\{\tan(e/2)+\cot(e/2)\}=\frac{8M}{\sin e}.
\]
For opposite coefficient signs the two same-sign cotangent supports
have monotone minima. Their endpoint heights are at most
$|D|\cot e$ and $|C|\cot e$. Including the jump to assigned value
zero gives total variation at most
$4(|C|+|D|)\cot e\leq8M/\sin e$. This proves the general bound
used before the smaller perturbation is considered.
For $g(C,D;t)$ the integrand of \eqref{uc:cancellation}, assigning
value zero at the poles,
\[
 \operatorname{Var}_{\rm per}g(C,D)\leq
       8\max(|C|,|D|)/\sin e.
\]
For equal coefficient signs, the two supports have unimodal minima
with heights at most $M\tan(e/2)$ and $M\cot(e/2)$;
$g=-2\min$ has variation at most four times these heights.
For opposite signs the two minima are monotone, with endpoint heights
$|D|\cot e,C\cot e$, giving the same bound including the pole jumps.
For $f=g(-w+c,w+d)-g(-w,w)$, $|c|,|d|\leq\varepsilon$, this implies
\begin{equation}\label{uc:ordinary-variation}
 \operatorname{Var}f\leq24\varepsilon/\sin e,
 \qquad |\hbox{shifted quadrature defect of }f|
                  \leq12\varepsilon/\sin e.
\end{equation}
If $|w|\leq\varepsilon$, apply the preceding bound twice.
Otherwise signs are preserved, $|f|\leq2\varepsilon\cot e$ on
its two supports. The equality numerator is a constant plus a sine,
with at most two zeros, and on each nonagreeing branch the derivative
has at most one zero because $\sin(t+e)/\sin t$ is monotone.
There are at most four interior cuts in total. Counting their changes
and the two pole jumps gives $12(2\varepsilon\cot e)$.
The quadrature kernel on centered cells has modulus at most $1/2$,
proving the second assertion.

The variation count can be checked without assuming the roots are
separated. Write $A=|w|>\varepsilon$, and let the perturbed positive
coefficient magnitudes be $b_1,b_2\in[A-\varepsilon,A+\varepsilon]$.
With $X=\cot t$, $Y=\cot(t+e)$, the two cancellations vanish off
\[
 I_1=(0,\pi/2-e),\qquad I_2=(\pi/2,\pi-e).
\]
On $I_1$, $X>Y>0$ and the reference selects $Y$; on $I_2$,
$Y<X<0$ and it selects $|X|$. Therefore
\[
 f=-2\{\min(b_1|X|,b_2|Y|)-A\min(|X|,|Y|)\},
 \qquad |f|\leq M_*:=2\varepsilon\cot e.
\]
Its branch equation has numerator
$[(b_1-b_2)\sin(2t+e)+(b_1+b_2)\sin e]/2$, with at most two
roots in one period. On an agreeing branch $f$ is a constant multiple
of the selected reference cotangent, hence monotone. On the other
branches the equations for a stationary point are
\[
\begin{array}{c|c|c}
 &f&\text{stationary equation}\\ \hline
 I_1&-2b_1X+2AY&[\sin(t+e)/\sin t]^2=A/b_1\\
 I_2&2b_2Y-2AX&[\sin(t+e)/\sin t]^2=b_2/A.
\end{array}
\]
The ratio has derivative $-\sin e/\sin^2t<0$, so each support has
at most one such stationary point. Altogether there are at most four
interior cuts. If $N_i$ of them lie on support $I_i$, its path runs
from zero to a finite endpoint $a_i$ with $|a_i|\leq M_*$. The
variation on that path is at most $|a_i|+2N_iM_*$; the pole jump
adds another $|a_i|$. The total is at most
$2(N_1+N_2+2)M_*\leq12M_*\leq24\varepsilon/\sin e$.
For $A\leq\varepsilon$, the separate general bounds give
$8(2A+\varepsilon)/\sin e\leq24\varepsilon/\sin e$.
At $e=\pi/2$ the sign-preserved supports disappear and $f=0$.
Thus all regimes and possible coincident cuts are covered. For physical rows take
$\varepsilon=P/4$ and sum $|\sin v_r|^{-1}\leq2mH_n$;
this gives exactly $6PH_n/n$. Together with the pole error and
\eqref{uc:reference-match}, it proves
$|E_{\rm row}|\leq D_\phi+K_DP$.

In this last combination, the reference matching contributes the
terms proportional to $a$, the ordinary perturbation contributes
$6PH_n/n$, and the weighted pole errors contribute the $n^{-2}$
terms. Replacing $H_n$ by $H=1+\log n$ uses integral comparison
$\sum_{j=2}^n1/j\leq\int_1^n dt/t$. The bound $K_DP$ therefore
has explicit linear phase size, while $D_\phi$ retains the symmetric
reference part. At zero $P=0$, so this part of the absolute argument
immediately gives the sharper origin row bound $D$.

\paragraph{Absolute charges in the same telescope.}
The remaining estimates in \eqref{uc:telescope} are
\begin{equation}\label{uc:absolute-components}
 |E_{\rm out}|\leq G_\phi,\qquad |E_{\rm ret}|\leq C_{6,\phi},
 \qquad |R_{\rm mesh}|\leq\mathcal M_\phi,
 \qquad |R_{\log}|/\pi\leq W_n.
\end{equation}
The absolute estimates are assembled without changing the exact
identity. Each term will be bounded in the normalization already used
by $S_n$. There are three different small parameters: the cell width
$h=\pi/m$, the coefficient displacement measured by $P$, and the
phase-independent reciprocal tail $1/n$. Keeping them separate
explains why some terms are $O((\log n)/n)$ and others are
$O(n^{-2})$. The final constants are deliberately explicit so that a
finite degree threshold can be checked later.

We supply their finite-row and quadrature details. The two centered
cell unions are
$[h(3/2+t),\pi-h(3/2-t)]$ and
$[h(5/2-t),\pi-h(1/2+t)]$; their initial holes total $4h$ and
central holes total $2h$. The subtracted $L$ coefficient has supremum
$M_A^\phi$ and combined variation $V_A^\phi$.
The centered cells in the two folded families do not cover all of
$[0,\pi]$. Their omitted initial intervals have total length $4h$,
and their omitted central intervals have total length $2h$. A quadrature
estimate on the cell union controls the variation term; the integrals
over these omitted intervals require separate endpoint envelopes.
For the $L$ coefficient, its bounded maximum and monotonic pieces
suffice. For the constant coefficient, root motion and the scalar cusp
also enter, producing the components in the next table.

For the constant term, the following components give its combined
variation $V_B^\phi$:
\[
\begin{array}{c|c}
\text{pole regular part}&23\\
\text{at most three monotone parameter pieces}&3\Gamma_\phi/\pi\\
\text{separation motion}&2p\ell/\pi\\
\text{logarithmic coefficient products}&6P\ell/\pi\\
\text{remaining regular terms}&(26a+20+2d_+\log2)/\pi.
\end{array}
\]
Here are derivations of these entries. For the pole numerator $N_s$,
$|N_s|\leq8$, $\operatorname{Var}N_s\leq12$ and
$\operatorname{Var}N_s'\leq47/(2\pi)$ follow from
\eqref{uc:scalar-bounds}, including its derivative jump. Writing
$P_s^0=(N_s-U_0)/z+N_s([2\sin(z/2)]^{-1}-z^{-1})$ bounds each
variation by $47/(2\pi)+20(1/2-1/\pi)$ and its supremum by six;
the sum of variations is $20+7/\pi<23$.
To see the three numerator bounds separately, each of the two scalar
channels has supremum at most four, so $|N_s|\leq2\cdot4=8$.
Each argument has derivative of magnitude $1/(2\pi)$, and
$|\Phi'|\leq12$, giving
\[
 \operatorname{Var}N_s\leq\pi\frac{2\cdot12}{2\pi}=12.
\]
Each argument traverses half a scalar period. Its ordinary curvature
contributes at most $35/2$ to the variation of $\Phi'$, and its
single possible derivative jump contributes at most six. Therefore
\[
 \operatorname{Var}N_s'\leq
    \frac{2(35/2+6)}{2\pi}=\frac{47}{2\pi}.
\]
The ordinary numerator and its derivative have bounded variation
because $f'$ does, including its one cusp jump per period. Translation
can move that jump through a cell but cannot enlarge its total size.
This is why variation estimates remain uniform in all phase shifts.
For a regularized reciprocal term $(N_s-U_0)/z$, the numerator vanishes
at zero. Its variation is controlled by integrating changes of $N_s'$
against bounded weights; the apparent extra reciprocal singularity
is removed by this vanishing.

For the regular numerator calculation, the useful exact identity is
\[
 P_s^0(z)=\int_0^1N_s'(tz)dt+N_s(z)d(z),
 \qquad d(z)=\frac1{2\sin(z/2)}-\frac1z.
\]
The first term has variation at most $\operatorname{Var}N_s'$
by integration of partition sums. The product rule for variation is
$\operatorname{Var}(AB)\leq\sup|A|\operatorname{Var}B+
\sup|B|\operatorname{Var}A$. Here $d$ increases from zero to
$1/2-1/\pi$, so the product term has variation at most
$(8+12)(1/2-1/\pi)$. This gives
\[
 \operatorname{Var}P_s^0\leq\frac{47}{2\pi}
                 +20(1/2-1/\pi)=10+\frac7{2\pi},
 \qquad
 \sup|P_s^0|\leq\frac{12}\pi+8(1/2-1/\pi)<6.
\]
Adding two sides yields $20+7/\pi<23$. The derivative jump is included
in $\operatorname{Var}N_s'$; integrating its translated positions
uses its full variation, just as for its ordinary derivative.

Integrating the $R_s$ bound in \eqref{uc:R-bounds} uses
\[
 \int_0^S\log\frac{v+s}{|v-s|}\,ds
 \leq2v\left(1+\log\frac{S+v}{v}\right),
\]
obtained by splitting at $s=v$. The three monotone pieces and the
separation derivative give the next two entries. On the initial half,
put $w=a\cot(z/2)$. The logarithmic coefficients after subtracting
$2w\log\sin z$ are $\max(q-p,-2w)$ and
$\max(-P,q-p-4w,-2w)$; their combined variation and supremum are
at most $4P$ and $2P$. The product rule gives $6P\ell$.
To evaluate the parameter integral, when $S\geq v>0$ split at $v$:
\[
 \begin{aligned}
 \int_0^S\log\frac{v+s}{|v-s|}ds
 &=(S+v)\log(S+v)-(S-v)\log(S-v)-2v\log v\\
 &=2v\log\frac{S+v}{v}
       +(S-v)\log\left(1+\frac{2v}{S-v}\right)\\
 &\leq2v\left(1+\log\frac{S+v}{v}\right).
 \end{aligned}
\]
The endpoint $S=v$ is its continuous value. At $v=0$ the whole
remainder vanishes. On the initial half, the left parameter curve is
monotone and the absolute-value right curve has at most two monotone
pieces; each costs at most this integral. Their separation derivatives
are at most $2v\cot z$, so two sides cost
$4v\int_\epsilon^{\pi/2}\cot z\,dz=2p\ell$.
The two logarithmic coefficients have summed suprema at most $2P$
and summed variations at most $4P$. The product inequality with
$\sup|\log\sin z|=\operatorname{Var}(\log\sin z)=\ell$
therefore gives $(2P+4P)\ell=6P\ell$.

Replacing $2w\log\sin z$ by $4a\log z/z$ costs $10a$ per side:
the two terms
$(2a\cot(z/2)-4a/z)\log\sin z$ and
$4a\log(\sin z/z)/z$ have variations at most $3a,7a$.
Centrally put $e=\pi-z$, $s_L=q/2+a\tan(e/2)$,
$s_R=|q/2-a\tan(e/2)|$. Their $R$ terms have total variation
less than $21/2$: parameter motion costs at most
$3(1+\log4)<15/2$, separation at most three.
For $k_s=2(s_s-p/2)_+$, $k_s(0)=d_+$,
$|k_s'|\leq2a$, $|k_s-d_+|\leq2a\tan(e/2)$.
The variation of $(k_s-d_+)\log\sin e+d_+\log(\sin e/e)$ is
less than $9/2$ per side, since
$\int_0^{\pi/2}-\log\sin e\,de=(\pi/2)\log2$ and
$\tan(e/2)\cot e\leq1/2$. Thus the central cancellation variation
is less than twenty and its supremum less than four.
The central subtraction of $4a\log z/z$ costs $3a$ per side;
subtracting $d_+\log(\pi-z)$ initially costs $d_+\log2$ per side.
This proves the table. Centered-cell BV quadrature costs half the
variation divided by $m$.

The remaining numerical entries follow by integrating their
product derivatives. If $f_a=2a\cot(z/2)-4a/z$, then
$|f_a'|\leq a$, $|f_a|\leq az$ on the initial half, whence
\[
 \operatorname{Var}(f_a\log\sin z)
 \leq a\int_0^{\pi/2}\{-\log\sin z+z\cot z\}dz<3a.
\]
For $g_a=4a\log(\sin z/z)/z$, use
$0\leq1/z-\cot z\leq\pi z/6$ and
$0\leq-\log(\sin z/z)\leq z^2/3$. Differentiation gives
$|g_a'|\leq4a(\pi/6+1/3)$ and thus
$\operatorname{Var}g_a<7a$. These give $20a$ over two sides.
On the central half, set $v=p/2$, $S=3/2$. Each parameter integral
is at most $1+\log4<5/2$, so three monotone pieces cost less than
$15/2$. Monotonicity of $R_e$ in its parameter before integration
bounds each separation contribution by $(S+v)\log2<3/2$.
Their combined remainder variation is below $21/2$.
For the explicit logarithmic part the product rule gives, per side,
\[
 \begin{aligned}
 &\int_0^{\pi/2}|k_s'|(-\log\sin e)de
 +\int_0^{\pi/2}|k_s-d_+|\cot e\,de\\
 &\qquad +d_+\log(\pi/2)\\
 &\leq a\pi\log2+a\pi/2+d_+\log(\pi/2)<9/2.
 \end{aligned}
\]
Thus the central variation is below $21/2+9<20$. The other
subtracted function has variation at most
$4a(1+\log\pi)\int_{\pi/2}^{\pi}z^{-2}dz<3a$ per side,
adding $6a$ to the initial $20a$. The two initial central-log
subtractions add $2d_+\log2$. This accounts for every term in
$26a+20+2d_+\log2$.
The central cancellation supremum is below four; subtracting
$4a\log z/z$ adds at most $4a/\mathrm e$, since
$\sup_{z>0}\log z/z=1/\mathrm e$. Together with the pole bound
six, this gives $C_B^\phi$. On the initial holes integrate
\eqref{uc:endpoint-envelope}; concavity of
$c[7+(P/\pi)(5+2\log(4\pi/c))]$ bounds the two holes by the
$4/m$ term of $G_\phi$. The initial central-log term costs
$4d_+\log\pi/(\pi m)$, the central constant supremum is $C_B^\phi$,
and the coefficient holes cost $6LM_A^\phi/m$. This proves the
outer estimate.

\paragraph{The retained endpoint rows in finite coordinates.}
The endpoint charts use distances $w=r\pi+b$ of order one and
sample coordinates $z=j\pi+x$ of order one. Their actual angles are
$v=w/m$ and $u_j=z/n$. Thus the ratio $\mu=n/m$ survives in the
exact sine arguments. It is close to one, but its difference from one
creates terms of order $1/n$ after normalization and has to be kept.
The sinc function defined above provides formulas that stay meaningful
at removable zeros.
The factor $\tau(t)=t\cot t$ likewise has value one at zero.

The reference $f_j$ below is the limit obtained by replacing the
small-angle factors by their leading values while retaining the whole
sinc numerator. Its infinite absolute sum is $\ell_r$. The exact finite
row still has exactly $n$ indices of the prescribed parity. Near
indices are compared with this reference directly; far indices and
omitted tails are then bounded by reciprocal-square summation.

For the three retained endpoint charts $(r,b)=(1,y),(1,-y),(2,-y)$,
put $w=r\pi+b$, $z=j\pi+x$, $\mu=n/m$, $v=w/m$, $u_j=z/n$.
The $n$ actual indices have parity $1-r$ and $-\pi\leq u_j<\pi$.
With $\sigma_-=(-1)^{(r-j-1)/2}$, $\sigma_+=(-1)^{(j+r-1)/2}$,
\begin{equation}\label{uc:sinc-row}
 \begin{aligned}
 \frac{F_j}{nm}=
 \frac{c_e\sigma_-\operatorname{sinc}((z-\mu w)/2)\tau((u_j-v)/2)
       +c_d\sigma_+\operatorname{sinc}((z+\mu w)/2)\tau((u_j+v)/2)}
 {2m\sin(w/(2m))},\\
 f_j=\frac{c_e\sigma_-\operatorname{sinc}((z-w)/2)
                  +c_d\sigma_+\operatorname{sinc}((z+w)/2)}w.
 \end{aligned}
\end{equation}
Here the cosine weights use $(x,b)$; their product is still $a$.

For $b=-y$ the two cosine weights exchange their symmetric roles,
but their product is unchanged. The signs $\sigma_\pm$ encode the
actual parity of each index. They are part of the exact row formula;
changing them independently would alter which terms cancel. This is
also why the reference is continued by the sinc formula as a complete
linked expression. At the particular boundary value given below, all
its mass is concentrated in one removable sample.
Partial fractions show $\sum_{j\in\mathbb Z}|f_j|=\ell_r$;
the sinc formula supplies all linked removable values, including
$f_0=1$, $f_j=0$ otherwise at $(r,b,x)=(1,-\pi,0)$.
For nearby indices, all arguments lie in a fixed bounded range.
Subtracting factors in an order avoids introducing products of errors.
For example $ABC-A_0B_0C_0$ can be written
$(A-A_0)BC+A_0(B-B_0)C+A_0B_0(C-C_0)$.
Here the chosen order is prefactor, sinc displacement, then $\tau$.
The prefactor error is quadratic in $1/m$, the displacement uses
$|\operatorname{sinc}'|\leq1/2$, and the $\tau$ error is quadratic
in its small argument. The finite sums of $z^2$ and $|z|$ below
make each resulting charge explicit.

For $|j|\leq4$, use $(c_e+c_d)/w\leq1$,
$0\leq[2m\sin(w/(2m))]^{-1}-1/w\leq w/(20m^2)$,
$|\operatorname{sinc}'|\leq1/2$, and $1-\tau(t)\leq2t^2/5$
for $|t|<1$. The even five-index sets have
$\sum z^2\leq45\pi^2$, $\sum|z|\leq13\pi$; the odd four-index
set has bounds $24\pi^2,8\pi$. Subtracting prefactor, sinc displacement,
and tau in that order costs in total $7/n+317/n^2$.
The near-index arithmetic can be displayed row by row. The prefactor
error per index is at most
$(w/(20m^2))(c_e+c_d)\leq\pi/(5m^2)<2/(3n^2)$.
The sinc arguments move by $w/(2m)$, so their contribution is at
most $(c_e+c_d)/(4m)\leq1/(2n)$ per index. Finally
\[
 E_\tau\leq\frac1{10}\sum_{\rm near}(|z|/n+v)^2.
\]
For an even five-index set, $\sum z^2\leq45\pi^2$,
$\sum|z|\leq13\pi$, and $v\leq2\pi/n$, giving
$E_\tau\leq(45+52+20)\pi^2/(10n^2)<117/n^2$.
For an odd four-index set the corresponding numbers are
$24+32+16=72$, giving $72/n^2$. Including prefactors, each of the
two even rows is bounded by $5/(2n)+121/n^2$ and the odd row by
$2/n+75/n^2$. Their sum is
\[
 2\left(\frac5{2n}+\frac{121}{n^2}\right)
       +\frac2n+\frac{75}{n^2}=\frac7n+\frac{317}{n^2}.
\]
The estimates for $\tau$ used here follow directly from
$\sin t-t\cos t=\int_0^t r\sin r\,dr\leq t^3/3$ and
$\sin t\geq5t/6$ for $0\leq t\leq1$; evenness handles negative
arguments. Also $\operatorname{sinc}'t=-\int_0^1r\sin(rt)dr$,
so its global absolute bound is $1/2$.

For the remaining indices $|z|\geq4\pi\geq2w$, put
$D_h=\mu^2w^2-z^2$, $D'=w^2-z^2$,
$K=\operatorname{sinc}((|u_j|+v)/2)\operatorname{sinc}((|u_j|-v)/2)$,
$H_h=(1+\cos v)/(2K)$. Sine chords give $K>1/6$,
$|H_h-1|\leq u_j^2$, and
\[
 \frac\mu{D_h}-\frac1{D'}
 =\frac{z^2+\mu w^2}{m(z^2-\mu^2w^2)(z^2-w^2)}
 \leq\frac{20}{9mz^2}.
\]
For far indices the separation $|z|\geq2w$ keeps both reference
and actual denominators away from zero. The displayed difference of
reciprocal factors records the effect of $\mu\ne1$ exactly; it is
bounded by a summable multiple of $1/(mz^2)$. The factor $K$ measures
the two sine chords, and $H_h$ their residual geometric correction.
A positive lower bound for $K$ permits division. The summable decay
bound controls the aggregate contribution of the far indices.

The complete numerator is
$2a(1+\cos v)-p\sin u_j+\sin b\sin v$ over
$2(\cos u_j-\cos v)$. Using
$\sum|z|^{-1}\leq2H_n/\pi$, $\sum z^{-2}\leq3/(8\pi^2)$,
$w\sum z^{-2}<1/\pi$, the three far parts cost
$18a/n+(48pH_n+24q)/(\pi n)$; their omitted reference tails cost
$3a/n$ by $|f_j|\leq81a/(20z^2)$.
The closest endpoint row has a useful signed identity. For any real
finite vector $(v_j)$,
\[
 \sum_j|v_j|=-\sum_jv_j+2\sum_j(v_j)_+.
\]
If its signed sum is known exactly and its positive samples can be
located, this can be sharper than bounding every absolute value.
The next estimates use precisely this identity: the positive part is
bounded by a harmonic sum, and the signed row is given explicitly.
Its finite continuation at the origin therefore stays in the retained
term with the correct multiplicity.

For the far sums, the worst parity denominators are
$4\pi,6\pi,8\pi,\ldots$ on each side, giving
\[
 \sum_{\rm far}z^{-2}
 \leq\frac1{2\pi^2}\sum_{k\geq2}k^{-2}
 \leq\frac1{2\pi^2}\left(\frac14+\int_2^\infty r^{-2}dr\right)
 =\frac3{8\pi^2}.
\]
The other parity starts farther away. Finite reciprocal sums on either
side are bounded by $H_n/\pi$, and $w\leq2\pi$ gives
$w\sum z^{-2}<1/\pi$. For the $a$ numerator the actual expression
is $4\mu aH_h/D_h$, compared with $4a/D'$. The geometric error
$H_h-1$ costs at most $16a/(3n^2)$ per index, hence
$16a/(3n)$ per row. The remaining displacement costs
\[
 4a\sum\left|\frac\mu{D_h}-\frac1{D'}\right|
 \leq\frac{40a}{9\pi^2m}\leq\frac{2a}{3n}.
\]
Thus three rows cost $3(16/3+2/3)a/n=18a/n$.
Since $\cos v-\cos u=(u^2-v^2)K/2\geq u^2/16$,
the sine part per index is at most
$8p/(n|z|)+8qw/(nz^2)$. Three rows and the preceding sums give
$(48pH_n+24q)/(\pi n)$.
Outside the actual $n$-index range, $|z|\geq\pi n$ and the spacing
is $2\pi$. Integral comparison bounds the two tail sums by
$2/(\pi^2n)$. Also $w^2/z^2\leq1/256$, so
$|f_j|\leq1024a/(255z^2)\leq81a/(20z^2)$.
Each omitted reference tail therefore costs less than $a/n$, giving
$3a/n$ over the three charts.

For the closest row, positive samples require $\sin u>0$, and their
positive part is at most $p/(u-y/m)$. The first positive denominator
is at least $2x/n$, the others at least $2k\pi/n$, so twice their
normalized sum is at most $(1+H_n/\pi)/m$.
The exact signed row is $q\cot(y/(2m))/(2m)-a/m$ after sign reversal.
Thus its difference from $\operatorname{sinc}y$ is at most
$(2+H_n/\pi)/m+\pi^2/(24m^2)$.
These bounds yield
\[
 |J_{\rm end}-C_2|\leq
 \frac{9+21a+\{(1+48p)H_n+24q\}/\pi}{n}
                  +\frac{317+5/12}{n^2}.
\]
The two central rows have $|v-\pi|\leq\pi/m$,
$\sin(v/2)>99/100$, $|\cot(v/2)|\leq\pi/m$.
The weighted pole bound applied before taking linked limits, with
$|C|+|D|\leq P/2+a\pi/m$, gives
$0\leq J_{\rm cen}\leq(12+2L)/n$. This proves the retained estimate.

\paragraph{Mesh tails and the uniform bound.}
The remaining mesh terms come from the two explicit cell series, so
second-derivative estimates give their tails directly. At the upper
end, two offsets appear with opposite displacements $c$. Keeping their
logarithms paired uses
$\log(1+c/n)+\log(1-c/n)=\log(1-c^2/n^2)$, which is quadratic
in $1/n$. Estimating them separately first would leave two linear
errors. The other displayed logarithmic bound handles the terms whose
coefficients prevent that pairing. These produce $\mathcal M_\phi$
with no unidentified remainder.

The reciprocal and logarithmic-reciprocal centered-cell tails in
\eqref{uc:mesh-series} are at most $1/(8w^2)$ and
$(\log w+2)/(8w^2)$; integrate the second-derivative remainder to
obtain them. At the paired upper endpoints $n\pm c$, $|c|\leq1/2$,
keep their logarithms paired before taking absolute values:
$|\log(1-c^2/n^2)|\leq(4n^2-1)^{-1}$ and
$|\log(1\pm c/n)|\leq(2n-1)^{-1}$.
The precise centered-unit-cell identity behind these tail radii is
\[
 F(x)-\int_{x-1/2}^{x+1/2}F(y)dy
 =-\frac12\int_{-1/2}^{1/2}(1/2-|r|)^2 F''(x+r)dr.
\]
The kernel is bounded by $1/8$. For $F=1/x$, summing all cells
beyond an upper endpoint $w$ bounds the tail by
$(1/8)\int_w^\infty2x^{-3}dx=1/(8w^2)$.
For $F=\log x/x$, one has
$|F''|\leq(2\log x+3)/x^3$ on $x\geq3/2$, and
\[
 \frac18\int_w^\infty\frac{2\log x+3}{x^3}dx
 =\frac{\log w+2}{8w^2}.
\]
At the paired upper endpoints put
$\eta_\pm=\log(1\pm c/n)$, $j_c=\eta_++\eta_-$, and
$\epsilon_n=\log(m/n)$. Inserting the finite cell sums in the
mesh remainder and cancelling powers of $L$ gives exactly
\[
 R_{\rm mesh}=-B_0\epsilon_n(2L+\mathcal H_2)
 +(D_0+B_0L-B_0\epsilon_n)j_c
 +\frac{B_0}{2}(\eta_+^2+\eta_-^2)+\mathrm{Tail}.
\]
The first three terms are bounded, respectively, by
\[
 \frac{B_0(2L+2r_0)}n,\qquad
 \frac{|D_0|+B_0(L+1/n)}{4n^2-1},\qquad
 \frac{B_0}{(2n-1)^2}.
\]
Both upper tails together are at most
\[
 \frac{|Z_n|+B_0\{\log(n+1/2)+2\}}{4(n-1/2)^2}.
\]
These four expressions give each line of $\mathcal M_\phi$ after
using $B_0=ab_0$ and the stated coefficient envelopes.
For $R_{\log}$, the variation of $\log(\pi-z)$ on either cell
union is at most $\log(2m)$; the two midpoint errors together are
therefore at most $\log(2m)/m$. Initial holes cost
$4\log\pi/m$. If the two central hole lengths sum to $2h$,
concavity of $r(1-\log r)$ bounds their log integrals by
$2h(1-\log h)$, hence $2(1+\log(m/\pi))/m$ after normalization.
Adding gives precisely $\pi W_n$.

Since $|\mathcal H_2|\leq2r_0$, $|D_0|\leq\overline D_\phi$,
$|Z_n|\leq\overline Z_\phi$, these are exactly $\mathcal M_\phi$.
Finally BV quadrature of $\log(\pi-z)$ costs $\log(2m)/m$;
its initial holes cost $4\log\pi/m$, central holes at most
$2(1+\log(m/\pi))/m$. This proves the last bound in
\eqref{uc:absolute-components}, and hence the first part of
\eqref{uc:global-error}.

To obtain its uniform constant, use $\Gamma_\phi\leq g_n$,
$\ell\leq\log m$, $C_B^\phi<8$ and
$\overline D_\phi<49/15<4$, $\overline Z_\phi<5$.
Termwise maximization gives the separated sum
$C_{E,{\rm sep}}=\overline C_E+3H/n$. Retain the correlations:
\[
 C_{E,{\rm sep}}-C_{E,\phi}>
 \left(\frac{52H}{\pi^2n}+\frac{28L}{\pi^2m}\right)(1-a)
                +\frac{48H}{\pi n}(1-p).
\]
The uniform bound keeps one final correlation. Replacing $a$ and
$p$ independently by one gives the separated upper sum, but the same
phase cannot have $a=1$ and $p=1$: the first requires $x=y=0$,
while the second requires $x=\pi/2$. The quantified saving uses
$1-a\geq(1-\cos x)/2$ together with $p=\sin x$. The lower bound
on $A(1-\cos x)+B(1-\sin x)$ follows from
$A\cos x+B\sin x\leq\sqrt{A^2+B^2}$. Thus the saving is
uniform over the full triangle and permits the simpler allowance
$\overline C_E$ without losing the needed margin.

Using $L>3$, $n/m\geq32/33$, $L/H>3/4$, its ratio to $H/n$
is at least $A(1-\cos x)+B(1-\sin x)$, where
$A=398/(11\pi^2)>18/5$, $B=48/\pi>15$.
The minimum is $A+B-\sqrt{A^2+B^2}>3$, since
$2(A-3)(B-3)>9$. This proves precisely the saving $3H/n$.
Each rational logarithmic term decreases for $n\geq32$;
for $p>0$, $m\Gamma_\phi'<p\leq\Gamma_\phi$, while at $p=0$
$\Gamma_\phi$ is identically zero. Together with
$m\ell'=\epsilon\cot\epsilon\leq1<\ell$, this proves the two less
immediate cases. Splitting the mesh terms into constants and logarithms
over their positive quadratic denominators proves their decrease.
The model bounds, $P\leq2$, and phase folding now give
\eqref{uc:analytic-norm}, including complete signed cardinal sums by
the finite reflection formula.

Monotonicity of the allowances is useful because a later argument
can check a lower degree and inherit the same upper bound for all
larger degrees. For a typical term $(\log n+c)/n$, differentiation
gives $(1-\log n-c)/n^2$, which is negative on the stated range.
The expressions involving $\ell$ and $\Gamma_\phi$ need the extra
derivative inequalities recorded above because their numerators also
depend on $m$. With the global model bounds, $P\leq2$ turns
$K_DP$ into $2K_D$, yielding the stated analytic norm. The signed
phase extension uses absolute phases separately, preserving the two
row directions and their finite reflection multiplicities.

\paragraph{The actual origin.}

The origin bound concerns the actual finite sum $S_n(0)$, including
its removable samples. We need a lower bound for
$E_n(0)=S_n(0)-M_\infty(0)$ because the relative comparison will
contain $-uE_n(0)$ with $u\geq0$. A lower bound for $E_n(0)$ gives
an upper bound for this contribution. This one-sided information can
be appreciably sharper than a global absolute error estimate.

Several terms simplify at zero. The coefficient perturbation vanishes,
$d_+=0$, and the four endpoint reference samples have fixed simple
coordinates. The two actual central rows have an exact positive total.
Keeping their signs and the sign of the outer quadrature error gives
a smaller origin allowance. This is why the following calculation is
made at the finite origin before any phase perturbation is estimated.
At zero, $|E_{\rm row}|\leq D$, $|R_{\rm mesh}|\leq\mathcal M$.
The four positive endpoint samples, with multiplicity, are
$(w,z)=(\pi,0),(\pi,0),(2\pi,\pi),(2\pi,-\pi)$.
The two central rows have exact total $J_{\rm cen}(0)=2/m$.
Their linked cyclic vectors below have respectively one sample $n$
and nonnegative cotangent differences. The identity
$\sum_k\cot((k\pi+\theta)/n)=n\cot\theta$ gives signed sum $n$
for the latter row as well.
Consequently
\[
 E_{\rm ret}(0)=-2/m+8\sum(\mu/D_h-1/D')
                         +8\mu\sum(H_h-1)/D_h.
\]
The two sums in this expression represent different corrections.
The first keeps the displacement $\mu=n/m$ of the actual denominator;
the second keeps the sine-chord correction $H_h-1$. At these four
samples $w>|z|$ and the denominators are positive, so the first
correction has a favorable positive sign. The exact identity below
separates its two reciprocal factors and makes this sign transparent.
The geometric correction is smaller, of order $n^{-2}$, and can be
bounded in absolute value. The resulting retained-row lower bound
therefore retains the positive term $22/(45m)$.

The exact identity
\[
 \frac\mu{D_h}-\frac1{D'}=\frac1{2m}
 \left\{\frac1{(\mu w-z)(w-z)}+
              \frac1{(\mu w+z)(w+z)}\right\}
\]
has bracket sum at least $56/(9\pi^2)$, by evaluating its decreasing
function of $\mu$ at one. At the four origin samples the reciprocal bracket can be summed
without estimating each term:
\[
 S_{\rm disp}(0)=\frac1{\pi^2}
   \left\{\frac4\mu+\frac2{2\mu-1}+\frac2{3(2\mu+1)}\right\}
 \geq\frac1{\pi^2}\left(4+2+\frac29\right)=\frac{56}{9\pi^2}.
\]
Each summand decreases with $\mu\in(1/2,1]$, proving the inequality.
The exact reciprocal identity multiplies this by $4/m$ in
$E_{\rm ret}(0)$, and the central-plus-endpoint signed constant is
$-2/m$. Since $\pi^2<10$,
$4(56/(9\pi^2))-2>22/45$. This is the source of the favorable
$22/(45m)$ before the geometric error is charged.

For these four samples
$K\geq49/50$, $|H_h-1|\leq19/n^2$, $\sum D_h^{-1}\leq25/54$;
hence the geometric error is at most $71/n^2$ and
$E_{\rm ret}(0)>22/(45m)-71/n^2$.
For the signed outer term define
\[
 \begin{gathered}
 b(z)=2\cot(z/2)-4/z,\quad d(z)=[2\sin(z/2)]^{-1}-1/z,
 \\ P_0(z)=2(f(z)-f(0))/z+2f(z)d(z),\quad
 Q_0(z)=2\cot(z/2)\log\sin z-4\log z/z,\\
 F_0=(L/\pi)b+P_0+Q_0/\pi.
 \end{gathered}
\]
The function $F_0$ is the signed outer remainder at zero. Its
components are already regularized: $b$ subtracts the reciprocal
cotangent pole, $d$ subtracts the reciprocal sine pole, $P_0$ subtracts
the scalar endpoint value, and $Q_0$ subtracts the logarithmic-reciprocal
term. Their apparent quotients have finite endpoint limits after the
subtractions. The derivative variation estimate is made on a truncated
interval because the remaining logarithmic derivatives grow near the
ends.

If a function has a derivative of bounded variation, centered-cell
quadrature gains a second power of the cell width. In usable form,
the error on a cell can be written as an integral of a continuous
quadratic kernel against the change measure of its derivative; that
kernel has magnitude at most $h^2/8$. Summing disjoint cells then
bounds the total error by $h^2/8$ times the derivative's total
variation, with the normalization of the outer integral included.
This gives the $O((\log n)/n^2)$ outer origin bound below.

On $[h/2,\pi-h/2]$,
\begin{equation}\label{uc:origin-variation}
 \operatorname{Var}F_0'\leq L/2+(25/2)\log(2m)+37/3=:V.
\end{equation}
Here are component estimates proving this bound. The cotangent
residual has $|(\cot z-1/z)''|<1$ on $[0,\pi/2]$;
$0\leq d\leq1/5$, $0\leq d'\leq1/3$, $|d''|\leq5/32$.
Thus $\operatorname{Var}(b'/\pi)\leq1/2$.
From $f'(0)=0$, $|f'|<2$, $|f''|<1$ one obtains
$|(2(f-f(0))/z)''|\leq8/z$ and
$\int_0^\pi|(2fd)''|<9$. Hence
$\operatorname{Var}P_0'\leq8\log(2m)+9$.
Initially write $Q_0=b\log\sin z+4\int_0^1(\cot(tz)-(tz)^{-1})dt$.
The second derivatives of the first and last regular factors cost one
and three; the other product terms cost $(9/2)\int dz/z$.
Centrally, $2\tan(e/2)\log\sin e$, $e=\pi-z$, costs
$4+9\int de/e$, while $-4\log z/z$ costs two.
Their derivatives match at $\pi/2$, giving
$\operatorname{Var}Q_0'\leq(27/2)\log(2m)+10$ and
\eqref{uc:origin-variation} after division by $\pi>3$.

\paragraph{Derivative calculations for the signed origin bound.}
These component bounds can be checked using ordinary derivatives.
For $c(z)=\cot z-1/z$, the preceding partial-fraction estimate gives
$|c''|<1$. Since $b(z)=2c(z/2)$ and
$d(z)=\{c(z/4)-c(z/2)\}/2$,
\[
 |b''|\leq\tfrac12,\qquad |d''|\leq\tfrac1{32}+\tfrac18=\tfrac5{32}.
\]
Also $d$ increases and $0\leq d\leq1/2-1/\pi<1/5$.
Writing $x=z/2$, direct differentiation gives
\[
 d'(z)=\frac{\sin^2x-x^2\cos x}{4x^2\sin^2x}.
\]
Its numerator is nonnegative: differentiating it gives
$-2\cos x(x-\sin x)+x^2\sin x\geq x^3(2/\pi-1/3)\geq0$,
using $x-\sin x\leq x^3/6$ and $\sin x\geq2x/\pi$.
For its upper bound use
$\sin^2x-x^2\cos x\leq x^2(1-\cos x)\leq x^4/2$ and
$\sin x\geq2x/\pi$ to get $d'\leq\pi^2/32<1/3$.

Set $r_0^*(z)=2(f(z)-f(0))/z$. Its second derivative is
\[
 (r_0^*)''=\frac{2f''}z-\frac{4f'}{z^2}
                            +\frac{4(f-f(0))}{z^3}.
\]
The inequalities $|f'|\leq z$ and $|f-f(0)|\leq z^2/2$
therefore give $|(r_0^*)''|\leq8/z$.
For the product term,
\[
 |(2fd)''|\leq2|f''|d+4|f'|d'+2|f||d''|.
\]
Its integral is bounded by
$2\pi/5+8/5+5\pi/4<9$, even using the weaker valid bound
$|f|\leq4$ in its last summand. Hence the two parts of $P_0$
give $8\log(2m)+9$.

On the initial half, write $Q_0=b\log\sin z+g$ with
$g=4\int_0^1c(tz)dt$. Thus $|g''|\leq4/3$, whose integral
is below three. The remaining second derivatives are
\[
 (b\log\sin z)''=b''\log\sin z+2b'\cot z-b\csc^2z.
\]
The first term has integral below one. Since $|b'|\leq1$,
$|b|\leq z$, $\cot z\leq1/z$, and
$\csc^2z\leq\pi^2/(4z^2)<5/(2z^2)$, the last two have
combined integral at most $(9/2)\int dz/z$.
On the central half put $e=\pi-z$, $T=2\tan(e/2)$.
In $(T\log\sin e)''$, the terms $T''\log\sin e$,
$2T'\cot e$, and $-T\csc^2e$ cost respectively less than four,
$4\int de/e$, and $5\int de/e$, using
$T\leq2e$, $T'\leq2$, $T''\leq2$. The remaining function
$-4\log z/z$ has second derivative bounded by $12/z^3$ and integral
$18/\pi^2<2$ on that half. Adding gives
$(27/2)\log(2m)+10$. Finally
\[
 \frac L2+8\log(2m)+9+
   \frac1\pi\{(27/2)\log(2m)+10\}
 \leq\frac L2+\frac{25}{2}\log(2m)+\frac{37}{3}=V.
\]
This is the complete component arithmetic of the origin derivative
variation estimate.
Centered-cell quadrature for a BV derivative has kernel at most
$h^2/8$; the two cell families therefore cost at most
$\pi V/(4m^2)$ downwards. On initial holes $z\leq5h/2$,
$d(z)<z/16$, $P_0<7z/6$, and
$F_0(z)\leq7z/6+(z/\pi)\log_+(\pi/(2nz))$.
Their total upper integral costs
$119\pi/(24m^2)+1/(8n^2)$, using
$\int_0^c z\log_+(A/z)dz\leq A^2/4$.
The omitted intervals at the initial endpoint are handled with an
upper integral bound because we need a lower bound for the
sample-minus-integral error. A positive omitted integral makes that
error more negative. The bound involving $\log_+$ restricts the
logarithm to the region where it can increase this upper integral;
here $\log_+r=\max(\log r,0)$. Its integral is elementary, and
its support shrinks like $1/n$. At the central endpoint the remainder
has the opposite sign, as checked next. Omitting a negative integral
helps the desired lower bound, so there is no need to replace it by
an absolute value.

The central holes have favorable sign: there $z\geq21\pi/22$,
$b/\pi<-7/20$, $P_0<101/100$, $Q_0\leq0$, so $F_0<-1/25$.
Thus
\[
 E_{\rm out}(0)\geq-B_{\rm out},\quad
 B_{\rm out}=\frac{\pi[L+25\log(2m)+193/3]}{8m^2}+\frac1{8n^2}.
\]
The hole constants follow from their actual lengths $3h/2,5h/2$.
The first initial envelope integrates to
\[
 \frac{7}{12\pi}\{(3h/2)^2+(5h/2)^2\}
 =\frac{119\pi}{24m^2}.
\]
The logarithmic envelope contributes at most
$2(\pi/(2n))^2/(4\pi^2)=1/(8n^2)$.
For the central sign, $v=(\pi-z)/2\leq\pi/44<1/14$ gives
$\sin(z/2)=\cos v>391/392$ and
$\cot(z/2)=\tan v<28/391<1/13$.
Since $z\geq21\pi/22$ and $\pi<22/7$,
\[
 \frac b\pi<\frac2{39}-\frac{49}{121}<-\frac7{20},
 \qquad P_0=\frac{f(z)}{\sin(z/2)}-\frac{2f(0)}z<\frac{101}{100}.
\]
Here $Q_0\leq0$ and $L>3$, so
$F_0<-21/20+101/100=-1/25$. The central omitted integrals are
therefore favorable. Combining the midpoint term with the initial holes,
\[
 \frac{\pi V}{4m^2}+\frac{119\pi}{24m^2}
 =\frac\pi{8m^2}\{L+25\log(2m)+193/3\},
\]
which explains every part of $B_{\rm out}$.

Together these prove
$E_n(0)\geq-D-\mathcal M-B_{\rm out}-[71/n^2-22/(45m)]_+$.
The bracket vanishes for $n\geq148$. Moreover $nB_{\rm out}$ decreases:
for $K=L+25\log(2m)+193/3$ its derivative is
$\pi[(1-n)K+m+25n]/(8m^3)-1/(8n^2)<0$, since $K>64$.
At $160$, $\log160,\log322<6$ gives $160B_{\rm out}<1$.
This proves \eqref{uc:origin-lower}.

The threshold $160$ has a specific arithmetic role. By that degree
the positive retained-row displacement has paid for its $71/n^2$
geometric allowance, while the signed outer contribution is less than
$1/n$. The row and mesh terms remain bounded by $D$ and
$\mathcal M$. Their sum is exactly the definition of $\bar e$.
The proof uses no assertion that the full error is positive at zero;
it gives the lower bound needed for the later weighted subtraction.

\paragraph{Relative charges and their two radii.}

We now subtract $a$ times each origin term before estimating its
change. Every exact increment $\delta_aX=X(x,y)-aX(0,0)$ vanishes
at zero. This eliminates common finite corrections that would otherwise
remain as phase-independent errors. The few remaining constants in the
table are explicit allowances for rows or tails that were estimated
more broadly; they are not new values of the exact increment.

The two domains have the same row sets but different lower bounds for
their denominators. On $\rho\leq1/2$ the retained sample distances
stay farther from their poles, giving the smaller constants. On
$\rho\leq1$ the same proof needs larger constants. Keeping these
two radii allows the later localization argument to use the sharp
near-origin comparison and still cover the surrounding region.

A charge below means an upper bound assigned to one piece of the
exact discrepancy. The absolute-value signs in the table identify
pieces estimated in both directions. The row and signed outer entries
need only upper bounds, since the objective is an upper bound for
$S_n-\Lambda_c$. In the last entry the central logarithm is included
with its original coefficient, preserving its contribution exactly once.
The same telescope, with its origin subtracted before estimation,
gives the following table. Every entry is valid for $n\geq32$;
the two columns retain their distinct finite-row denominators.
\begin{equation}\label{uc:relative-components}
\begin{array}{c|c|c}
 &\rho\leq1/2&\rho\leq1\\ \hline
\delta_aE_{\rm row}&\multicolumn{2}{c}{aA_4+20/n^2+\rho(8H+10)/n}\\
|\delta_aE_{\rm ret}|&a\rho J_N+P(7/n+4/n^2)
                         &a\rho J_W+P(11/n+4/n^2)\\
|\delta_aR_{\rm mesh}|&\multicolumn{2}{c}{\rho/(2n)}\\
\delta_aE_{\rm out}+(d_+/\pi)R_{\log}
 &\rho((23/3)L+47)/n+yd_\beta
 &\rho((23/3)L+50)/n+yd_\beta
\end{array}
\end{equation}
We derive these entries, in particular the signed last one.
On the wide domain $c_e,c_d\geq7/8$, $a\geq49/64$ and
$2a-P\geq17/32$. For the endpoint reference use fixed label sets
$j=2,\ldots,\lfloor n/2\rfloor$ on the left and
$j=3,\ldots,\lfloor n/2\rfloor$ on the right. Put
$\xi=x/\pi$, $s=\pm y/\pi$, $\nu=j+s$, $q_\pm=(1+s\pm\xi)/2$.
The normalized endpoint comparison in \eqref{uc:reference-match} is
$a e_j$, where
\[
 \begin{gathered}
 e_j=\frac2{\pi^2}\left\{\frac{\tau_j-1}{\nu}B_j
                              +\frac{\tau_j}{\nu}U_j\right\},
 \quad\tau_j=\frac{\pi\nu}{2m}\cot\frac{\pi\nu}{2m},\\
 U_j=\sum_\pm[\psi(j+q_\pm-\nu/(2m))-\psi(j+q_\pm)]
                                      -2\log(n/m).
 \end{gathered}
\]
For this comparison the row labels $j$ are fixed while the phase
moves along a straight path from $(0,0)$ to $(x,y)$. One may parametrize
it by $(\lambda x,\lambda y)$, $0\leq\lambda\leq1$; a dot denotes
differentiation with respect to $\lambda$. Fixed labels ensure that
differentiating a finite row sum differentiates the same terms at every
point of the path. A cutoff based on a moving pole location would
create membership changes requiring extra endpoint terms.

The bounds for $\psi'$ and $\psi''$ come from differentiating its
positive reciprocal series:
$\psi'(z)=\sum_{k\geq0}(k+z)^{-2}$ and
$\psi''(z)=-2\sum_{k\geq0}(k+z)^{-3}$. The lower argument bound
$z\geq j$ allows integral comparison to supply the stated uniform
estimates. The four prefactor inequalities then implement the ordinary
product rule for $e_j$. Integrating $|\dot e_j|$ from zero to one
bounds the actual phase increment.

Along the straight phase path, $5j/6\leq\nu\leq7j/6$,
$j+q_\pm\geq j+1/3$, and the displaced arguments are at least $j$.
The positive psi derivative series gives
$\psi'\leq3/(2j)$, $|\psi''|\leq2/j^2$; hence, with
$b=\xi+|s|$, $|\dot B_j|\leq4b/j$, $|\dot U_j|\leq3b/(mj)$,
$|B_j|\leq2/\nu$, $|U_j|\leq5/n$.
The four prefactor bounds are
$|({\tau_j-1})/\nu|\leq4\nu/m^2$,
$|[({\tau_j-1})/\nu]'|\leq12/m^2$,
$|\tau_j/\nu|\leq1/\nu$, $|(\tau_j/\nu)'|\leq3/\nu^2$.
They yield
\[
 |\dot e_j|\leq\frac{2b}{\pi^2}
 \{29/(m^2j)+19/m^2+22/(nj^2)+4/(mj^2)\}.
\]
Using harmonic sum at most $2H_n$, square sum at most $3/2$, and
$H_n\leq n/4$, its sum is at most $145\rho/(27n)<8\rho/n$.
The infinite tails cost $2a\rho/n$, by differentiating $\Xi_j$.

\paragraph{The four product-rule coefficients.}
Let $F_j=(\tau_j-1)/\nu$, $G_j=\tau_j/\nu$, and
$c_*=2/\pi^2$, so $e_j=c_*(F_jB_j+G_jU_j)$. Along the path,
$|\dot\nu|=|s|\leq b$ and the four differentiated products have
bounds
\[
 \begin{aligned}
 |\dot e_j|/c_*
 &\leq |F_j'|b|B_j|+|F_j||\dot B_j|
                  +|G_j'|b|U_j|+|G_j||\dot U_j|\\
 &\leq b\left\{\frac{24}{m^2\nu}+\frac{16\nu}{m^2j}
                 +\frac{15}{n\nu^2}+\frac3{mj\nu}\right\}.
 \end{aligned}
\]
Using $5j/6\leq\nu\leq7j/6$ changes these four coefficients to
$144/5<29$, $56/3<19$, $108/5<22$, and $18/5<4$,
respectively. This produces the displayed derivative bound.
To justify the ingredients, $\partial_s B_j$ has magnitude at most
$3/(2j)+2/\nu\leq4/j$, and its $\xi$ derivative at most
$3/(2j)$. In a displaced $\psi$ difference, the derivative changes
by at most $2\delta/j^2$ with $\delta=\nu/(2m)$, by the
$\psi''$ bound. Differentiating both its argument and $\delta$
gives $|\partial_\xi U_j|\leq2/(mj)$ and
$|\partial_s U_j|\leq3/(mj)$. These imply the path bounds used.
The prefactor derivatives follow by differentiating $u\cot u$ with
$u=\pi\nu/(2m)<\pi/4$; the inequalities
$0\leq1-u\cot u\leq u^2$ and $|(u\cot u)'|\leq2u$
give the relaxed constants $4,12,1,3$ listed above.

Summing over the two fixed label sets gives
\[
 \sum_j|\dot e_j|
 \leq\frac{2b}{\pi^2}
       \left\{\frac{58H_n}{m^2}+\frac{25}m+\frac{33}n\right\}
 \leq\frac{145\rho}{\pi^3n}<\frac{145\rho}{27n}.
\]
Here fewer than $m$ labels multiply the constant $19/m^2$,
$\sum j^{-1}\leq2H_n$, $\sum j^{-2}\leq3/2$, and
$H_n\leq n/4$. For the missing infinite tails,
\[
 \left|\frac d{d\lambda}\frac{\Xi_j}a\right|
 \leq\frac{2b}{\pi^2}
             \left(\frac4{j\nu}+\frac2{\nu^3}\right)
 \leq\frac{12b}{\pi^2j^2}.
\]
Summing both tails from $J+1$ with $J=\lfloor n/2\rfloor$ costs
$24a\rho/(\pi^3J)\leq2a\rho/n$. No phase-dependent label was
added or removed in these differentiations.

We also estimate the increment of the half-line residual. Define $r(0)=0$,
$r(w)=\cot w-1/w$ for $0<w<\pi/2$, and $r(w)=-1/w$ afterwards.
Direct differentiation gives
$\operatorname{Var}r'=5/3$, $\int|r'|=4/\pi$.

The variation numbers have a simple endpoint calculation.
On $(0,\pi/2)$, $r''<0$, $r'(0)=-1/3$, and
$r'(\pi/2-)=-1+4/\pi^2$. At $\pi/2$ its derivative jumps upward
by one to $4/\pi^2$, then decreases to zero on the tail. Thus
\[
 \operatorname{Var}r'=(2/3-4/\pi^2)+1+4/\pi^2=5/3.
\]
The function itself decreases from zero to $-2/\pi$ and increases
back to zero, so $\int|r'|=2/\pi+2/\pi=4/\pi$. These are
statements about its derivative and matched differences; its separate
$-1/w$ tail integral diverges.
The piecewise definition of $r$ is continuous at $\pi/2$ because
$\cot(\pi/2)=0$. For $w\geq\pi/2$, its tail is $r(w)=-1/w$
and its derivative is $r'(w)=1/w^2$. Thus the derivative has an
integrable tail and finite total variation, including the finite
derivative change at the joining point. The matched-cell defect
subtracts an integral from the sample in each cell before summing to
infinity. The integrability of $r'$ makes the resulting series converge;
its variation also controls how the defect changes with phase.
The two estimates have different uses: one gives the correction at a
fixed offset, while the other controls its increment as the row and
the fractional offset move.

For its matched-cell defect $\mathcal E(h_*,v,t)$,
defined by the absolutely convergent series
\[
 \sum_{k\geq0}\left\{r(v+h_*(k+t))
       -h_*^{-1}\int_{v+kh_*}^{v+(k+1)h_*}r(w)\,dw\right\},
\]
\[
 \begin{gathered}
 |\mathcal E-(1/2-t)r(v)|\leq5h_*/24,\\
 |\mathcal E(v,t)-\mathcal E(v_0,t_0)|
 \leq(5/3)|v-v_0|+(4/\pi+5h_*/3)|t-t_0|.
 \end{gathered}
\]
The cell kernel is piecewise linear with one downward jump. Its
mean records the first-order offset bias. After removing this mean,
its primitive is bounded by $1/8$, independently of the offset.
Integrating that primitive against the finite change measure of $r'$
gives a bound proportional to $h_*\operatorname{Var}r'$. The phase
increment is bounded by differentiating finite matched sums and using
$\int|r'|$ and $\operatorname{Var}r'$. Their common summable
majorants justify passage to the infinite limit. Thus a derivative
jump is included in these estimates just as an ordinary derivative
change would be.

Indeed the cell kernel $K_t(w)=w-\mathbf1_{w>t}$ has mean
$t-1/2$ and a zero-mean primitive of modulus at most $1/8$;
integration against $dr'$ proves the first bound.

In detail, put $\mu_t=t-1/2$ and define on the unit cell
\[
 B_t(s)=s^2/2-\mu_ts-(s-t)_+.
\]
Then $B_t'=K_t-\mu_t$, $B_t(0)=B_t(1)=0$, and
$|B_t|\leq1/8$. The positive extremum is $t(1-t)/2\leq1/8$;
the possible negative extremum is no smaller than
$-(t-1/2)^2/2\geq-1/8$. The exact first-order cell identity gives
\[
 r(v+h_*(k+t))-h_*^{-1}\int_{v+kh_*}^{v+(k+1)h_*}r
 =\int_{v+kh_*}^{v+(k+1)h_*}
               K_t((w-v-kh_*)/h_*)r'(w)dw.
\]
The mean part sums to
$\mu_t\{r(v+Nh_*)-r(v)\}$ over $N$ cells. Integrating the
zero-mean part by parts gives an error at most
$h_*\operatorname{Var}r'/8=5h_*/24$. Let $N\to\infty$;
$r(v+Nh_*)\to0$, leaving $(1/2-t)r(v)$.
Differentiation of a finite matched sum in $v$ gives a matched defect
of $r'$, bounded by $\operatorname{Var}r'=5/3$. Differentiation in
$t$ gives $h_*\sum r'$; comparison with its integral bounds its
absolute value by $\int|r'|+h_*\operatorname{Var}r'$. These are
the two Lipschitz constants stated for $\mathcal E$. Differentiation of
finite matched sums gives the second, and absolute convergence permits
the limit. Here the unwrapped actual offsets are
$\theta_\pm=(1-s\pm\xi)/2+(j+s)/(2m)$, strictly in $(0,1)$.
They satisfy $|v-v_0|=y/m$ and
$\sum|\theta_\pm-\theta_0|\leq\rho/\pi$, $v_0=\pi j/m$.
Since $|r(v)|\leq\pi v/6$,
$\cot(v/2)\leq8/(3v_0)$, and
$|\cot(v/2)-\cot(v_0/2)|\leq8\pi^2y/(9mv_0^2)$,
the single-row residual increment is at most
\[
 a\left\{\frac{16\rho m}{3\pi^2j}
 +\left(\frac{4\pi^2}{27}+\frac{80}{9\pi}\right)\frac yj
 +\frac{10\pi m}{27n}\frac y{j^2}\right\}.
\]
Normalized summation bounds it by $2a\rho H_n/n$
(the coefficient is below $32/27+32/51+85/612<2$).

This bound follows because the combined harmonic sum is at most $2H_n$
and the combined square sum is at most $3/2$. For example, dividing the three displayed
single-row terms by $nm$ yields a coefficient below
\[
 \frac{32}{27}+\frac{32}{51}+\frac{85}{612}<2
\]
for $a\rho H_n/n$, using $y\leq\rho$, $m/n\leq17/16$,
$m\geq17$, $H_n\geq1$, and $3<\pi<4$. The assumptions here are
weaker than $n\geq32$, so they apply on both radii.
The complementary central references cost $4a/n$, the ordinary
perturbation $6PH_n/n$, and the two pole errors
$20/n^2+40aH/(\pi n^2)$. These give the first table row.

Indeed the endpoint reference motion and its infinite tails give
$a\rho(2H_n+10)/n$. Add $4a/n$ for the complementary central
references, $6PH_n/n$ for the ordinary perturbation, and the two
weighted-pole allowances $20/n^2+40aH/(\pi n^2)$. Since
$a\leq1$, $P\leq\rho$, and $H_n\leq H$, their sum is bounded by
\[
 a\left(\frac4n+\frac{40H}{\pi n^2}\right)+\frac{20}{n^2}
                  +\frac{\rho(8H+10)}n,
\]
which is the first row with $A_4$ as defined.

\paragraph{Why the finite endpoint sets stay exact.}
The next step locates the positive samples of the retained endpoint
rows before differentiating their contributions. The bound
$2a-P\geq17/32$ keeps the relevant numerators strictly positive.
The lower bounds for $w\pm z$ and $\mu w\pm z$ protect their
denominators on each of the two phase domains. The first excluded
index has the opposite denominator sign, proving that no positive
sample has been missed. This exact membership is needed for the signed
row identity: an extra or missing sample would change an order-$1/n$
term, the scale at which the relative comparison operates.

\paragraph{The exact endpoint error before differentiation.}
The signed-row sum and its positive samples give, for each of the three
charts, the finite identity
\[
 \begin{aligned}
 e_r&:=R_r/(nm)-\ell_r\\
 &=-a/m+\frac{\sin b}{w}\{\tau(v/2)-1\}
       +8a\sum\left(\frac\mu{D_h}-\frac1{D'}\right)\\
   &+8\mu a\sum\frac{H_h-1}{D_h}
       +2\mu\sum\frac{g_j}{D_hK},\\
 g_j&=-p\sin u_j+\sin b\sin v.
 \end{aligned}
\]
The sums here run over that chart's positive samples. The closest row
has $e_0=-a/m+\operatorname{sinc}y\{\tau(y/(2m))-1\}$.
Define $S_{\rm disp}$ to be the sum of
$[(\mu w\pm z)(w\pm z)]^{-1}$ over the four samples and both
signs, and $Q_H=\sum(H_h-1)/D_h$. The exact origin subtraction is
\[
 \delta_a(J_{\rm end}-C_2)
 =\frac{4a}m\{S_{\rm disp}(\phi)-S_{\rm disp}(0)\}
       +8\mu a\{Q_H(\phi)-Q_H(0)\}+T_\tau+T_g.
\]
This identity explains why $a$ is held outside the differentiated
quantities. The four $-a/m$ terms cancel exactly; $T_\tau,T_g$
vanish at the origin. We next bound these four remaining terms.

The positive endpoint sample sets on $\rho\leq1$ are exactly
$\{0\},\{0\},\{-1,1\}$ in \eqref{uc:sinc-row}; their full
numerators are positive because $2a-P\geq17/32$. On the wide domain
$w\pm z\geq(17/8)k$, $\mu w\pm z\geq(19/10)k$,
$k=r\pm j\in\{1,3\}$; the near bounds are $(21/8)k,(12/5)k$.
The first excluded index has $|z|>w>\mu w$, proving completeness.
In the exact row error, the four terms $-a/m$ cancel upon subtracting
the actual origin. The remaining displacement, geometric, sine and
tau charges are respectively
\[
\begin{array}{c|cccc}
 &\text{displacement}&\text{geometric}&\text{sine}&\text{tau}\\ \hline
\rho\leq1/2&16a\rho/(5n)&78a\rho/n^2&7P/n&4P/n^2\\
\rho\leq1&6a\rho/n&148a\rho/n^2&11P/n&4P/n^2.
\end{array}
\]
For explicit checks, a reciprocal-product derivative is at most
$(25760/104329)\rho/k^3$ on the wide domain. The eight $k$ values
are six ones and two threes, giving displacement-sum increments
$3\rho/2$ wide and $4\rho/5$ near. Also
$K\geq49/50$, $|H_h-1|\leq19/n^2$, $|\dot H_h|\leq5\rho/n^2$,
$|\dot D_h^{-1}|\leq(20/19)\rho/D_h$, and
$\sum D_h^{-1}\leq800/1083$ wide, $25/54$ near.
Differentiating $H_h/D_h$ gives exactly the geometric column.
The sine column follows by inserting the same denominators into the
last numerator in \eqref{uc:sinc-row}; the tau column uses
$\sin y\geq5y/6$, $y^2\leq6q/5$.

The finite multiplicities and the product-rule arithmetic are
\[
\begin{array}{c|c|c|c}
 (r,b)&j&k_-=r-j&k_+=r+j\\ \hline
 (1,y)&0&1&1\\
 (1,-y)&0&1&1\\
 (2,-y)&-1&3&1\\
 (2,-y)&1&1&3.
\end{array}
\]
Thus the eight reciprocal products have six $k=1$ terms and two
$k=3$ terms. On the near domain a product derivative is at most
\[
 \frac{10}{63}\left(\frac5{12}+\frac8{21}\right)
       \frac\rho{k^3}=\frac{335}{2646}\frac\rho{k^3}.
\]
On the wide domain it is at most
$(80/323)(10/19+8/17)\rho/k^3=(25760/104329)\rho/k^3$.
Multiplying by $6+2/27$ bounds the displacement-sum increments by
$4\rho/5$ and $3\rho/2$. The factor $4a/m$ in the exact
subtraction produces $16a\rho/(5n)$ and $6a\rho/n$.

For the geometric term the denominator sum is bounded by its four
$k_-k_+$ products:
\[
 \sum D_h^{-1}\leq
 \begin{cases}(25/144)(1+1+1/3+1/3)=25/54,&\rho\leq1/2,\\
 (100/361)(1+1+1/3+1/3)=800/1083,&\rho\leq1.
 \end{cases}
\]
From $|\operatorname{sinc}'u|\leq|u|/3$ and
$|\dot u|+|\dot v|\leq\rho/n$ one gets
$|\dot K|\leq7\rho/(6n^2)$. Differentiating
$H_h=(1+\cos v)/(2K)$ then gives
\[
 |\dot H_h|\leq
 \left\{\frac{50}{49}\frac72+
                   \left(\frac{50}{49}\right)^2\frac76\right\}
       \frac\rho{n^2}\leq\frac{5\rho}{n^2}.
\]
For $Q_H=\sum(H_h-1)D_h^{-1}$, the product rule therefore bounds
its derivative by $(5+19c_\sigma)(\rho/n^2)\sum D_h^{-1}$,
where $c_N=5/6$, $c_W=20/19$. After multiplication by $8\mu a$
the two constants are
\[
 8(5+19\cdot5/6)\frac{25}{54}=\frac{6250}{81}<78,
 \qquad
 8(5+19\cdot20/19)\frac{800}{1083}=\frac{160000}{1083}<148.
\]
This is the geometric column, with its origin already subtracted.

Since $|g_j|\leq7P/n$, its two costs are bounded by
\[
 \frac{14P}{n}\frac{50}{49}\frac{25}{54}<\frac{7P}n,
 \qquad
 \frac{14P}{n}\frac{50}{49}\frac{800}{1083}<\frac{11P}n.
\]
Finally the three nonzero $w$ values sum to $4\pi-y$; the full
$\tau$ error, including the closest row, is at most
$\{q(4\pi-y)+y^2\}/(4m^2)<4P/n^2$ using $y^2\leq6q/5$.
These calculations reproduce all four columns on both radii.
The four $-a/m$ contributions cancel exactly under $\delta_a$.
The remaining endpoint charges are then grouped according to what
changed: the denominator displacement, the sine-chord geometry, the
sine numerator, or the $\tau$ factor. Each group uses the same finite
positive-sample set. This organization prevents a small movement of a
removable sample from being charged as a whole new row. The central
rows require an additional argument because their limiting mass can
concentrate in a single residue.

For the two actual central rows use the linked cyclic vector
\[
 L_{n,k}(\theta)=\sin\theta\cot((k\pi+\theta)/n),\qquad
 L_{n,0}(0)=n,\quad L_{n,k}(0)=0\ (k\not\equiv0\pmod n).
\]
It is polynomial because $\sin(nt)\cot t=\cos t\,U_{n-1}(\cos t)$.

Here $U_{n-1}$ is the polynomial specified by
$\sin(nt)=\sin t\,U_{n-1}(\cos t)$. Thus the product
$\sin(nt)\cot t$ has the polynomial expression
$\cos t\,U_{n-1}(\cos t)$ through every apparent cotangent pole.
For the vector $L_n$, the variable is $t=(k\pi+\theta)/n$ and
$\sin\theta=(-1)^k\sin(nt)$. This gives a smooth continuation in
$\theta$ for each residue. At $\theta=0$ its zero residue has
value $n$ by $\sin\theta\cot(\theta/n)\to n$; all other residues
in one cycle have value zero. Hence the whole concentrated sample is
included in the vector and its derivative bounds.

The norm $\|L_n'\|_1$ means the sum of the absolute derivatives over
one full cycle of residues. Summing $1/|k|$ and $1/k^2$ over the
nonzero residues gives its harmonic bound. Integrating that norm along
the phase path bounds the change of the full row, including collisions.
The origin shifts in the second row come from its actual position
$m-1$, and remain $\pm\pi/(2m)$ in this calculation.
The exact vectors are
\[
 \begin{gathered}
 \frac{c_eL_{n,k}((x-\mu y)/2)+c_dL_{n,k}((x+\mu y)/2)}
                {2\cos(y/(2m))},\\
 \frac{c_eL_{n,k}((x-\mu y-\pi/m)/2)
             +c_dL_{n,k-1}((x+\mu y+\pi/m)/2)}
                {2\cos((\pi-y)/(2m))}.
 \end{gathered}
\]
All residues are included. For $|\theta|\leq11/20$, the centered
nonzero residues satisfy
$|L'_{n,k}|\leq2n/(5|k|)+11n/(50k^2)$; the zero residue costs
at most $n/5$. Hence $\|L_n'\|_1\leq(9/8)nH_n$, using
$H_n\geq7/2$. The row $m-1$ origin shifts are $\pm\pi/(2m)$,
with norms below $25n/16$. Since
$(c_e+c_d)/(2c)\leq8a/(7c)$ and $c>99/100$, each row's vector
motion costs at most $(2/3)a\rho nH_n$.
The remaining coefficient errors are $a\rho^2/6$ for row $m$ and
$a\rho^2/13+ay/m^2$ for row $m-1$. With the origin norms these
give at most $(3/4)a\rho H_n/n$ per normalized row.
Thus $|\delta_aJ_{\rm cen}|\leq(3/2)a\rho H_n/n$, proving the
retained entries, with both meshes and their finite collisions intact.

Here are details for the cyclic-vector constants. For a centered
nonzero residue the cotangent argument has distance at least
$(\pi|k|-11/20)/n\geq(5/2)|k|/n$ from $\pi\mathbb Z$.
The bounds $|\cot u|\leq1/\operatorname{dist}(u,\pi\mathbb Z)$
and $\csc^2u\leq\pi^2/(4\operatorname{dist}(u,\pi\mathbb Z)^2)$
then give
\[
 |L'_{n,k}|\leq\frac{2n}{5|k|}+\frac{11n}{50k^2}.
\]
At $k=0$ use $L_{n,0}=n\operatorname{sinc}\theta\,
\tau(\theta/n)$; $|\operatorname{sinc}'\theta|\leq|\theta|/3$
and $|\tau'(u)|\leq2|u|$ give
$|L'_{n,0}|\leq11n/60+11/(10n)\leq n/5$.
Summing both sides of the centered residue range gives
\[
 \|L_n'\|_1\leq\frac45nH_n+\frac{27}{25}n
 \leq\frac{194}{175}nH_n<\frac98nH_n,
\]
using $\sum_{k\geq1}k^{-2}<2$ and $H_n\geq7/2$.
The latter lower bound follows by grouping $1$, $2$, $3$--$4$,
$5$--$8$, $9$--$16$, $17$--$32$: after the first term each
group contributes at least $1/2$. Integrating the derivative bound
from zero to $\epsilon=\pi/(2m)$ gives
$\|L_n(\pm\epsilon)\|_1\leq n+(9/8)nH_n\epsilon<25n/16$,
using $H_n\leq n/4$.

Write $A=c_e/(2c)$, $B=c_d/(2c)$ for either row's weights.
Their sum is at most $8a/(7c)$, and each argument moves by at most
$\rho/2$. The vector-motion contribution before normalization is
therefore at most
\[
 \frac98nH_n\frac\rho2\frac{8a}{7c}
 =\frac9{14c}a\rho nH_n\leq\frac23a\rho nH_n.
\]
For row $m$, $A+B-a\leq a\rho^2/6$ follows from
$1/c_e-1,1/c_d-1\leq\rho^2/7$ and
$1/c-1\leq\rho^2/(8m^2c)$. Multiplying the residual coefficient
by $\|L_n(0)\|_1=n$ costs at most $a\rho H_n/(21n)$ after
normalization. Together with vector motion this is at most
$(5/7)a\rho H_n/n<(3/4)a\rho H_n/n$.
For row $m-1$, each coefficient differs from $a$ times its own origin
coefficient by at most $a\rho^2/13+ay/m^2$. Multiplying their sum
by the origin norms $25n/16$ and dividing by $nm$ costs
\[
 a\left\{\frac{25\rho^2}{104m}+\frac{25y}{8m^3}\right\}
 \leq\frac{a\rho}{4n}\leq\frac{a\rho H_n}{14n}.
\]
Here $25/104+25/(8m^2)\leq1/4$ for $m\geq33$.
Adding vector motion gives $31/42<3/4$ times $a\rho H_n/n$.
Thus both actual central rows fit the common $3/4$ bound separately.

\paragraph{The relative mesh correction.}
The leading mesh coefficients at a phase are proportional to $a$.
Weighted subtraction removes their large common parts. What remain
are differences of smooth offset series and the smaller coefficient
$D_0-aD_0(0)$. The notation $\delta\mathcal H_j$ in the next display
means an ordinary offset difference, whereas $\delta_aR_{\rm mesh}$
means the weighted phase difference. The two operations have different
roles. Differentiating the absolutely convergent series gives small
offset derivatives at the upper tail $j=n$; at $j=2$ it gives a
bounded derivative. Multiplication by $\epsilon_n=O(1/n)$ then
makes the complete relative mesh error linear in $\rho/n$.

For the mesh row, cancel its leading terms first:
\[
 \begin{aligned}
 \delta_aR_{\rm mesh}=-B_0\epsilon_n\delta\mathcal H_2
 &-aZ_n(0)\delta\mathcal H_n-B_0\delta\mathcal Q_n\\
 &-(D_0-aD_0(0))(\mathcal H_n+2L),
 \end{aligned}
\]
where $\delta\mathcal H_j=\mathcal H_j(t)-\mathcal H_j(0)$.
Differentiating \eqref{uc:mesh-series} gives bounds
$|\mathcal H_2'|\leq2$, $|\mathcal H_n'|\leq(n-1)^{-2}$,
$|\mathcal Q_n'|\leq(L-1)/(n-1)^2$, $|\mathcal H_n+2L|<n^{-2}$.
With $B_0\leq1/2$, $|Z_n(0)|<5$,
$|D_0-aD_0(0)|\leq5\rho/3$, $t\leq\rho/3$, this is at most
$\rho\{1/(3n)+(L+9)/(6(n-1)^2)+5/(3n^2)\}<\rho/(2n)$.

The differentiated tail series behind these bounds are
\[
 \begin{aligned}
 h_j'(v)&=-\sum_{k\geq j}(k+v)^{-2},&
 h_j''(v)&=2\sum_{k\geq j}(k+v)^{-3},\\
 g_j'(v)&=\sum_{k\geq j}\frac{1-\log(k+v)}{(k+v)^2},&
 g_j''(v)&=\sum_{k\geq j}\frac{2\log(k+v)-3}{(k+v)^3}.
 \end{aligned}
\]
For a paired function, differentiation subtracts the two first
derivatives at $v$ and $1-v$. The mean value theorem and integral
comparison therefore give
$|\mathcal H_n'|\leq2\int_{n-1}^\infty r^{-3}dr=(n-1)^{-2}$ and
\[
 |\mathcal Q_n'|\leq\int_{n-1}^\infty
                     \frac{2\log r-3}{r^3}dr
 =\frac{\log(n-1)-1}{(n-1)^2}\leq\frac{L-1}{(n-1)^2}.
\]
At $j=2$ the two reciprocal-square sums give
$|\mathcal H_2'|\leq2$. The actual paired upper endpoints
$n\pm c$, $|c|\leq1/2$, give
\[
 |\mathcal H_n+2L|\leq\frac1{4n^2-1}
                  +\frac1{4(n-1/2)^2}<\frac1{n^2}.
\]
Finally multiplying the displayed coefficient of $\rho$ by $n$
and using $(n-1)^2>7n^2/8$, $(L+9)/n\leq13/32$, bounds it by
$1/3+13/168+5/96<1/2$. This explains the finite constant in the
relative mesh row.

\paragraph{The signed outer increment.}
The outer term is the most delicate relative estimate because a
uniform absolute bound would charge the full size of several common
profiles. We separate its $L$ coefficient, regular scalar pole part,
complete cancellation, and the shift of the origin grid. The $L$
coefficient will have a favorable sign. The central logarithm already
appearing in the telescope is retained in the cancellation expression
and then charged once with $R_{\log}$.

The numerators $N_L,N_R$ keep the two actual folded scalar profiles.
Subtracting $2af(z)$ defines $V_s$, the motion relative to the weighted
origin. Removing its endpoint value $V_0$ in $P_s$ cancels the
reciprocal singularity. This is the same endpoint regularization used
in the absolute estimate, now applied after the origin subtraction.

It remains to prove the last row with its scalar logarithm included
once. Put $e=\rho$, $\delta=x-y$, and on either folded side set
\[
 \begin{gathered}
 N_L=c_ef(z+\delta)+c_df(z-e),\quad
 N_R=c_ef(z-\delta)+c_df(z+e),\\
 V_s=N_s-2af(z),\quad V_0=V_s(0),
 \quad P_s=(V_s-V_0)/z+V_sd(z).
 \end{gathered}
\]
With $Z=2a\cot(z/2)$, define
$\pi g_L=\max(q,p-Z)$, $\pi g_R=\max(-q,p-Z,q-2Z)$ and
$\mathcal C_s=I_s(\phi,z)-aI_s(0,z)-d_+\log(\pi-z)$.
The exact outer increment is
\begin{align}\label{uc:signed-outer}
 T_{\rm out}={}&L\{Q_t[g_L]+Q_{1-t}[g_R]\}
 +Q_t[P_L]+Q_{1-t}[P_R]
 +\{Q_t[\mathcal C_L]+Q_{1-t}[\mathcal C_R]\}/\pi\notag\\
 &+\mathrm{Shift}+(d_+/\pi)R_{\log},\notag\\
 \mathrm{Shift}={}&\frac am\sum_{k=2}^{m-2}
 \{F_0(h(k+t))+F_0(h(k+1-t))-F_0(hk)-F_0(h(k+1))\}.
\end{align}
The functions $g_L-q/\pi=(p-q-Z)_+/\pi$ and
$g_R+q/\pi=\max(0,p+q-Z,2q-2Z)/\pi$ are nonnegative and
nondecreasing. Their disjoint right rectangles lie in $[0,\pi]$,
so each $Q_\tau$ is nonpositive. Since $Q_\tau[1]=-3/m$,
the opposite constants cancel. Thus the whole $L$ term is favorable.

For clarity, the sign argument uses intervals to the right of each
sample. If $g\geq0$ is nondecreasing, then
$hg(z_k)\leq\int_{z_k}^{z_k+h}g(z)dz$. These right intervals are
disjoint and lie in $[0,\pi]$, so their sample sum is bounded by the
whole integral. Thus $Q_\tau[g]\leq0$. The two functions obtained
by subtracting the opposite constants have exactly these properties.
Also $Q_\tau[1]=(m-3)/m-1=-3/m$, and the constants $q/\pi$
and $-q/\pi$ cancel. This proves the favorable sign without replacing
it by an absolute variation bound.

\paragraph{Deriving the pole-regular table from its weights.}
Use $e=\rho$, $\delta=x-y$, and $r=c_e+c_d-2a$. Each numerator
difference has the form
\[
 V_s=c_e\{f(z\pm\delta)-f(z)\}
       +c_d\{f(z\mp e)-f(z)\}+rf(z),
 \qquad 0\leq r\leq\rho^2/4.
\]
On the near domain $r\leq\rho/8$; on the wide domain
$r\leq\rho/4$. Lipschitz continuity bounds the two translated
terms by $2(|\delta|+e)\leq4\rho$.
At $z=0$, evenness and $f'(0)=0$ improve this to
\[
 |V_0|\leq(\delta^2+e^2)/2+r\leq5\rho^2/4.
\]
On the initial half all shifted arguments stay in $(-\pi,\pi)$,
where the derivative bounds are ordinary bounds. Differentiating the
translated differences gives the following four inputs:
\[
\begin{array}{c|cccc}
 &\sup|V_s|&|V_0|&\sup|V_s'|&\sup|V_s''|\\ \hline
 \rho\leq1/2&(33/8)\rho&(5/8)\rho&(9/4)\rho&(57/40)\rho\\
 \rho\leq1&(17/4)\rho&(5/4)\rho&(5/2)\rho&(31/20)\rho.
\end{array}
\]
Here $|V_s'|\leq|\delta|+e+2r$ and
$|V_s''|\leq(13/20)(|\delta|+e)+r$; these display how the
second and third scalar derivatives enter.

For $A_s(z)=(V_s(z)-V_0)/z$, integration by parts gives
\[
 A_s'(z)=z^{-2}\int_0^z uV_s''(u)du,
 \quad |A_s'|\leq\tfrac12\sup|V_s''|,
 \quad |A_s|\leq\sup|V_s'|.
\]
Over an interval of length $\pi/2$, the near and wide variations are
therefore at most $57\pi\rho/160<57\rho/50$ and
$31\pi\rho/80<31\rho/25$. The other term is $V_sd$ with
$d\leq1/5$ increasing. Its two variation bounds are
\[
 \begin{aligned}
 \operatorname{Var}(V_sd)
 &\leq\tfrac15\sup|V_s|+\tfrac15(\pi/2)\sup|V_s'|\\
 &\leq\begin{cases}309\rho/200,&\rho\leq1/2,\\
                         33\rho/20,&\rho\leq1.
          \end{cases}
 \end{aligned}
\]
Thus the initial variations are $57/50+309/200=537/200$ and
$31/25+33/20=289/100$ times $\rho$. The initial suprema are
$9/4+33/40=123/40$ and $5/2+17/20=67/20$ times $\rho$.

On the central interval $I=[\pi/2,\pi]$, measure translation
bounds the translated terms' combined variation by
$4(|\delta|+e)\leq8\rho$. The residual $rf$ has variation
at most $\pi r<4r$. Hence
$\operatorname{Var}_I V_s\leq17\rho/2$ near and $9\rho$ wide.
Now write $P_s=w_1V_s-w_2V_0$, where
\[
 w_1(z)=\frac1{2\sin(z/2)},\qquad w_2(z)=\frac1z.
\]
Both decrease on $I$. Their endpoint values give
\[
 \begin{array}{c|cc}
 &\sup_I w_i&\operatorname{Var}_I w_i\\ \hline
 w_1&1/\sqrt2<3/4&1/\sqrt2-1/2<1/4\\
 w_2&2/\pi<2/3&1/\pi<1/3.
 \end{array}
\]
The product variation inequality therefore gives
\[
 \begin{aligned}
 \operatorname{Var}_IP_s
 &\leq\tfrac34\operatorname{Var}_IV_s
              +\tfrac14\sup|V_s|+\tfrac13|V_0|,\\
 \sup_I|P_s|&\leq\tfrac34\sup|V_s|+\tfrac23|V_0|.
 \end{aligned}
\]
For the near domain these coefficients are
$51/8+33/32+5/24=731/96$ and
$99/32+5/12=337/96$. For the wide domain they are
$27/4+17/16+5/12=395/48$ and
$51/16+5/6=193/48$. This derives every entry in the table below,
including the effects of a translated cusp.

Here are the pole-regular bounds for each side, each entry multiplied
by $\rho$; the columns give variation and supremum on the initial
and central half intervals:
\[
\begin{array}{c|cc|cc}
 &\operatorname{Var}P_s&\sup|P_s|&\operatorname{Var}P_s&\sup|P_s|\\ \hline
\rho\leq1/2&537/200&123/40&731/96&337/96\\
\rho\leq1&289/100&67/20&395/48&193/48.
\end{array}
\]
To verify the wide row use $r=c_e+c_d-2a\leq\rho^2/4$ and
\eqref{uc:scalar-bounds}: initially
$|V_s|\leq17\rho/4$, $|V_0|\leq5\rho/4$,
$|V_s'|\leq5\rho/2$, $|V_s''|\leq31\rho/20$.
The identity $((V_s-V_0)/z)'=z^{-2}\int_0^z wV_s''(w)dw$
gives variation $31\rho/25$, while $V_sd$ costs $33\rho/20$.
The relative scalar bounds use both smooth motion away from the
cusp and finite variation across it. Initially the shifted arguments
stay in a range where ordinary derivatives provide the listed bounds.
The identity for $((V_s-V_0)/z)'$ follows by integrating $V_s''$
twice; it removes the apparent division by $z^2$. Centrally a phase
shift may cross the cusp. The finite derivative measure records its
derivative jump as a point mass. Thus the
variation of a translate difference is at most the translation distance
times $\operatorname{Var}f'$. This follows by writing a translate
difference as an integral of translated derivatives and then summing
absolute changes. The uniform bound $\operatorname{Var}f'<4$
therefore supplies the next estimate.

Centrally the measure translation bound
$\operatorname{Var}[f(\cdot+b)-f]\leq4|b|$ gives
$\operatorname{Var}V_s\leq9\rho$. In
$P_s=V_s/[2\sin(z/2)]-V_0/z$, the two weights have
supremum/variation $(3/4,1/4)$ and $(2/3,1/3)$, proving the row.
Near the origin retain $r\leq\rho/8$; the corresponding four
initial bounds are $33\rho/8,5\rho/8,9\rho/4,57\rho/40$ and the
central variation is $17\rho/2$. This proves the sharper row.
BV midpoint quadrature and holes of total $4h,2h$ now give
$|Q_t[P_L]+Q_{1-t}[P_R]|\leq30\rho/m$ near and $33\rho/m$ wide;
for example the wide sum is $39073\rho/(1200m)<33\rho/m$.

The conversion of that table into quadrature charges uses both
sides and all omitted intervals. If $V_i,V_c$ are its per-side
initial and central variation coefficients and $M_i,M_c$ its supremum
coefficients, the two midpoint errors sum to
$\rho(V_i+V_c)/m$. Initial holes have total length $4h$ and
central holes $2h$, giving an additional
$\rho(4M_i+2M_c)/m$. Thus the two sums are
\[
 \begin{aligned}
 \rho\leq1/2:\quad&
 \frac\rho m\left\{\frac{537}{200}+\frac{731}{96}
                   +4\frac{123}{40}+2\frac{337}{96}\right\}
 =\frac{71089}{2400}\frac\rho m<\frac{30\rho}m,\\
 \rho\leq1:\quad&
 \frac\rho m\left\{\frac{289}{100}+\frac{395}{48}
                   +4\frac{67}{20}+2\frac{193}{48}\right\}
 =\frac{39073}{1200}\frac\rho m<\frac{33\rho}m.
 \end{aligned}
\]
The midpoint factor $1/2$ and the two sides cancel each other in the
variation contribution; the hole lengths retain their full factors
four and two.

For the complete cancellation, the initial half is root-free because
$(a\cot(z/2)-q/2)\sin z\geq a-q/2\geq p/2+17/64$.
Its exact increments are $\pm q\log\sin z$.
Centrally put $e=\pi-z$, $s_0=a\tan(e/2)$,
$s_L=s_0+q/2$, $s_R=|s_0-q/2|$. Then
\[
 \begin{gathered}
 \mathcal C_s=R(p/2,s_s,e)+b_s(e)\log\sin e
                                  +d_+\log(\sin e/e),\\
 b_L=\max(q-p,-2s_0)-d_+,
 \quad b_R=\max(-p-q,q-p-4s_0,-2s_0)-d_+.
 \end{gathered}
\]
The functions $b_s$ separate the explicit logarithmic coefficient
from the bounded remainder $R$. Near $e=0$ their value is zero, so
the product with $\log\sin e$ remains integrable even when its
derivative is large. The term $d_+\log(\sin e/e)$ has its singular
logarithm cancelled already. The remaining issue is how much $R$
changes as its coefficient parameters pass through the equality
threshold. Integrating its logarithmic parameter derivative gives the
bound below.

The $b_s$ decrease from zero, with combined supremum and variation
at most $3\rho$. The $s_s$ have at most three monotone pieces and
are at most $3/2$. For $\epsilon=h/2$, $\ell=-\log\sin\epsilon$,
if $\rho\leq\epsilon^2/8$ then on every cell
$(a\tan(e/2)-q/2)\sin e\geq7e^2/80\geq\rho/2\geq p/2$,
so there are no central roots. Otherwise the parameter integral in
\eqref{uc:R-bounds} is at most
$p[1+\log(4/p)]\leq\rho(11/2+2\ell)$, using
$p\log(\rho/p)\leq\rho/\mathrm e$.
The split at $\rho=\epsilon^2/8$ reflects the quadratic central
root threshold. For very small phase the first available cell is too
far from the center to contain a root, so a direct root-free formula
applies there. For larger phase the logarithmic parameter integral is
used. The factor $p\log(\rho/p)$ stays bounded even when $p$ is
much smaller than $\rho$: putting $r=p/\rho$ reduces it to
$\rho r\log(1/r)\leq\rho/\mathrm e$. Thus the estimate is linear
in $\rho$ times a mesh logarithm and remains continuous on the axis
$p=0$. This two-case argument covers the moving root region without
assuming differentiability of the root locations at their merger.

In either case the full two-side variation is at most
$\rho(39/2+16\ell)$. Initial and central hole envelopes are
$\rho[\log(\pi/(2z))+\log\pi]$ and
$\rho[7+4(-\log\sin e)]$, respectively. Integrating them and
using midpoint BV error yields, respectively,
\[
 \frac\rho{\pi m}(39/4+8\log m),\quad
 \frac{4\rho}{\pi m}(3+\log m),\quad
 \frac\rho{\pi m}(22+8\log m).
\]
Thus the full cancellation costs
$\rho(175/4+20\log m)/(\pi m)$, including every root switch.

Here is the complete variation ledger for this last cancellation
charge. On the initial half the exact increments are
$\pm q\log\sin z-d_+\log(\pi-z)$, with combined variation at
most $2\rho\ell+2\rho$. On the central half the two parameter
curves have at most three monotone pieces, so their $R_s$ contribution
is at most $3G$, where $G=p[1+\log(4/p)]$. The separation part is
at most $2p\ell$.
If $\rho\leq\epsilon^2/8$, then
\[
 (a\tan(e/2)-q/2)\sin e
 =a(1-\cos e)-(q/2)\sin e
 \geq\frac{3e^2}{20}-\frac\rho2
 \geq\frac{7e^2}{80}\geq\frac\rho2\geq\frac p2.
\]
Here $1-\cos e\geq2e^2/\pi^2>e^2/5$ and $a\geq3/4$.
Both parameters are root-free on the truncated interval, and their
remainders are exactly $p\log\sin e$. Otherwise
\[
 G=p\{1+\log(4/\rho)\}+p\log(\rho/p)
 \leq\rho\{3/2+\log(4/\rho)\}
 \leq\rho(11/2+2\ell),
\]
using $\rho>\epsilon^2/8$, $\log32<4$, and
$\log(1/\epsilon)\leq\ell$.
In either case the central remainder variation is at most
$\rho(33/2+8\ell)$. The two $b_s$ have combined supremum and
combined variation at most $3\rho$, so their products with
$\log\sin e$ cost $6\rho\ell$. The two
$d_+\log(\sin e/e)$ terms cost at most
$2d_+\log(\pi/2)<\rho$. Adding the initial half gives
\[
 \rho\{(33/2+8\ell)+6\ell+1+2\ell+2\}
 =\rho(39/2+16\ell).
\]
The midpoint charge is consequently
$\rho(39/4+8\log m)/(\pi m)$.

For the holes, integrating $\log(A/z)$ gives
$z\{1+\log(A/z)\}$. This function is concave, so two lengths
with a fixed sum are bounded by twice their common-average value.
The initial lengths sum to $4h$, giving at most
\[
 \frac{4\rho}{\pi m}\{1+\log(\pi m/4)\}
 \leq\frac{4\rho}{\pi m}(3+\log m).
\]
The central lengths sum to $2h$, giving at most
$2\rho\{11+4\log(m/2)\}/(\pi m)$, bounded by
$\rho(22+8\log m)/(\pi m)$. Thus constants and logarithms add to
$39/4+12+22=175/4$ and $8+4+8=20$, respectively.
The logarithm contributes $d_+W_n$, already proved above.
The shift term compares the same origin profile on two slightly
translated row meshes. Bounding each translated sample separately
would add an unnecessary error over all rows. Instead, neighboring
first differences cancel in the finite sum. Only boundary differences
and disjoint interior second differences remain. The former are
controlled by $\sup|F_0'|$ and the latter by
$\operatorname{Var}F_0'$. The actual offset is proportional to
$y$, which explains the separate $y(17L+27)/m^2$ term and the
small difference between the $x$ and $y$ coefficients at the end.

Finally first differences in the shift sum telescope to two boundary
terms and disjoint interior second differences. On $[2h,\pi-h]$,
the differentiation used in \eqref{uc:origin-variation} gives
\[
 \sup|F_0'|\leq(4/3)L+(2/3)\log m+20/3,\quad
 \operatorname{Var}F_0'\leq L/2+(25/2)\log m+37/3.
\]
Their sum $2\sup|F_0'|+\operatorname{Var}F_0'$ is less than
$17L+27$, using $\log m\leq L+1/32$. Hence
$|\mathrm{Shift}|\leq y(17L+27)/m^2$.

The finite cancellation in the shift can be written explicitly.
Put $\eta=ht=y/m$ and $z_j=hj$. Regrouping the actual finite sum gives
\[
 \begin{aligned}
 \frac m a\mathrm{Shift}={}&F_0(2h+\eta)-F_0(2h)
       +F_0(\pi-h-\eta)-F_0(\pi-h)\\
 &+\sum_{j=3}^{m-2}
       \{F_0(z_j+\eta)+F_0(z_j-\eta)-2F_0(z_j)\}.
 \end{aligned}
\]
Equivalently multiply this identity by $a$, so no division is needed
when a zero weight is considered. The intervals
$[z_j-\eta,z_j+\eta]$ are disjoint because $2y\leq2<\pi$.
For any such interval,
\[
 F_0(z+\eta)+F_0(z-\eta)-2F_0(z)
 =\int_0^\eta\{F_0'(z+r)-F_0'(z-r)\}dr,
\]
whose absolute value is at most $\eta$ times the derivative variation
on that interval. The two boundary differences cost $2\eta M$,
where $M=\sup|F_0'|$, and the disjoint interior terms cost $\eta V$.
Hence $|\mathrm{Shift}|\leq ay(2M+V)/m^2$.
The derivative estimates used here give
\[
 2M+V\leq\frac{19}{6}L+\frac{83}{6}\log m+\frac{77}{3}
 \leq17L+27,
\]
since $\log m\leq L+1/32$ and $77/3+83/192<27$.
To check $M$ directly, the cotangent residual derivative contributes
at most $4L/3$, the quotient derivative in $P_0$ at most one, and
its product derivative at most $4/5+8/3$, altogether less than five.
On the initial half $|Q_0'|\leq\log m+5$; centrally differentiating
$2\tan(e/2)\log\sin e$ gives $|Q_0'|\leq2\log m+5$.
Division by $\pi>3$ gives
$M\leq(4/3)L+(2/3)\log m+20/3$, as used above.
Adding the last four estimates gives
\[
 \begin{gathered}
 T_{\rm out}\leq\frac\rho m
 \left\{c_P+\frac{23\log m+175/4+\log2+2\log\pi+2}\pi\right\}
                  +\frac{y(17L+27)}{m^2},\\ c_P=30\text{ or }33.
 \end{gathered}
\]
The elementary bounds $\pi>3$, $\log2<7/10$,
$\log\pi<23/20$, $\log m\leq L+1/32$ give the constants $47,50$
in \eqref{uc:relative-components}. All exact increments vanish at
zero. The arguments use integrable complete cancellations and linked
finite rows; a dense collision-free sequence therefore proves the
same estimates on all stated closed domains.

\paragraph{Assembling the comparison with the finite corner.}
The estimates have now been proved for all pieces of the original
identity. Passing to an axis or collision uses complete expressions:
the polynomial finite rows are continuous, the matched series have
summable tails, and the integrals have the common endpoint and root
envelopes already proved. Taking a collision-free sequence therefore
preserves the estimates at its limit. This also explains why all
finite retained values were specified before the estimates began.

The weighted subtraction is converted to an ordinary error difference
by the exact algebra
\[
 E_n(x,y)-E_n(0)=\delta_aE_n-uE_n(0)\leq\delta_aE_n+u\bar e.
\]
The inequality uses $u\geq0$ and the one-sided origin bound. In the
table the factor $aA_4$ becomes $A_4-uA_4$, and
$a\rho J_\sigma$ becomes $\rho J_\sigma-u\rho J_\sigma$.
These negative $u$ terms strengthen the final model loss. The
unweighted allowances combine into $C_4$ and the linear phase terms.
Replacing $P$ by $\rho=x+y$ keeps their explicit coordinate form;
the shift charge adds $d_\beta$ only to the $y$ coefficient.

Summing \eqref{uc:relative-components}, using $P\leq\rho$ and
$a=1-u$, gives
\[
 E_n(\phi)-E_n(0)\leq C_4+k_{\sigma x}x+k_{\sigma y}y
              +u(\bar e-A_4-\rho J_\sigma),\qquad \sigma=N,W,
\]
where we used $E_n(\phi)-E_n(0)=\delta_aE_n-uE_n(0)$ and
\eqref{uc:origin-lower}. The constants match term by term:
near $16/5+10+7+47+1/2=677/10$, wide
$6+10+11+50+1/2=155/2$; the square-denominator terms are
$78+4=82$, $148+4=152$, and both harmonic coefficients are
$8+3/2=19/2$. Finally \eqref{uc:model-loss} and
\eqref{uc:regular-loss} give the model drop at least $u(T+j)-E_c$.
Since $S_n(0)=\Lambda_c(n)$, this proves
\eqref{uc:relative-comparison} and completes the finite comparison proof.

The model part of the last step can also be read directly:
\[
 M_\infty(x,y,L)-M_\infty(0,0,L)
 \leq-u(T+j)+E_c.
\]
The leading coefficient contributes $-4uL^2/\pi^2$, the possible
linear increase contributes at most $3uL/\pi^2$, and the regular
constant contributes at most $-ju$. Adding the finite error difference
gives $-uU_\sigma$ with precisely the $U_N,U_W$ defined at the
start. Finally $S_n(0)=\Lambda_c(n)$ is an exact finite identity.
The result therefore compares the full finite board with its actual
corner, while retaining the central root contribution and the distinct
near and wide allowances for the next localization argument.
